\chapter{Le Razze e le Etnie}
%% \newcommand\TitoloRazza[2]{\subsubsection{#1}\noindent\includegraphics{bandierine.eps}\par\nobreak\margintitle{#2}}
\newcommand\TitoloRazza[2]{\medskip\par\nobreak\subsubsection{#1}\nobreak\ \margintitle{#2}}
\newcommand{\Politica} {\TitoloRazza{Politica}{\iconcorona}}
\newcommand{\Economia} {\TitoloRazza{Economia}{\iconsacco}}
\newcommand{\Religione} {\TitoloRazza{Religione}{\icontiara}}
\newcommand{\Magia} {\TitoloRazza{Magia}{\iconsfera}}
\newcommand{\Moda} {\TitoloRazza{Moda, cultura, costumi}{\icongiacca}}
\newcommand{\Geografia} {\TitoloRazza{Geografia}{\iconmappa}}
\newcommand{\Esercito} {\TitoloRazza{Esercito e Armi}{\iconspada}}
\newcommand{\Fisico} {\TitoloRazza{Caratteristiche fisiche}{\iconbobbie}}
\newcommand{\Stendardo}[1] {\begin{center}\includegraphics{#1}\end{center}}

\newcommand {\minmaxumani} {
\subsubsection{Massimi e minimi alle Caratteristiche}

\vspace{0.3cm}
{\centering \begin{tabular}{|l|c|c|} \hline & Max& Min\\
    \hline \hline INT& 20& 3\\
    \hline CONC& 20& 3\\
    \hline OSS& 20& 3\\
    \hline AGI& 20& 3\\
    \hline COS& 20& 5\\
    \hline FOR& 20& 3\\
    \hline BEL& 20& 3\\
    \hline CAR& 20& 3\\
    \hline
\end{tabular}\par}
\vspace{0.3cm}
}

\newcommand {\minmaxelfi} {
\subsubsection{Massimi e minimi alle Caratteristiche}

\vspace{0.3cm}
{\centering \begin{tabular}{|l|c|c|} \hline & Max& Min\\
    \hline \hline INT& 20& 3\\
    \hline CONC& 20& 3\\
    \hline OSS& 22& 3\\
    \hline AGI& 20& 3\\
    \hline COS& 16& 5\\
    \hline FOR& 20& 3\\
    \hline BEL& 22& 12\\
    \hline CAR& 20& 3\\
    \hline
\end{tabular}\par}
\vspace{0.3cm}
}

\newcommand {\minmaxgiganti} {
\subsubsection{Massimi e minimi alle Caratteristiche}

\vspace{0.3cm}
{\centering \begin{tabular}{|l|c|c|} \hline & Max& Min\\
    \hline \hline INT& 12& 3\\
    \hline CONC& 20& 3\\
    \hline OSS& 22& 3\\
    \hline AGI& 16& 3\\
    \hline COS& 25& 15\\
    \hline FOR& 25& 15\\
    \hline BEL& 20& 3\\
    \hline CAR& 20& 3\\
    \hline
\end{tabular}\par}
\vspace{0.3cm}
}

\newcommand {\minmaxnani} {
\subsubsection{Massimi e minimi alle Caratteristiche}

\vspace{0.3cm}
{\centering \begin{tabular}{|l|c|c|} \hline & Max& Min\\
    \hline \hline INT& 20& 3\\
    \hline CONC& 20& 3\\
    \hline OSS& 20& 3\\
    \hline AGI& 20& 3\\
    \hline COS& 22& 9\\
    \hline FOR& 20& 5\\
    \hline BEL& 18& 3\\
    \hline CAR& 20& 3\\
    \hline
\end{tabular}\par}
\vspace{0.3cm}
}

\newcommand {\minmaxgnomi} {
\subsubsection{Massimi e minimi alle Caratteristiche}

\vspace{0.3cm}
{\centering \begin{tabular}{|l|c|c|} \hline & Max& Min\\
    \hline \hline INT& 20& 3\\
    \hline CONC& 22& 6\\
    \hline OSS& 22& 9\\
    \hline AGI& 20& 3\\
    \hline COS& 16& 3\\
    \hline FOR& 18& 3\\
    \hline BEL& 20& 3\\
    \hline CAR& 22& 3\\
    \hline
\end{tabular}\par}
\vspace{0.3cm}
}

\newcommand {\minmaxorchi} {
\subsubsection{Massimi e minimi alle Caratteristiche}

\vspace{0.3cm}
{\centering \begin{tabular}{|l|c|c|} \hline & Max& Min\\
    \hline \hline INT& 18& 3\\
    \hline CONC& 20& 3\\
    \hline OSS& 21& 5\\
    \hline AGI& 21& 5\\
    \hline COS& 22& 9\\
    \hline FOR& 22& 9\\
    \hline BEL& 12& 2\\
    \hline CAR& 20& 3\\
    \hline
\end{tabular}\par}
\vspace{0.3cm}
}

\newcommand {\minmaxreuben} {
\subsubsection{Massimi e minimi alle Caratteristiche}

\vspace{0.3cm}
{\centering \begin{tabular}{|l|c|c|} \hline & Max& Min\\
    \hline \hline INT& 20& 3\\
    \hline CONC& 20& 3\\
    \hline OSS& 20& 3\\
    \hline AGI& 20& 3\\
    \hline COS& 22& 6\\
    \hline FOR& 22& 6\\
    \hline BEL& 9& 1\\
    \hline CAR& 20& 3\\
    \hline
\end{tabular}\par}
\vspace{0.3cm}
}

\section*{Introduzione}

In questo capitolo verranno descritte le razze e le etnie che popolano l'Arcipelago,
nonch\'e le terre da loro abitate.

Per \textbf{ogni razza} ed etnia verranno indicati:
\begin{itemize}\itemsep -3pt
\item La \textbf{geografia} delle terre in cui sono stanziati e le citt\`a
  pi\`u importanti;
\item il \textbf{sistema di governo} su cui poggiano le fondamenta della vita
  politica;
\item le risorse su cui si basa il \textbf{sistema economico};
\item l'\textbf{organizzazione militare} adibita alla difesa del
  territorio;
\item le \textbf{pratiche religiose}, i culti e le credenze in uso presso gli
  abitanti;
\item gli \textbf{insegnamenti magici} a cui \`e possibile accedere;
\item le \textbf{mode}, la cultura e i costumi pi\`u diffusi;
\item le \textbf{caratteristiche fisiche} degli abitanti.

\end{itemize}

Vengono sucessivamente indicati:

\iffullversion
\begin{table*}[hbt]
  \begin{center}\Large
    \begin{tabular}{|l|c|c|c|c|c|c|c|} \hline
            & Umani & Elfi & Orchi & Gnomi & Nani & Reuben & Giganti \\ \hline\hline
      INT   & 3-20  & 3-20 & 3-18  & 3-20  & 3-20 & 3-20   & 3-12 \\ \hline
      CONC  & 3-20  & 3-20 & 3-20  & 6-22  & 3-20 & 3-20   & 3-20 \\ \hline
      OSS   & 3-20  & 3-22 & 5-21  & 6-22  & 3-20 & 3-20   & 6-22 \\ \hline
      AGI   & 3-20  & 3-20 & 5-21  & 3-20  & 3-20 &12-20   & 3-16 \\ \hline
      COS   & 5-20  & 3-16 & 9-22  & 3-16  & 9-22 & 6-22   &15-25 \\ \hline
      FOR   & 3-20  & 3-20 & 9-22  & 3-18  & 3-20 & 6-22   &15-25 \\ \hline
      BEL   & 3-20  &12-22 & 2-12  & 3-20  & 3-18 & 1-9    & 3-20 \\ \hline
      CAR   & 3-20  & 3-20 & 3-20  & 6-22  & 3-20 & 3-20   & 3-20 \\ \hline
    \end{tabular}
  \end{center}
  \caption{Tabella comparativa minimi e massimi per le Caratteristiche}
\end{table*}
\fi  

\textbf{I massimi ed i minimi delle caratteristiche}, cio\`e le
soglie per ogni caratteristica, che non possono essere superate, nella
divisione dei Punti CAT (\ref{creazione} al momento della creazione,
vedi pagina \pageref{creazione})
  
Le \textbf{Abilit\`a di Razza} con i rispettivi VAL di base, cio\`e
quelle capacit\`a connaturate alla razza e all'etnia considerata.

Vengono poi evidenziate le \textbf{Classi sociali} a cui \`e
possibile appartenere in una determinata Razza ed Etnia. Per ognuna
delle Classi sociali sono indicati:

\begin{itemize}

\item Le \textbf{probabilit\`a di farvi parte} se si \`e un
  personaggio giocante (\textbf{PG}), espresse come intervallo di
  valori ottenuti dal lancio di 1d100; 
  
  le probabilit\`a indicate non rispecchiano necessariamente le
  proporzioni della popolazione;
\item il \textbf{Fattore Istruzione (FI)} proprio di ogni classe
  sociale (a destra);
\item una breve \textbf{descrizione} delle Classi sociali che mostra
  le professioni tipiche e gli introiti mensili;
\item l'elenco delle \textbf{Abilit\`a di Classe sociale}.
\end{itemize}

\subsection*{Mini-glossario}

In questo capitolo troverete ``sigle strane'' che verranno spiegate in
seguito. Comunque, per farvi un'idea, tenete conto che:

\begin{itemize}\itemsep -3pt  
\item per ATL si intende l'abilit\`a di combattimento Armi da Taglio
  Lunghe.
\item per ABL si intende l'abilit\`a di combattimento Armi da Botta
  Lunghe.
\item per ATC si intende l'abilit\`a di combattimento Armi da Taglio
  Corte.
\item per ABC si intende l'abilit\`a di combattimento Armi da Botta
  Corte.
\item per AIA si intende l'abilit\`a di combattimento Armi In Asta.
\item per ADT si intende l'abilit\`a di combattimento Armi da Tiro.
\item per ADL si intende l'abilit\`a di combattimento Armi da Lancio.
\item per AO si intende l'abilit\`a di combattimento Armi Occasionali.
\item per CAC si intende l'abilit\`a di combattimento Corpo A Corpo.
\item per \emph{Abilit\`a Professionale} si intende un'abilit\`a
  correlata al mestiere svolto o alla professione esercitata. Devono
  essere concordate con il Master in funzione della professione scelta
  dal giocatore per il suo PG.
    
\item per \emph{Abilit\`a a scelta del Master}, questi deve tener
  conto delle attitudini, delle inclinazioni e del Background del
  Personaggio.
    
\item per \emph{Lingua Straniera} si intende una lingua diversa dalla
  lingua Madre e dalla Lingua Levante (vedi pag. 82 Abilit\`a Lingua) 

\end{itemize}

\iffullversion
\pinupbig{ragazzo_falco2.eps}{}{tbh}
\fi

Qualora siano previste pi\`u abilit\`a alternative (per esempio:
Ventriloquio o Giocoleria o Acrobazia), \`e necessario optare per UNA
soltanto di esse.

\es{Abilit\`a professionali tipiche del Medico possono essere:
  Medicina, Pronto soccorso, Alchimia, Conoscere Piante, ecc.}

\es{ Abilit\`a professionali di un Falegname: Carpenteria, Conoscere
  Piante, Architettura, Ingegneria, ecc.}


\iffullversion
\section*{Mezze Razze}

\`E possibile giocare personaggi ``meticci'' derivati dall'unione di
genitori di razza diversa.

I massimi (e i minimi) alle caratteristiche in quel caso sono pari
alla media dei massimi (minimi), arrotondata per difetto, delle
caratteristiche delle razze dei genitori.

Per quanto riguarda l'estrazione sociale, si potr\`a scegliere quella
di uno qualunque dei due genitori.

I meticci non possono procreare.


\es{Brimbrim \`e un Mezzorco, nato da madre Umana Nomade e padre
  Orco Haroka.  I suoi massimi saranno la media tra i massimi degli
  umani e quelli degli orchi, e cos\`{\i} minimi. 
  
  Brimbrim ha vissuto fin da piccolo con i Nomadi e la classe sociale
  verr\`a scelta fra quelle disponibili fra i Nomadi.}
\fi

\clearpage
\section{Gli Umani}

\subsection{La Terra di Levante}

\Stendardo{stendardo_levante.eps}

I Levanti sono un'etnia umana stanziata nella Terra di Levante. 

\Geografia La Terra di Levante occupa la parte meridionale dell'Isola
di Dyshu-Thandy; confina a nord con la Teocrazia di Ryless ed \`e
separata a Sud da Bahuney dallo Stretto di Samar. Le origini del nome
``Terra di Levante'' si perdono nella notte dei tempi.

La superficie del territorio risulta ugualmente divisa tra pianura, collina,
montagna, con prevalenza di monti al Nord-Ovest e di pianura a Sud-Est. 

Le formazioni montuose pi\`u importanti sono la Catena Muraria a nord
(cima pi\`u alta: Monte Rama, 4412 m), che costituisce il confine
naturale con la Teocrazia di Ryless, e i monti Dolani (cima pi\`u
alta: Altavir, 6001 metri) che percorrono tutta la costa occidentale
da nord a Sud. Rilevante la presenza, tra Dolani e Catena Muraria,
dell'altopiano desertico di Jeomath, di origine vulcanica, circondato
da alte ed impervie montagne.

I Dolani danno vita al fiume Majirsus, il fiume pi\`u lungo ed
importante di tutto il continente, lungo oltre 2000 km ed in gran
parte navigabile, che sfocia nel mare Interno. La vastissima pianura
di Maja, originata dai fiumi della rete idrografica del Majirsus, \`e
fertile e produttiva.

Le coste sono basse e ricche di porti naturali a Sud e ad Est (se si
esclude Capo Samar), alte e scoscese in quasi tutta la costa
occidentale, fatte salve le piccole pianure che costituiscono il
Prospero Ovest. Importante a Est il Golfo di Thandy, nel quale sfocia
ad estuario il Majirsus.

La vegetazione naturale \`e costituita da boschi di latifoglie che
crescono al centro, foreste di conifere sulle montagne ad Ovest e
nord, e macchia mediterranea e pinete al Sud. Le citt\`a pi\`u
importanti sono: Nuova Giassa, Lorn, Tentsa, Rubia e la capitale
Dyshu.

\Politica Terra di Levante \`e una repubblica democratica
parlamentare, la cui popolazione \`e costituita in maggioranza da
Levanti. L'organo pi\`u importante \`e il Gran Consiglio, composto da
199 membri eletti ogni 6 anni a suffragio universale tra tutti i
cittadini Levanti che abbiano compiuto 23 anni, che ha il compito di
legiferare su questioni che riguardano l'intera nazione.

\iffullversion
Una volta insediatosi, il Gran Consiglio elegge tra i suoi membri: 

\begin{itemize}\itemsep -3pt
\item il \textbf{Presidente del Gran Consiglio}, che dirige, coordina tutto il
  lavoro del Gran Consiglio rivestendo quindi la carica istituzionale
  pi\`u importante;
\item la \textbf{Commissione Governativa} che detiene il potere esecutivo,
  composta da 6 ministri (tesoro, trasporti, corporazioni, esteri,
  difesa, cultura) a cui si aggiunge il Governatore, capo della
  Commissione;
\item la \textbf{Commissione Giudiziaria} di 3 membri il cui compito \`e
  coordinare e dirigere il lavoro delle Corti di Giustizia (potere
  giudiziario) e di tutte le forze di polizia civile (la Guardia).
\end{itemize}
\fi

Compito del Gran Consiglio \`e nominare i due Condottieri
dell'esercito.

\iffullversion
La capitale, ospitante il Consiglio e le Commissioni, \`e la citt\`a di
Dyshu.

Dal punto di vista amministrativo, la Terra di Levante \`e composta da
12 Zone, che godono di ampia autonomia legislativa, avendo propri
Consigli e Commissioni con mandato di durata variabile a seconda della
Zona.

La cittadinanza Levante pu\`o essere ottenuta da chiunque risieda sul
territorio da almeno 10 anni e dai nuovi nati, aventi almeno un
genitore di cittadinanza Levante, entro i confini del territorio. La
cittadinanza stessa pu\`o essere revocata se il cittadino non risiede
pi\`u nel territorio da 10 anni. L'immigrazione \`e regolamentata da
un apposito ufficio dipendente dal Gran Consiglio.

Le Corti di Giustizia hanno il compito di emettere e far eseguire le
condanne per i contravventori al Codice Civile (il testo che raccoglie
tutte le leggi emesse dal Gran Consiglio) o al Codice di Zona (testo
che contiene le leggi locali). \`E prevista la pena di morte per reati
gravissimi (omicidi plurimi, sequestri, alto tradimento) ma viene
rarissimamente comminata.

I titoli nobiliari non danno diritto ad alcun privilegio; anche se,
generalmente, le famiglie nobili possiedono vasti territori e sono
perci\`o ricche.

Le donne godono degli stessi diritti e hanno gli stessi doveri degli
uomini, escluso il servizio di leva.
\fi

\Esercito La Terra di Levante ha un esercito composto da tre settori:
Leva, Professionisti, Mercenari, suddiviso in 12 guarnigioni (una per
zona). Ciascuna delle guarnigioni si avvale di un organico che
comprende elementi dei tre settori.

\iffullversion
La Leva \`e costituita dai giovani cittadini maschi abili (min. 9 in
AGI, FOR, COS) che hanno compiuto 19 anni e vengono obbligatoriamente
richiamati per prestare un anno di servizio, in cui \`e compreso un
periodo di addestramento.

I Professionisti sono soldati che, dopo aver terminato il servizio di
leva, giurano fedelt\`a al Gran Consiglio e alla Popolazione e
intraprendono la carriera militare come professione, dopo aver
sostenuto una prova. Solo i Professionisti possono aspirare ai gradi
pi\`u alti del comando. Tra i Professionisti spiccano i Corpi Speciali
dei Guastatori (fanti specializzati in spionaggio e infiltrazione) e
degli Incursori specializzati in missioni d'assalto.

I Mercenari sono soldati, maschi o femmine, che non hanno la
cittadinanza Levante, assunti e pagati ``a giornata''
quando necessario tramite bando emesso dal Ministro della Difesa.

Le donne non vengono chiamate a prestare il servizio di Leva
obbligatorio, ma sono ammesse come volontarie, potendo, dopo un anno
di addestramento parificato alla leva, accedere al settore
Professionisti.

Il settore pi\`u potente dell'esercito \`e senza dubbio la Flotta, che
come tutto l'esercito della Terra di Levante \`e numeroso, coeso e ben
addestrato. La Flotta riveste importanti funzioni di controllo di
tutto il mare interno, a livello internazionale.

Durante l'addestramento viene insegnato l'uso di tutte le armi,
comprese quelle tipiche degli eserciti degli altri paesi, poich\'e
molti Professionisti sono di origine straniera e tra i Mercenari \`e
facile trovare veterani in grado di svolgere il ruolo di istruttore.

Il comando dell'esercito \`e affidato a due Condottieri, scelti ogni 6
anni dal Gran Consiglio tra gli ufficiali di grado pi\`u alto,
ciascuno dei quali controlla l'operato dell'altro mediante il veto.

Le funzioni di ordine pubblico vengono svolte da un apposito corpo
armato chiamato Guardia, dipendente dalla Commissione Giudiziaria ed
indipendente dall'esercito.

Presso la citt\`a di Thyesi si trova la principale accademia di Arti
Marziali.
\fi


\Economia L'economia Levante \`e basata sull'agricoltura e
sull'allevamento intensivi, sul commercio e sull'artigianato. Tutta la
pianura di Maja \`e fertile e coltivabile; met\`a di essa \`e per\`o
protetta da leggi internazionali volute dal Conclave che impediscono
di disboscare le vaste aree protette (ed in alcuni casi di accedervi).

\iffullversion
Le zone non coltivabili sono comunque utilizzate come territori di
caccia, raccolta di erbe, piante curative, castagne, noci, nocciole e
taglio del legname (in quantit\`a limitata e legata al mantenimento
dello stato di salute dei boschi).

Tra le colture si enumerano cereali, legumi, ortaggi e alberi da
frutto. I pascoli del Meridione e delle Colline Pre-dolane forniscono
cibo sufficiente per grandi mandrie di bestiame (cavalli, bovini) e
greggi di ovini. \`E diffuso anche l'allevamento del fortral (grosso
uccello corridore da sella).

L'artigianato \`e legato ai prodotti della terra e in misura minore
alla lavorazione dei minerali e metalli importati.  Famosi e fiorenti
l'artigianato tessile e calzaturiero, i salumifici e caseifici e in
genere tutte le attivit\`a di manipolazione dei generi alimentari. Sui
monti Dolani, ad alta quota, sono situate miniere di metalli preziosi,
che non possono essere sfruttate al meglio per via del clima freddo.

Le produzioni agricola e artigiana sono molto maggiori del consumo
interno, cos\`{\i}  che gran parte dei prodotti \`e destinata
all'esportazione, soprattutto verso Rylex, Bahuney e Terranova.
L'importazione riguarda soprattutto materie prime di origine
mineraria: metalli, minerali e carbone.
\fi

Il commercio \`e estremamente agevolato dalla morfologia del
territorio, soprattutto dalla presenza di fiumi larghi e navigabili.
Fittissima anche la rete stradale, che supera il confine
settentrionale arrivando sino a Terranova, passando per la Teocrazia
di Ryless.

Cos\`\i\  come la Flotta Navale Militare, anche la Marina Mercantile gode
di meritatissima fama. La pesca \`e praticata nel mare interno.  Il
grande sfruttamento avvenuto in passato ha, purtroppo, impoverito i
fondali costieri rendendo questa attivit\`a molto meno importante e
redditizia delle altre.

Le attivit\`a produttive sono organizzate dalle \textbf{Corporazioni},
associazioni sorte per difendere le categorie di lavoratori e
organizzarne il lavoro. L'iscrizione alle Corporazioni \`e formalmente
facoltativa, ma di fatto obbligatoria, perch\'e altrimenti non si \`e
n\`e protetti, n\`e rappresentati.

\iffullversion
Ogni Corporazione, in cambio di una quota trimestrale variabile,
istituisce Casse di Soccorso e Previdenza per garantire una dignitosa
sopravvivenza ai membri che si infortunano o restano vittime di
imprevisti. Ogni Corporazione garantisce comunque la segretezza sulle
informazioni fornite dai propri iscritti.

Le corporazioni costituiscono degli ``stati nello stato'' paralleli,
con elezioni proprie e lotte per il potere al loro interno. Essere a
capo di una corporazione porta spesso ad un potere maggiore di quello
detenuto dal Capo del Gran Consiglio di Zona. Ad esempio, Il capo
della Corporazione dei Commercianti \`e una delle personalit\`a pi\`u
in rilievo in tutto lo Stato.

Per iscriversi ad una corporazione non \`e richiesta la cittadinanza,
ma serve essere residenti da almeno tre anni. Gli interessi di alcune
corporazioni spesso sconfinano nell'illegalit\`a.
\fi

\Religione La religione di stato, per ragioni storiche, \`e quella
Rileana. Essa viene insegnata nelle scuole e professata dalla maggior
parte della popolazione originaria, anche se non con il rigore tipico
di Novesi e Rileani.

\iffullversion
Molti altri culti di minoranza sono per\`o diffusi e tollerati, a
causa della presenza di numerose etnie sul territorio, tanto che le
religioni sono mischiate fra di loro in un particolare sincretismo.

La personalit\`a religiosa pi\`u carismatica \`e il Vescovo di Tamat
che, nonostante nelle gerarchie ufficiali sia un subalterno del Sommo
Pontefice e del Vescovo di Rylex, nelle Terre Levanti gode di maggior
prestigio rispetto a questi ultimi.
\fi

\Magia La magia \`e legalizzata. Qualsiasi mago iscritto al Conclave
pu\`o esercitare la ``libera professione'' iscrivendosi alla
Corporazione dei Maghi. Le leggi da rispettare sono le stesse imposte
dal Conclave, oltre, naturalmente, alle leggi del Codice Civile e del
Codice di Zona.

\iffullversion
I maghi vengono considerati dalla popolazione alla stessa stregua di
medici, alchimisti, ingegneri e sono frequentemente ammirati e
stimati.  Sovente vengono assunti dalla Guardia per risolvere casi
particolarmente complicati o in cui sono coinvolti altri maghi.
\fi

\`E vietato l'insegnamento della Magia al di fuori delle quattro
Accademie dipendenti dal Conclave. Le scuole insegnate nelle Accademie
sono l'Illusionismo, l'Esistenza e la Stregoneria. L'accademia
principale si trova nella citt\`a di Seemur, ed \`e retta da un
Arcimago.

\iffullversion
L'ingresso alle Accademie \`e subordinato al versamento di una onerosa
tassa (5 Piastre), il cui ricavato \`e diviso tra il Conclave, la
Corporazione dei Maghi, il Gran Consiglio e il Consiglio di Zona.
Questa tassa, insieme alla pesante retta, rende le Accademie
accessibili quasi esclusivamente ai ricchi (0.6 Fattore Istruzione).

La Corporazione dei Maghi mette ogni anno a disposizione dei meno
ricchi che dimostrano capacit\`a fuori dal comune borse di studio
attribuite in seguito alle segnalazioni di un membro della
Corporazione.
\fi

\Moda Le leggi della Terra di Levante consentono una grande libert\`a di
comportamento alla popolazione, essendo ben bilanciate le libert\`a
individuali con i doveri sociali.

Non esistono vere e proprie ``usanze levanti'', dato il continuo
viavai di mercanti di tutti i paesi del Continente che introducono
continui cambiamenti nelle abitudini dei locali.

\iffullversion
Qualunque tipo di attivit\`a che riguardi adulti consenzienti \`e
concesso, purch\'e non disturbi il lavoro, il riposo, il divertimento
e la libert\`a altrui e sia praticato in luoghi privati. \`E vietato
lo sfruttamento, di qualsiasi tipo, dei bambini.

Sono illegali (ma vengono praticati clandestinamente) i giochi
molto violenti ed i combattimenti all'ultimo sangue tra uomini
o umanoidi.

Esistono scuole elementari pubbliche che garantiscono 3 anni di
istruzione gratuita facoltativa ad ogni cittadino che ne faccia
richiesta e scuole professionali e universit\`a private finanziate
dalle Corporazioni.  Chi possiede grossi patrimoni acquisisce
prestigio finanziando scienziati ed artisti, cosicch\'e lo sviluppo
della cultura \`e notevole.

La Terra di Levante \`e uno degli stati pi\`u evoluti e socialmente
avanzati dell'intero continente.

La lingua ufficiale \`e il Levante. Questa \`e anche la lingua
pi\`u diffusa nei commerci e nei rapporti interrazziali.
\fi

\Fisico I Levanti autoctoni sono umani, hanno un'altezza media di 1
metro e 75/1 metro e 80, pelle bianca, carnagione chiara, capelli
castani, occhi solitamente scuri e corporatura media. Sono frequenti
gli incroci con altre razze ed etnie

\minmaxumani

\subsubsection{Classi sociali e abilit\`a di classe sociale}

\classesoc{01-10}{Servi}{0.2}

Vivono nella casa del loro signore, ricevono vitto e alloggio
ricevendo un compenso mensile variabile da 1 a 3 Corone.

\begin{abilist}
\item Calcolare 
\item Cantare o Danzare o Suonare o Cucinare o Abilit\`a professionale collegata
alla mansione 
\item Cavalcare quadrupedi 
\item Galateo 
\item Lingua Madre
\item Lingua Levante
\item Muoversi in silenzio
\end{abilist}

\classesoc{11-27}{Pescatori}{0.3}

Vivono della pesca e dello sfruttamento del mare e dei suoi proventi.
Producono cibo per la propria famiglia e possono guadagnare al netto
delle spese di sostentamento da 1 a 20 scudi al mese.

\begin{abilist}
\item Carpenteria
\item Conoscere animali 
\item Controllo respirazione
\item Fare nodi 
\item Lingua levante 
\item Lingua madre
\item Navigare
\item Nuotare
\item Orientamento 
\end{abilist}

\classesoc{28-44}{Contadini}{0.3}

Coltivano campi ed allevano bestiame. Producono cibo e vestiti per la
loro famiglia e possono guadagnare al netto delle spese di
sostentamento da 1 a 10 scudi al mese.

\begin{abilist}
\item Carpenteria 
\item Cavalcare bipedi
\item Conoscere animali
\item Conoscere piante 
\item Lingua levante
\item Lingua madre
\item Resistenza alla fatica 
\end{abilist}

\classesoc{45-74}{Professionisti}{0.5}

Praticano i pi\`u svariati mestieri e vivono di essi. Possono essere:
artigiani, giuristi, cartografi, erboristi, carpentieri, medici, maghi,
guerrieri, mercanti, macellai ecc. Il loro introito mensile pu\`o
variare da 5 a 30 corone al mese.

\begin{abilist}
\item Abilit\`a professionale \#1
\item Abilit\`a professionale \#2
\item Abilit\`a professionale \#3 
\item Lingua levante 
\item Lingua madre
\item Altre abilit\`a a scelta del Master 
\end{abilist}

\classesoc{75-90}{Benestanti}{0.6}

Praticano molteplici mestieri. Possono essere proprietari di grandi
esercizi commerciali, come alberghi, gioiellerie, medi allevamenti e
via dicendo, che forniscono loro gran parte del reddito, variabile tra
15 e 300 corone al mese.

\begin{abilist}
\item ATL o ADT 
\item Calcolare 
\item Cavalcare bipedi o quadrupedi 
\item Danzare 
\item Galateo 
\item Lingua levante 
\item Lingua madre 
\item Lingua straniera 
\item Mitologia 
\item Religione 
\item Storia 
\item Trattare 
\item Valutare
\end{abilist}

\classesoc{91-100}{Borghesi ricchi}{0.7}

Sono imprenditori di molteplici attivit\`a ed hanno molti interessi,
in cui investono il proprio capitale. Sono perlopi\`u banchieri,
armatori, grossi allevatori, possessori di miniere, mercanti di
schiavi, ecc.  Ricavano ogni mese da 100 a 2000 corone.

\begin{abilist}
\item ADT 
\item Calcolare 
\item Cavalcare bipedi o quadrupedi 
\item Danzare 
\item Galateo 
\item Lingua levante 
\item Lingua madre 
\item Lingua straniera 
\item Mitologia 
\item Religione 
\item Storia 
\item Trattare 
\item Valutare
\end{abilist}

\filler{corona.eps}

\iffullversion
\subsection{La Teocrazia di Ryless}

\Stendardo{stendardo_rylex.eps}

I Rileani occupano la Teocrazia di Ryless, detta anche Il Pontificato.

\Geografia La Teocrazia di Ryless si estende nella fascia centrale
nell'isola di Dyshu-Thandy. Confina a nord con Terranova e a Sud con
la Terra di levante tramite la Catena Muraria che agisce da confine
naturale. 

La costa Ovest \`e percorsa dai Monti Dolani la cui cima pi\`u alta
nel territorio della Teocrazia \`e Monte Nay con i suoi 5560 metri. Il
territorio degrada poi dolcemente verso est fino a creare la piana di
Sowsathow. 

Il fiume pi\`u importante \`e Sarhyu che segna il confine con
Terranova ed \`e totalmente navigabile nel suo corso in pianura.  Non
trascurabile \`e anche il Ryo Mothow anch'esso navigabile che sfocia
nel Mare Interno dando vita alla Baia di Keeah dove sorge il Porto di
Madaow.

Le foreste sono prevalentemente popolate da conifere, ma qualche bosco
di latifoglie \`e presente all'estremo Sud, a ridosso della catena
Muraria. 

La capitale \`e Leem-Bara, ma la citt\`a pi\`u importante \`e Rylex,
arroccata sui monti Dolani e circondata da mura alte sino a 30 metri.
Altre citt\`a importanti sono: Y'nhoy, Casthya, Y'nhuee e Aquileia.

\Politica Nella Teocrazia di Ryless le supreme autorit\`a sono i
vescovi, uno per ognuno dei sette vescovadi in cui \`e suddiviso il
territorio. 

I vescovi avocano a s\`e tutto il potere, ma lo delegano
per il normale esercizio ad un governatore; questi \`e nominato dal
vescovo fra i membri delle famiglie pi\`u ricche ed influenti (e che
elargiscono pi\`u offerte al vescovado).  

Il governatore amministra tutti gli affari correnti ad eccezione
dell'ordine pubblico che \`e gestito dal corpo delle guardie vescovili
che rendono conto direttamente al vescovo.

Il governatore resta in carica per 3 anni (salvo revoca del vescovo) e
pu\`o essere rinominato. \`E facile capire perci\`o che il grado di
servilismo verso il Vescovo da parte delle famiglie pi\`u ricche \`e
elevatissimo.

Il governatore ha il compito di trattare con le parti sociali
(benestanti, borghesi ricchi e professionisti) e raggiunto l'accordo
deve presentarlo al Vescovo per la ratifica. Se essa non viene
concessa, il governatore \`e costretto a riaprire le trattative. 

La quasi totalit\`a del Clero \`e composta da membri delle famiglie
influenti, poich\'e di solito i primogeniti vengono affidati alla cura
dei seminari fra gli 8 e i 13 anni.

Il potere giuridico \`e amministrato per i reati comuni dal
governatore che si avvale di tre giureconsulti clericali; per i reati
pi\`u gravi la giurisdizione viene rimandata al Vescovo.
l'amministrazione periferica del vescovado \`e affidata ai Capitolari,
autorit\`a a capo delle parrocchie (cio\`e le unit\`a amministrative
in cui \`e suddiviso il territorio di ogni vescovado).  I Capitolari
riferiscono direttamente al Vescovo.

\Esercito L'esercito rileano \`e un esercito compatto e pesante
finanziato al 90\% dalle Casse Episcopali e al 10\% dalle famiglie
ricche. L'esercito \`e composto per 1/3 da chierici e per 2/3 da
soldati tutti addestrati dall'et\'a di 15 anni fino al compimento dei 
23 anni in accademie militari
religiose.  Spesso la carriera militare \`e l'unica via per
affrancarsi da situazioni economiche di estrema povert\`a. I candidati
laici sono scelti in base a dure selezioni (e a forti
raccomandazioni). L'esercito \`e organizzato in 3 divisioni: fanteria
leggera, fanteria pesante e cavalleria.

La prima divisione impiega solitamente spade, mazze, asce e scudi.
La seconda divisione armi pesanti quali: spadoni a due mani, asce da
battaglia, morning star, alabarde e picche.  La terza divisione scudi
grandi, spade e lance. 

Capo dell'esercito \`e il comandante crociato nero (capitano delle
guardie di Rylex) che riferisce direttamente al Consiglio episcopale.
Suoi luogotenenti i comandanti crociati degli altri Vescovadi.

Ogni Guardia Vescovile, corpo di fanteria pesante,
porta sopra l'armatura una tunica crociata (la croce \`e la
stilizzazione del pugnale che uccise il profeta Ryless). 

I colori delle tuniche caratterizzano i Vescovadi: Rylex nera, Madaow
grigia, Leem-Bara viola, Y'nhoy giallo, Aquileia azzurro, Casthya
arancione, Y'nhuee marrone. Il colore delle croci indica la gerarchia:
blu: comandante, rosso: tenente, verde: sergente, bianco: guardia
semplice.

\Economia L'economia rileana \`e molto fiorente sia all'interno, che
nei traffici con gli altri regni.  Ciononostante quasi tutta la
ricchezza \`e incamerata dalle Casse Episcopali (tramite i Gabellieri
che riscuotono le imposte).

Ogni famiglia (escluse quelle dei Borghesi Ricchi, che si limitano ad
offerte di solito cospicue) \`e soggetta a 3 imposte: l'imposta
ecclesiastica del 10\%, l'imposta vescovile del 20\% e la tassa sul
reddito del 5\%.  L'imposta ecclesiastica viene destinata per il 70\%
al Papato e per il restante 30\% rimane nelle casse Episcopali insieme
all'imposta vescovile.  La tassa sul reddito viene destinata al
governatore. Il divario di reddito fra le classi medio-basse e le
ricche famiglie \`e notevole.

La Teocrazia produce legname, miele e filati; alleva ovini, bovini,
equini, suini, fortral e nortral. Sono diffuse anche la pesca
fluviale, le colture di ortaggi e la fabbricazione di monili e
gioielli.

La Teocrazia esporta miele e gioielli e importa metalli da cui
ricavare le armature e le armi. Una volta l'anno, nella citt\`a di
Rylex si svolge la pi\`u importante fiera mercantile delle terre
conosciute. In occasione di essa la citt\`a si riempie di commercianti
provenienti da ogni parte dell'Arcipelago.

\Religione La religione ha un ruolo fondamentale nella vita di
ogni rileano. Il rileanesimo prende il nome dal profeta Ryless
considerato il Messia, il Salvatore che avrebbe scacciato la
malvagit\`a da tutte le terre conosciute.

\filler{pugnale_ryless.eps}

Al tempo in cui egli diffondeva la sua parola, il culto pi\`u diffuso
era l'adorazione degli Dei del Sole e delle Lune, del Mare e della
Terra.  I sacerdoti di questo culto vedevano Ryless ed il consenso che
riscuoteva fra la gente come una minaccia e decisero cos\`{\i}  di
eliminarlo; ma facendolo alla luce del sole avrebbero rischiato di
inimicarsi la popolazione.

Contattarono un suo seguace, Burdhus e gli offrirono ben 15 corone per
uccidere Ryless. Burdhus si rifiut\`o e riusc\`{\i} a sfuggire
miracolosamente all'ira dei sacerdoti.  Essi per\`o non si arresero,
rapirono la moglie di Burdhus e lo costrinsero a sottomettersi ai loro
voleri in cambio della vita della sua consorte. Fu cos\`{\i}  che una
notte Burdhus si arm\`o di un pugnale e si accost\`o al capezzale di
Ryless.  Ryless si dest\`o, la leggenda narra che Burdhus pianse e che
Ryless gli disse: --So perch\'e sei qui e non ho rancore. Uccidimi e
trasmetti al mondo la mia parola--  Burdhus conficc\`o la lama nel
cuore di Ryless e con le lacrime agli occhi maledisse i sacerdoti. Tre
giorni dopo il tempio dei sacerdoti croll\`o e li seppell\`{\i} tutti.
Burdhus divenne il ``Portatore di Luce'', colui che svel\`o al popolo
le verit\`a di Ryless.

L'assoluto capo religioso \`e il Sommo Pontefice che vive nella
capitale Leem-Bara e decide lo svolgersi della vita religiosa, dei
riti e delle festivit\`a. Il calendario di questi eventi \`e deciso
collegialmente con i vescovi entro la prima decade dell'anno.  Il
Pontefice viene eletto, a vita, dai vescovi, che bench\'e scelgano i
candidati fra tutti gli abitanti, spesso fanno risultare vincitore un
appartenente all'alto Clero.

Le feste principali sono Paskha, il 9 di Kallentee che \`e il giorno
della Morte di Ryless e Paskha-Manha il 25 di Gehlow, il giorno della
nascita di Ryless. Numerose festivit\`a sono poi consacrate ai
principali santi e si celebrano in date stabilite.
Festa comandata di ogni settimana \`e il giorno di Deebhona ed \`e
obbligatorio per tutti i fedeli partecipare alla liturgia.

Ognuna delle parrocchie in cui \`e diviso il vescovado, affidata
all'amministrazione dei Capitolari, comprende pi\`u chiese
amministrate da un Officiante che si occupa della celebrazione dei
riti e della lettura dei testi sacri. 

I Capitolari si riuniscono in un collegio detto Capitolo Cattedrale
ogni qualvolta il Vescovo lo ritenga opportuno. Essi rappresentano il
senato religioso del vescovado con compiti consultivi in materia di
fede.

I morti vengono sepolti con appositi riti all'interno di cimiteri.

I religiosi sono vincolati dall'obbligo di celibato.

\Magia In tutto il pontificato l'uso dell'arte magica \`e severamente
vietato, eccezion fatta per la scuola di Esistenza, considerata
l'unica magia buona, per esercitare la quale \`e comunque necessaria
l'autorizzazione da parte di un vescovo.  La magia dell'Esistenza
viene largamente utilizzata per proteggere l'esercito in guerra; essa
\`e insegnata solo nei seminari.

La presenza del Conclave \`e fortemente osteggiata e di conseguenza
molto scarsa; solo poche citt\`a di frontiera vedono la presenza di un
delegato.

L'Anticonclave, data la sua non ufficialit\`a, vanta pi\`u di una
sede nel territorio, ma si guarda bene dal fomentare disordini tramite
l'uso della magia come accade invece in altri regni.  

A controllare che non venga praticata la magia \`e preposto il corpo
degli inquisitori, costituito da chierici nominati dal vescovo. Il
corpo \`e presente in ogni vescovado ed \`e costituito 9 membri,
fra essi uno \`e il capo ed ha il titolo di Grande Inquisitore.
Compito degli inquisitori \`e reprimere e punire severamente l'uso
della magia proibita; le pene vanno dalla fustigazione, in caso di
incantesimi che non provocano danni, al rogo per le infrazioni pi\`u
gravi.

Gli inquisitori sono insieme giudici e carnefici e non hanno bisogno
di autorizzazioni particolari per condannare a morte una persona. Di
solito gli inquisitori, eccezion fatta per il Grande Inquisitore, si
dividono in gruppi di due e pattugliano le strade, soli durante il
giorno e insieme alle guardie vescovili la notte. Sono vestiti con un
saio nero stretto in vita da una fune rossa. Le loro armi predilette 
sono il pugnale e la spada.

\Moda Il pontificato non garantisce pienamente le libert\`a
individuali. La libert\`a di espressione \`e tollerata sinch\'e non
offenda la religione. La libert\`a sessuale \`e limitata, di
conseguenza \`e reato gestire bordelli.

A questo riguardo le prostitute ricorrono solitamente all'adescamento
nelle taverne. Il reato di prostituzione \`e punito con la
fustigazione, ma non \`e infrequente che si proceda anche alla
mutilazione di dita e orecchie. Il gioco d'azzardo \`e proibito e
punito col carcere; se al reato di gioco d'azzardo si somma anche
l'imbroglio, la pena consta anche dell'amputazione del pollice.

L'unico gioco d'azzardo tollerato \`e la scommessa sugli incontri di
Skownha-Kownhoo, un gioco che si pratica in un campo rettangolare di
40x15 dove due squadre di 6 elementi ciascuna devono contendersi il
cranio trattato e imbalsamato dei criminali giustiziati, per poi
infilzarlo su una punta metallica posta nel lato corto del campo
avversario. La partita dura quaranta minuti ed \`e proibito l'uso
della magia e delle armi da taglio.  Solitamente si utilizzano
particolari armi da botta lunghe che inabilitano senza causare
lacerazioni; sono comunque ammesse armi come catene e bastoni. Il
campionato ufficiale \`e composto da 12 squadre e l'iscrizione costa
150 corone. La squadra vincitrice del campionato intasca 1500 corone
(pi\`u i soldi eventualmente scommessi).

Poich\'e per il clero \`e vantaggioso tenere il popolo nell'ignoranza,
le scuole, gestite dalla chiesa, hanno rette altissime (sino a 5
corone al mese).  Le classi povere istruiscono in famiglia la prole.
L'insegnamento non \`e libero, tant'\`e che \`e considerato reato
tenere lezioni collettive. 

La lingua ufficiale \`e il Levante.

\Fisico I Rileani sono umani, hanno un'altezza media di 1 metro e 75/1
metro e 80, pelle bianca, carnagione chiara, capelli chiari, occhi
solitamente chiari e corporatura media.

\minmaxumani


\subsubsection{Abilit\`a di razza}

\begin{abilist}
\item Resistenza al freddo 2
\end{abilist}

\subsubsection{Classi sociali e abilit\`a di Classe sociale}

\classesoc{01-10}{Servi}{0.2}

Vivono nella casa del signore, ricevendo vitto e alloggio insieme ad
un compenso mensile variabile da 1 a 3 Corone.

\begin{abilist}
\item Calcolare
\item Cantare o Danzare o Suonare o Cucinare o abilit\`a professionale
  collegata alla mansione
\item Cavalcare quadrupedi
\item Galateo
\item Lingua Madre
\item Lingua Levante
\item Muoversi in silenzio
\end{abilist}

\classesoc{11-40}{Contadini}{0.3}

Coltivano campi ed allevano bestiame.  Producono cibo e vestiti per la
propria famiglia e possono guadagnare al netto delle spese di
sostentamento da 1 a 10 scudi al mese.

\begin{abilist}
\item Carpenteria
\item Cavalcare bipedi
\item Conoscere animali
\item Conoscere piante
\item Lingua levante
\item Lingua madre
\item Resistenza alla fatica
\end{abilist}

\classesoc{41-70}{Professionisti}{0.5}

Praticano i pi\`u svariati mestieri e vivono di essi. Possono essere:
artigiani, giuristi, cartografi, erboristi, carpentieri, medici,
guerrieri, sacerdoti, macellai ecc. Il loro introito mensile pu\`o
variare da 5 a 30 corone al mese.

\begin{abilist}
\item Abilit\`a professionale \#1
\item Abilit\`a professionale \#2
\item Abilit\`a professionale \#3
\item Lingua levante
\item Lingua madre
\item Altre abilit\`a a scelta del Master
\end{abilist}

\classesoc{71-90}{Benestanti}{0.6}
 
Praticano molteplici mestieri. Hanno di solito grandi esercizi
commerciali, come alberghi, gioiellerie, medi allevamenti e via
dicendo, che gli forniscono gran parte del reddito, variabile tra 100
e 1000 corone al mese.



\begin{abilist}
\item ATL o ADT
\item Calcolare
\item Cavalcare bipedi o quadrupedi
\item Danzare
\item Galateo
\item Lingua levante
\item Lingua madre
\item Lingua straniera
\item Mitologia
\item Religione
\item Storia Generale
\item Trattare
\item Valutare
\end{abilist}

\classesoc{91-100}{Borghesi ricchi}{0.7}

Sono proprietari di molteplici attivit\`a, in cui investono il proprio
capitale. Sono perlopi\`u banchieri, armatori, grossi allevatori,
possessori di miniere, mercanti di schiavi, ecc.  Ricavano al mese da
200 a 2500 corone.

\begin{abilist}
\item ATL o ADT
\item Calcolare
\item Cavalcare quadrupedi
\item Danzare
\item Lingua levante
\item Lingua madre
\item Lingua straniera
\item Religione
\item Storia Generale
\item Trattare
\item Valutare
\end{abilist}


\subsection{Terranova}

\Stendardo{stendardo_terranova.eps}

Terranova \`e abitata dagli umani Novesi.

\Geografia Il Regno di Terranova occupa tutta la parte occidentale
dell'isola di Norda, confinando a Est con le Terre del Nord ed a
Sud-Est con Mahn-Hesa ed occupando inoltre la parte settentrionale
dell'isola di Dyshu-Thandy. Le coste dell'isola che fanno parte del
Regno sono frastagliate e ricche di fiordi, molti dei quali,
soprattutto a Settentrione, sono sovrastati da ghiacciai. Il clima \`e
artico nell'estremo nord, mentre a Sud la temperatura \`e appena pi\`u
mite.  Nell'entroterra, invece, assume caratteristiche tipicamente
continentali, con estati tiepide e secche e inverni gelidi e molto
nevosi.

Il territorio \`e prevalentemente montuoso o collinare se si
eccettuano le pianure costiere e fluviali del meridione; i fiumi sono
brevi e hanno pendenze medie relativamente alte, nascono solitamente
da sorgenti di tipo glaciale e sfociano ad estuario.  Generalmente i
fiumi che sfociano nel mare interno sono per\`o pi\`u larghi e
lenti, consentendo la navigazione.

La vegetazione naturale \`e formata in gran parte da foreste di
conifere, in cui predomina la tipica sagoma a forma di cono
dell'imponente abete novese; nelle zone pi\`u secche e pianeggianti
\`e diffusa la tundra. 

La cima pi\`u alta, facente parte del Massiccio centrale del Dehrberg,
al confine con le Terre del Nord, \`e il monte Kahlfahr, che raggiunge
i 3520 metri. 

Il fiume pi\`u importante, il Nother, \`e lungo circa 250 km, nascendo
proprio dal Massiccio del Dehrberg. Altra catena montuosa di rilevante
importanza \`e la Catena delle Sei Cime (Hesfahrts), all'estremo Nord,
la cui cima pi\`u alta, il la Terza Montagna (Tihrfahr), \`e alta 2200
metri.

\Politica Terranova \`e uno stato ad organizzazione tipicamente
feudale, il cui territorio \`e diviso in tre Ducati (suddivisi a loro
volta gerarchicamente in Marche, Contee, Baronie) governati ciascuno
da una potente famiglia di nobili: Ducato degli Arethin, Ducato degli
Homer, ducato dei Van Metz.

L'organo di governo pi\`u importante \`e il Consiglio delle Famiglie,
costituito da nove rappresentanti, tre per ciascuno dei Ducati, che si
occupa di quanto riguarda la politica estera, le guerre e l'elezione
del re. All'interno di ciascun Ducato ogni Duca, capofamiglia, \`e di
fatto un monarca assoluto, pur dovendo rispondere al Consiglio delle
Famiglie di eventuali questioni che possano riguardare anche uno solo
degli altri ducati.

Il capo dello Stato \`e il Re, che viene nominato dal Consiglio
delle Famiglie e rimane in carica a vita. Il Re non ha di fatto poteri
e funge quasi esclusivamente da ambasciatore e rappresentante
diplomatico. Egli viene scelto tra i pi\`u importanti esponenti di
queste tre famiglie, ma di fatto viene scelto da quella pi\`u
potente e pi\`u ricca, il cui Capofamiglia \`e colui che detiene
concretamente il potere.

Le famiglie meno importanti sono organizzate gerarchicamente,
facendo capo ad una delle tre maggiori. Capita spesso che, durante le
guerre, alcune famiglie cambino ``sfera di influenza'', spostandosi,
nei limiti dovuti alla posizione geografica, da un Ducato all'altro
modificando in questo modo l'equilibrio politico nel Consiglio delle
Famiglie. 

All'interno di ogni feudo, il Signore ha diritto di vita e di morte
sui suoi sudditi, e i suoi diritti e poteri hanno come unico limite
l'obbligo di non creare disturbi ai feudi vicini o ai propri
superiori. Capita quindi che, per evitare rivolte e sommosse che
potrebbero uscire dai confini, i Signori preferiscano non abusare del
loro potere.


Il titolo nobiliare, cos\`{\i} come il feudo, si trasmette di padre in
figlio in ``linea diretta'' (legge salica). I figli cadetti vengono in
genere indirizzati a carriere militari o ecclesiastiche.

I Signori o i loro rappresentanti amministrano la giustizia mediante
le pubbliche udienze alle quali chiunque pu\`o essere ammesso, previa
iscrizione ad un'apposita lista. Gli unici codici sono quelli che
stabiliscono i poteri del Re e quelli del Consiglio delle Famiglie,
mentre non esistono leggi scritte che regolino la vita pubblica,
lasciando spesso il popolo in bal\`{\i}a dei capricci dei signori.

Bench\'e non godano dei diritti politici, i ricchi borghesi vengono
trattati con riguardo, visto che dalla loro ricchezza (e dalle tasse
da loro pagate) dipende la potenza dello stato, cos\`{\i}  che i
commerci, le imprese e l'artigianato sono incentivati con favori,
lavori pubblici (strade, ponti) e concessioni di terre. Chi ci
rimette, in tutto questo, \`e ovviamente la classe inferiore,
costituita dai piccoli agricoltori, braccianti e per ultimi da servi
della gleba, che sono poco pi\`u che puri e semplici animali da soma.

Le donne della nobilt\`a non hanno diritti politici e non possono
trasmettere i titoli nobiliari, essendo quasi sempre subalterne dei
loro padri o mariti.  Le donne borghesi sono pi\`u emancipate, ma agli
occhi dei nobili sono esseri inferiori e capricciosi, perci\`o non
affidabili; capita perci\`o raramente di vedere una donna a capo di
una grossa ``impresa''.

La capitale, residenza del Re, \`e la citt\`a di Rambleen.  Le citt\`a
pi\`u importanti dei tre ducati, dove si trovano le corti dei Duchi di
Arethin, Homer, Van Metz sono A\-reth, Hoss, Gott\-metz,
rispettivamente.

\Esercito L'esercito di Terranova gode di ottima fama per quanto
riguarda le capacit\`a di combattimento. La sua organizzazione
riflette strettamente la struttura politica dello stato. L'esercito,
unitario, \`e diviso in tre parti, una per ciascuna delle famiglie,
ognuna delle quali provvede al reclutamento e all'addestramento dei
propri uomini, presi a loro volta dagli eserciti e dai plotoni dei
propri sottomessi. 

Uno dei tre Capifamiglia nominato dal Consiglio delle Famiglie, in
caso di guerra, diviene Comandante in Capo dell'esercito unificato. In
genere viene eletto il pi\`u capace, ma molte volte i giochi di potere
fanno s\`\i\  che venga eletto il pi\`u ``potente''. Il Comandante in
Capo pu\`o a sua volta nominare un Generale di sua fiducia a cui pu\`o
rimettere il comando. Questa operazione viene svolta molto di
frequente, poich\'e non sempre i Duchi sono esperti di tattica
militare e quindi non sono in grado di guidare un grosso esercito.

La forza dell'esercito \`e dovuta in prevalenza al corpo dei Cavalieri
di Ferro, costituita da un'\'elite di cavalieri in armatura di piastre
che conducono cavalli anch'essi dotati di armatura di piastre.

Chiunque dimostri di avere sufficiente destrezza o intelligenza
pu\`o entrare nell'esercito, compresi i servi della gleba, ma solo i
nobili possono aspirare ai gradi pi\`u alti; mentre, essendo le
armature e le cavalcature a spese di chi le usa, solo i pi\`u ricchi
possono far parte della Cavalleria di Ferro.

Avere un membro della famiglia in questo corpo \`e generalmente
motivo di grande orgoglio e prestigio. I Cavalieri di Ferro sono
investiti direttamente dal Re in una solenne cerimonia che avviene
ogni due anni.

I Cavalieri di Ferro che si distinguono in modo particolare durante le
guerre possono essere investiti di un titolo nobiliare qualunque sia
la loro origine.

Esistono accademie in cui i giovani cadetti imparano le arti del
comando e l'uso delle armi tipiche: lo Spadone a due mani, la Spada
Bastarda, la Spada a una Mano, l'Alabarda, la Mazza Ferrata, la Lancia
e la Picca. 

L'esercito svolge in tempo di pace le mansioni di ordine pubblico.
Esiste inoltre un corpo speciale, dipendente direttamente dal Duca di
Arethin, costituito da mercenari stranieri, sul quale vige il pi\`u
rigoroso riserbo. Le donne non sono ammesse a far parte dell'esercito.

\Economia Chi tiene le redini dell'economia \`e la classe dei
borghesi, dai quali i Nobili dipendono in modo quasi parassitario; i
nobili non gestiscono infatti direttamente i loro possedimenti (anzi,
considerano disonorevole occuparsene), ma vendono le concessioni al
miglior offerente in cambio di una ``dote'' iniziale pi\`u una
percentuale variabile tra il 10\% e il 60\% del prodotto. I
commercianti e gli artigiani invece sono costretti a pagare una tassa
fissa del 30\% su tutti i loro introiti.  Ogni feudatario, inoltre,
paga una ``tassa'' fissa annuale al suo diretto superiore.
La ``cassa'' del feudo e la ``borsa'' del feudatario coincidono.

Alla morte del concessionario oppure allo scadere della concessione
(generalmente della durata di 25 anni), viene solitamente bandita
un'asta. I concessionari possono servirsi liberamente dei servi della
gleba come manodopera a costo zero (i servi fanno parte della
concessione) ma possono anche pagare (poco) braccianti liberi o
subaffittare la concessione a piccoli coltivatori, allevatori o
minatori. Ogni servo della gleba \`e per\`o messo nelle condizioni di
poter avere un suo ``gruzzolo'' personale e poter comprare a caro
prezzo la sua libert\`a.

Il piatto forte dell'economia Novese \`e costituito dall'industria
mineraria: i monti di Dehrberg sono infatti praticamente costituiti da
minerali di ferro e giacimenti di carbone, inframmezzati da piccoli
giacimenti di metalli preziosi. I minerali vengono lavorati sul posto
da abili artigiani oppure esportati.  l'abilit\`a degli artigiani
giustifica l'alto costo delle armi e delle armature decorate,
esportate in tutto il Continente.  I trasporti avvengono per lo
pi\`u per via fluviale mediante alaggio lungo i maggiori fiumi
navigabili. Gran parte del mercato del carbone \`e costituito dalle
esportazioni verso Loydi-Genya (terra degli Gnomi).

Le maggiori miniere risiedono nel territorio degli Arethin
giustificandone in gran parte la loro secolare potenza. l'agricoltura
e l'allevamento costituiscono una buona parte degli introiti dei
feudatari. Al Sud si allevano cavalli, bovini e capre, al nord sono
importanti l'allevamento della renna e del Nortral (grosso uccello
bipede simile allo struzzo, utilizzabile come cavalcatura).
Relativamente poco importanti sono la pesca delle balene nei mari del
Nord-Ovest e del pesce azzurro nei mari meridionali. 

\Religione I Novesi professano rigidamente la religione Rileana,
riconoscendo il Sommo Pontefice come capo supremo della Chiesa e il
Vescovo di Rylex come massima autorit\`a esecutiva. Sono diffusissime
le chiese di questa religione (ogni paese ne ha almeno una) e le
abbazie. Le altre religioni non sono tollerate.  L'appartenenza a
religioni diverse \`e usato per giustificare talvolta la
discriminazione di alcune categorie di persone: i mercanti orientali
che non accettano di convertirsi, ad esempio, pagano ``sottobanco''
una tassa aggiuntiva ai feudatari per non essere perseguitati.

Tutti i fedeli, cio\`e tutti i Novesi, sono tenuti a seguire la
Cerimonia ogni Giorno Sacro (il primo, l'undicesimo e il ventunesimo
giorno di ogni mese) e a versare un obolo, per quanto piccolo, alla
Chiesa. Chi si assenta dalla Cerimonia senza una valida
giustificazione \`e costretto a pagare una multa, che viene divisa
tra la Chiesa e il feudatario, e ad una penitenza di tipo corporale.

La magia \`e vista con sospetto dai religiosi ma viene tollerata dalla
maggior parte dei Signori, nonostante chi venga sorpreso a recitare
incantesimi in pubblico possa essere arrestato e interrogato.  Tutti
gli uomini e le donne liberi possono accedere a ordini monastici,
mentre, per ragioni economiche, la carriera ecclesiastica \`e
riservata ai ceti pi\`u alti (min. 0.4 Fattore Istruzione).

I religiosi cercano in tutti i modi di preservare almeno all'apparenza
l'integrit\`a morale della popolazione, condannando apertamente (ma
non sempre) la prostituzione, l'adulterio, i rapporti prematrimoniali,
la mancata osservanza delle regole religiose.

\Magia Terranova \`e l'unico Stato in cui possono essere insegnate
tutte e sei le Scuole, ma siccome ognuno dei Ducati ne ospita soltanto
due, non facenti parte dello stesso regno, questo limita di fatto il
potere della maggior parte dei maghi di basso rango. 

Nel ducato degli Homer vengono insegnate le scuole di Stregoneria ed
Illusionismo; nel Ducato degli Arethin sono insegnate Elementalismo e
Necromanzia; nel Ducato dei Van Metz sono insegnate Demonologia ed
Esistenza.

La principale Accademia di Magia si trova a Barlina ed \`e retta da
un Arcimago. 

Le scuole dipendono direttamente dal Conclave, sono accessibili, per
via delle alte rette, soltanto a ricchi e nobili (min. 0.5 Fattore
Istruzione), anche se questi ultimi considerano molto pi\`u
prestigiosa una carriera di tipo ecclesiastico o militare. Chi si
iscrive ad una scuola viene immediatamente registrato in apposite
liste presso i palazzi Ducali.

\`E vietato lanciare incantesimi in pubblico se non per ``ragioni
validissime'' (la cui validit\`a \`e valutata, ovviamente, dal Signore
del posto), e ci si preoccupa di tenere il popolino il pi\`u possibile
all'oscuro dell'esistenza della Magia.

\Moda Le attivit\`a di base sono tipiche della classe sociale di
appartenenza, cos\`{\i}  che essendo le classi praticamente chiuse,
sono sempre gli stessi gruppi di persone che svolgono gli stessi
ruoli. I Nobili di alto rango, in tempo di pace, si preoccupano di
gestire le concessioni ai borghesi e i rapporti tra feudatari, oltre
che di sbrigare le faccende giuridiche, ma non dedicano a queste
attivit\`a che la mattina, avendo cos\`{\i}  a disposizione tutto il
pomeriggio e la sera per svolgere attivit\`a di svago.

 Tra i divertimenti pi\`u alla moda si enumerano la caccia
alla volpe, la caccia col falcone, la scherma simulata, i balli, le
rappresentazioni teatrali e musicali e le feste.  Pochi Signori amano
circondarsi di poeti, musici, scienziati, dando cos\`{\i}  una piccola
spinta allo sviluppo artistico-letterario e intellettuale. Sono
pochissimi invece i Signori che dedicano il pomeriggio alle
attivit\`a culturali e di studio. In tempo di guerra i Nobili
occupano alte cariche di comando nell'esercito, rinunciando cos\`{\i} 
ai divertimenti, escluse le notti ``da osteria''. 

I borghesi hanno generalmente meno tempo da dedicare a se stessi,
poich\'e sono costantemente impegnati nella gestione delle concessioni
o nella direzione delle ``imprese'' artigiane, ma tendono a curare di
pi\`u l'aspetto ``culturale'' della loro vita, studiando musica,
filosofia o arti magiche.

A causa della presenza della Chiesa Rileana i rapporti tra uomini e
donne sono rigidamente regolati da severe e restrittive norme di
comportamento, che impediscono gli incontri tra ragazzi non fidanzati,
i rapporti prematrimoniali e qualunque tipo di promiscuit\`a
sessuale.

Alcuni Signori usufruiscono ancora del Diritto della Prima Notte,
bench\'e questo costume sia condannato e stia ormai cadendo in disuso.

Quasi nessuno svago \`e riservato alle classi pi\`u disagiate
(servi e braccianti).

I Novesi si vestono con pesanti panni abbastanza sobri, usano
moltissimo il cappello a tesa larga e il mantello per uscire quando
c'\`e freddo. I Nobili esibiscono volentieri l'effigie della loro
famiglia, sia sugli abiti che sulle cavalcature e nelle loro
abitazioni.

La loro cucina \`e molto grassa e ricca di calorie per via dei
freddi inverni; i cibi pi\`u prelibati sono a base di carne di
animale allevato o, meglio ancora, di cacciagione. l'arte culinaria,
comunque, non \`e molto elaborata. La lingua ufficiale \`e il
Novese.

\Fisico I Novesi sono di razza umana, generalmente alti, robusti ed
asciutti (1.80 di altezza media), di carnagione chiara, naso dritto e
stretto, fronte alta e labbra sottili, occhi solitamente chiari e
capelli castani. Dato il clima, non \`e infrequente incontrare persone
sovrappeso.

\minmaxumani

\subsubsection{Abilit\`a di razza}

\begin{abilist}
\item Resistenza al freddo 8
\end{abilist}

\subsubsection{Classi sociali e abilit\`a di Classe sociale}

\classesoc{01-07}{Servi della gleba}{0.1}

Non vantano alcuna libert\`a, e devono lavorare per tutta la vita alle
dipendenze del loro padrone. I figli di un servo della gleba
nasceranno servi della gleba.  I padroni pi\`u benevoli possono
concedere ai servi della gleba piccole somme mensili con cui potersi
un giorno affrancare.

\begin{abilist}
\item Cantare
\item Conoscere Piante
\item Lingua Madre
\item Lingua Levante
\item Resistenza alla fatica
\item Resistenza al freddo
\item Resistenza alle malattie
\item Sollevare pesi
\end{abilist}


\classesoc{08-15}{Servi}{0.2}

Vivono nella casa del proprio signore. Ricevono vitto e alloggio e un
compenso mensile variabile da 1 a 3 corone.

\begin{abilist}
\item Calcolare
\item Cantare o Danzare o Suonare o Cucinare o abilit\`a professionale
  collegata alla mansione
\item Cavalcare quadrupedi
\item Galateo
\item Lingua Madre
\item Lingua Levante
\item Muoversi in silenzio
\end{abilist}

\classesoc{16-22}{Contadini}{0.3}

Producono cibo e vestiti per la propria famiglia, coltivano i campi ed
allevano bestiame. Possono guadagnare al netto delle spese di
sostentamento da 1 a 10 scudi al mese.

\begin{abilist}
\item Carpenteria 
\item Cavalcare bipedi
\item Conoscere animali 
\item Conoscere piante
\item Lingua levante
\item Lingua madre
\item Resistenza alla fatica
\end{abilist}

\classesoc{23-29}{Pescatori}{0.3}

Vivono della pesca e dello sfruttamento del mare e dei suoi proventi.
Producono cibo per la propria famiglia e possono guadagnare al netto
delle spese di sostentamento da 1 a 20 scudi al mese.

\begin{abilist}
\item Carpenteria
\item Conoscere animali
\item Controllo respirazione
\item Fare nodi
\item Lingua levante
\item Lingua madre
\item Navigare
\item Nuotare
\item Orientamento
\end{abilist}

\classesoc{30-55}{Professionisti}{0.5}

Praticano i pi\`u svariati mestieri e vivono di essi. Possono essere:
artigiani, giuristi, cartografi, erboristi, carpentieri, medici,
macellai ecc. Il loro introito mensile pu\`o variare da 5 a 30
corone al mese.

\begin{abilist}
\item Abilit\`a professionale \#1
\item Abilit\`a professionale \#2
\item Abilit\`a professionale \#3
\item Lingua levante
\item Lingua madre
\item Altre abilit\`a a scelta del Master
\end{abilist}

\classesoc{56-68}{Benestanti}{0.6}

Praticano molteplici mestieri. Possiedono di solito grandi esercizi
commerciali, come alberghi, gioiellerie, medi allevamenti e via
dicendo, che gli forniscono gran parte del reddito, variabile tra 15 e
300 corone al mese.

\begin{abilist}
\item ATL o ADT
\item Calcolare
\item Cavalcare bipedi o quadrupedi
\item Danzare
\item Galateo
\item Lingua levante
\item Lingua madre
\item Lingua straniera
\item Mitologia Religione
\item Storia Generale
\item Trattare
\item Valutare
\end{abilist}

\classesoc{69-78}{Borghesi ricchi}{0.7}

Sono proprietari di molteplici attivit\`a in cui investono il proprio
capitale. Sono perlopi\`u banchieri, armatori, grossi allevatori,
possessori di miniere, mercanti di schiavi, ecc.  Ricavano al mese da
100 a 2000 corone.

\begin{abilist}
\item ATL o ADT
\item Calcolare
\item Cavalcare quadrupedi
\item Danzare
\item Lingua levante
\item Lingua madre
\item Lingua straniera
\item Religione
\item Storia Generale
\item Trattare
\item Valutare
\end{abilist}

\classesoc{79-84}{Baroni}{0.7}

Oltre alle entrate dovute al gettito fiscale possono incamerare da 200
a 1200 corone al mese come rendita.

\begin{abilist}
\item ATL
\item ADT
\item Calcolare
\item Cavalcare bipedi o Cavalcare quadrupedi
\item Danzare
\item Galateo
\item Lingua levante
\item Lingua madre
\item Lingua straniera
\item Mitologia
\item Oratoria
\item Religione
\item Storia Generale
\item Valutare
\end{abilist}

\classesoc{85-91}{Conti}{0.7}

Oltre alle entrate dovute al gettito fiscale possono incamerare da 500 a
2000 corone al mese come rendita.

\begin{abilist}
\item ATL
\item ADT
\item Calcolare
\item Cavalcare bipedi o Cavalcare quadrupedi
\item Danzare
\item Galateo
\item Lingua levante
\item Lingua madre
\item Lingua straniera
\item Mitologia
\item Oratoria
\item Religione
\item Storia Generale
\item Valutare
\end{abilist}

\classesoc{92-96}{Marchesi}{0.8}

Oltre alle entrate dovute al gettito fiscale possono incamerare
da 1000 a 7000 corone al mese come rendita delle proprie terre.

\begin{abilist}
\item ATL
\item ADT
\item Calcolare
\item Cavalcare bipedi o quadrupedi
\item Danzare
\item Galateo
\item Lingua levante  
\item Lingua madre
\item Lingua straniera
\item Mitologia
\item Oratoria
\item Religione
\item Storia Generale
\item Valutare
\end{abilist}

\classesoc{97-00}{Duchi}{0.8}

Oltre alle entrate dovute al gettito fiscale possono incamerare da 2000
a 15000 corone al mese come rendita delle propriet\`a.

\begin{abilist}
\item ATL
\item ADT
\item Calcolare
\item Cavalcare bipedi o Cavalcare quadrupedi
\item Danzare
\item Galateo
\item Lingua levante
\item Lingua madre
\item Lingua straniera
\item Mitologia
\item Oratoria
\item Religione
\item Storia Generale
\item Valutare
\end{abilist}


\filler{pugnale_rileano.eps}


\subsection{Le Terre del Nord}

\Stendardo{stendardo_nordici.eps}

I Nordici sono umani e vivono nelle Terre del Nord.

\Geografia Le Terre del Nord occupano la parte Est dell'Isola di Norda
e confinano ad Ovest con Terranova e a Sud-Ovest con Mahn-Hesa.

Le coste delle Terre del Nord sono frastagliate e ricche di fiordi,
soprattutto a Nord.

Il clima \`e artico nell'estremo nord, mentre a Sud la temperatura \`e
pi\`u mite, grazie all'influenza benefica del Mare Interno.
Nell'entroterra assume caratteristiche continentali, con estati
tiepide e secche e inverni gelidi e nevosi. 

Il territorio \`e prevalentemente montuoso se si esclude la piana
Deys-Mathas. La cima pi\`u alta, facente parte del Massiccio centrale
del Dehrberg, al confine con Terranova, \`e il monte Kahlfahr, che
raggiunge i 3520 metri

I fiumi sono brevi e nascono solitamente da sorgenti di tipo glaciale
e sfociano ad estuario. Il pi\`u importante \`e il Kowdian, che \`e
navigabile nel suo tratto inferiore.

La vegetazione naturale \`e composta in maggioranza da foreste di
conifere, in cui predomina l'abete Novese; nelle zone pi\`u secche e
pianeggianti \`e diffusa la tundra.

Le citt\`a pi\`u importanti sono Norda, che \`e anche la capitale,
Finnan e Garhish dove sorge il porto pi\`u importante.

\Politica La nazione \`e frammentata in numerosi clan che sono spesso
in guerra tra loro. Ogni clan commercia e produce in maniera
indipendente. I capi clan si incontrano solo per le feste religiose
pi\`u importanti o nel caso di guerra contro altre nazioni sotto la
guida del Grande Sacerdote.

Il Clan \`e la base della societ\`a nordica. Il nucleo del Clan
\`e la Coorte di guerrieri costituita dalla famiglia del capoclan e
dalle famiglie dei suoi parenti. Gli appartenenti a tutte le altre
classi vivono sotto la loro protezione e vengono pertanto chiamati
Protetti.

I figli seguono sempre la carriera del padre anche se \`e possibile
chiedere al Capoclan una variazione o elevazione di classe sociale in
funzione dei propri meriti.

I poteri giuridico e religioso sono detenuti dal Druido. Il capoclan
detiene il potere militare e gestisce le risorse economiche.

\Esercito Ogni Clan ha la sua classe di guerrieri navigatori che
difendono il Clan, combattono contro i clan avversari, saccheggiano le
coste e prestano il loro servizio come mercenari presso le altre
nazioni. I Nordici sono i pi\`u temibili guerrieri sul mare interno,
 solo i Levanti sono loro degni antagonisti.

Sulla terraferma \`e temuta la loro falange mentre sono scadenti come
cavalieri e arcieri.

L'educazione militare viene impartita all'interno delle famiglie fin
dalla pi\`u tenera et\`a. Durante l'addestramento viene somministrata,
al fine di evitarne gli effetti collaterali in seguito, una droga chiamata
Rajina. 

Al compimento del quindicesimo anno di et\`a, ovvero la maturit\`a, i
giovani guerrieri ricevono l'ampolla di Rajina (contenente tre dosi)
che portano sempre al collo. In battaglia, ad un comando del capoclan,
tutti i guerrieri ne assumono una dose determinando in s\`e quella
terribile sete di sangue che ha fatto meritare ai Nordici
l'appellativo di Lupi Rabbiosi del Nord ed il rispetto delle nazioni
confinanti.

Le armi pi\`u diffuse sono: lo Spadone a due mani, la Spada Bastarda,
l'Ascia da battaglia, l'Ascia bastarda, il Martello da Guerra.

Nella Capitale si trova la principale Accademia di Arti Marziali.

\Economia All'interno di ogni Clan, la gestione delle risorse (che
comprendono anche le attivit\`a dei Protetti) \`e realizzata dai
capoclan.  Ogni famiglia di guerrieri protegge specificamente
determinate famiglie di Protetti e ha diritto al 10\% dei loro
proventi.  Il capoclan riceve il 5\% degli introiti di tutte le
famiglie del suo clan. 

Il druido vive delle ricche offerte che gli vengono elargite, si interessa
personalmente dell'acquisto della Rajina che viene venduta
ai guerrieri.

Le pecore, le capre e gli yak vengono allevati per ricavare pelli, lana,
latte e corna. 

Gli animali selvatici che sono oggetto di caccia riforniscono i
Nordici di pelli e pellicce pregiate.
Le grandi foreste lasciano ben poco spazio alle colture, ma il legno,
ricavato dagli abeti novesi e dai larici nordici, \`e di ottima
qualit\`a.

Il settore pi\`u fiorente dell'economia \`e la pesca, del pesce
azzurro e della balena, al largo delle coste del Mare Interno.

Le esportazioni riguardano pesce, lana, pellicce e formaggi, mentre si
importano cereali, ortaggi e prodotti finiti come armi e armature.

\Religione Ogni Clan ha
il suo druido, rappresentante degli spiriti del Clan, delle bestie che
lo proteggono e degli eroi che ne hanno fatto la storia. Il loro capo
\`e il sommo Druido, rappresentante dell'Uno che raccoglie in s\`e
tutti gli spiriti, della realt\`a che si cela dietro mille
apparenze. 

Ogni Druido sceglie il proprio successore, solitamente il figlio o un
bambino adottato ed educato a tale scopo. Il sommo Druido decide il
proprio successore in punto di morte, solitamente un altro Druido. I
giorni sacri corrispondono agli equinozi ed ai solstizi. In questi
giorni si svolgono le cerimonie di passaggio alla maggiore et\`a.

I matrimoni si svolgono ogni luna piena, i funerali ogni luna nuova.
In attesa di tale giorno i cadaveri vengono sottoposti ad un
trattamento speciale da parte dei Druidi.

\Magia Esistono solo 2 accademie
di magia nelle Terre del Nord, una si trova a Finnan e l'altra a
Hersh. A tutti coloro che sono in grado di pagare le rette di tali
accademie (2 piastre), vengono insegnate le basi della magia. Tuttavia,
grazie a borse di studio, coloro che vengono ritenuti pi\`u
meritevoli hanno accesso gratuito all'apprendimento. 

Le uniche scuole
insegnate sono Necromanzia e Demonologia. Presso la citt\`a di
Finnan si trova l'accademia retta dall'Arcimago.

\Moda I Nordici vivono in grosse abitazioni di pietra con tetti di legno e ricoperti
da zolle di terra.  Vestono con pesanti tessuti di lana. Le donne
indossano lunghe tuniche, gli uomini gonne lunghe fino al ginocchio,
stivali e camicie o maglioni molto pesanti.  Tutti utilizzano grossi
mantelli di pelliccia. 

Il disegno dei tessuti indossati e il loro colore \`e tipico del
Clan di provenienza ed \`e simbolo della raggiunta maturit\`a.

Verso i diciassette anni ci si sposa e le famiglie sono solitamente
molto numerose.

I divertimenti pi\`u comuni sono le gare di canoa, la lotta contro
lo yak (a mani nude o con un'arma a scelta) o il Breheg-by: quattro
squadre di 5 individui all'interno di un campo quadrato di 50x50
metri cercano di portare la carogna di una pecora oltre il lato del
campo appartenente alla propria squadra. Non \`e valido naturalmente
il ricorso ad armi e magia.

La lingua ufficiale \`e il nordico. 

\Fisico Alti in media 1.80 m, pesano solitamente attorno ai 90 kg; dato il clima
rigido, non \`e rara l'obesit\`a. La carnagione, i capelli e gli
occhi sono chiari. Portare i baffi e i capelli lunghi \`e segno di
grande virilit\`a. Il piercing e il tatuaggio sono permessi solo ai
guerrieri come segno di forza e di coraggio.  

\minmaxumani

\subsubsection{Abilit\`a di razza}

\begin{abilist}
\item Resistenza al freddo
\end{abilist}

\subsubsection{Classi sociali e abilit\`a di Classe sociale}

\classesoc{01-05}{Schiavi}{0.1}

Non vantano alcuna libert\`a, devono lavorare per tutta la vita alle
dipendenze del padrone e i loro figli nasceranno schiavi. I padroni
pi\`u benevoli possono concedere loro piccole somme mensili
con cui poter, un giorno, affrancarsi.

\begin{abilist}
\item Cantare
\item Conoscere Piante
\item Lingua Madre
\item Lingua Levante
\item Resistenza alla fatica
\item Resistenza al freddo
\item Resistenza alle malattie
\item Sollevare pesi
\end{abilist}

\classesoc{06-20}{Contadini}{0.3}

Producono cibo e vestiti per la propria famiglia, coltivano campi ed
allevano bestiame. Possono guadagnare al netto delle spese di
sostentamento da 1 a 10 scudi al mese.

\begin{abilist}
\item Ammaestrare
\item Carpenteria
\item Cavalcare quadrupedi
\item Conoscere animali
\item Conoscere piante
\item Lingua levante
\item Lingua madre
\item Resistenza alla fatica
\end{abilist}

\classesoc{21-35}{Cacciatori}{0.3}

Vivono della caccia e dei suoi proventi. Producono cibo e vestiti per
la propria famiglia e possono guadagnare da 2 a 7 corone al mese.

\begin{abilist}
\item Allarme
\item ADT o ADL
\item ATC o ABC
\item Conoscere animali
\item Conoscere piante
\item Fare nodi
\item Fiutare
\item Imitare
\item Individuare tracce
\item Lingua levante
\item Lingua madre
\item Muoversi in silenzio
\item Trappole
\item Udire
\end{abilist}

\classesoc{36-50}{Pescatori}{0.3}

Vivono della pesca e dello sfruttamento del mare e dei suoi proventi.
Producono cibo per la propria famiglia e possono guadagnare al netto
delle spese di sostentamento da 1 a 20 scudi al mese.

\begin{abilist}
\item Carpenteria
\item Conoscere animali
\item Controllo respirazione
\item Fare nodi
\item Lingua levante
\item Lingua madre
\item Navigare
\item Nuotare
\item Orientamento
\end{abilist}

\classesoc{51-70}{Professionisti}{0.5}

Praticano i pi\`u svariati mestieri e vivono di essi. Possono
essere: artigiani, cartografi, erboristi, carpentieri, macellai ecc.
Il loro introito mensile pu\`o variare da 5 a 20 corone al mese.

\begin{abilist}
\item Abilit\`a professionale \#1
\item Abilit\`a professionale \#2
\item Abilit\`a professionale \#3
\item Lingua levante
\item Lingua madre
\item Altre abilit\`a a scelta del Master
\end{abilist}
 
\classesoc{71-86}{Guerrieri navigatori}{0.4}

Ricevono dal clan un contributo mensile variabile tra 2 e 8 corone
oltre al vitto. Hanno poi diritto a prendere parte alla spartizione
del bottino di guerra.

\begin{abilist}
\item ADT
\item AIA
\item ATC o ABC
\item ABL
\item CAC
\item Maestria in un'arma
\item Allarme
\item Fare nodi
\item Lingua levante
\item Lingua madre
\item Navigare
\item Orientamento
\item Prontezza
\item Tattica militare
\end{abilist}

\classesoc{87-90}{Druidi}{0.7}

Sono custodi della memoria e dei riti religiosi dei clan. Vivono di
ci\`o che viene loro offerto dal popolo e della raccolta di erbe e
dei loro prodotti.  Il loro introito pu\`o variare da 1 a 50 scudi
al mese.

\begin{abilist}
\item Alchimia
\item Astronomia
\item Astrologia
\item Cavalcare bipedi
\item Conoscere piante
\item Lingua Madre
\item Lingua Levante
\item Medicina
\item Mitologia
\item Religione
\end{abilist}

\classesoc{91-100}{Capiclan}{0.5}

Provvedono alle esigenze del clan mantenendovi l'ordine e
amministrandovi la giustizia. I loro introiti possono variare da 5 a
30 corone al mese.

\begin{abilist}
\item ATC
\item Calcolare
\item Cavalcare quadrupedi
\item Giurisprudenza
\item Lingua levante 
\item Lingua madre
\item Mitologia 
\item Navigare
\item Oratoria
\item Raggirare
\item Recitare
\item Religione
\end{abilist}

\subsection{Le Terre Orientali}

\Stendardo{stendardo_orientali.eps}

Gli Orientali sono gli umani che abitano le terre orientali.

\Geografia Le terre orientali sono costituite da quattro isole che si
trovano a Sud di Malagana e a Nord Est di Loydi-Genya. Le quattro
isole sono chiamate Sangokushi (I tre Regni, per via dell'esistenza
delle Tre capitali itineranti), Kowa, Baori e Shibuta.

L'isola di Sangokushi (quella principale) presenta un terreno
prevalentemente montuoso con coste alte e scoscese. A Nord i rilievi
sono coperti dalle foreste di conifere; la cima pi\`u elevata \`e
costituita dal Monte Tama, con i suoi 5730 metri, facente parte della
catena del Kujitama. La sua cima \`e perennemente innevata e dal suo
ghiacciaio prendono vita i numerosi corsi d'acqua che si gettano al
mare nelle coste Est ed Ovest. A Sud il territorio degrada dolcemente
creando un'area collinosa anch'essa coperta da foreste chiamata
Kala-roweena, che significa ``Cuscino dei Demoni''.

Il fiume pi\`u importante, il Kijuku, \`e navigabile solo nei pressi
della foce. Le citt\`a stabili pi\`u importante sono costituite dalle
Capitali dei Potentati e sono Shiburu, Yerutsu, Koyada, Sashida,
Akutsu e Sikada. A Koyada si trova il porto pi\`u importante, sede dei
traffici con le altre isole dell'Arcipelago. Il clima \`e temperato.

\Politica La massima autorit\`a politica, economica e religiosa \`e
l'imperatore, che trasmette la sua carica al proprio primogenito
maschio. Egli ha il potere di assegnare e revocare terre e titoli
nobiliari (che sono ereditari), di riscuotere le tasse, arruolare i
soldati, legiferare e giudicare i criminali (solitamente si interessa
per\`o del giudizio dei nobili).

I diretti subordinati dell'imperatore sono i Reggenti che hanno il
controllo sulle unit\`a territoriali in cui \`e suddivisa la nazione,
detti Potentati.  Essi attualmente sono 6, ma possono variare di
numero per annessioni o attribuzioni imperiali. Il controllo sui
Potentati \`e realizzato tramite messi, poich\'e per legge chi \`e
insignito del titolo di Reggente \`e tenuto a far parte della Corte
Errante dell'Imperatore (in tale maniera egli \`e in grado di
controllare l'operato dei suoi pi\`u diretti subordinati). Il
controllo diretto dei Reggenti sui Potentati si verifica soltanto
quando la corte viaggiante li attraversa. Il Reggente ha il potere di
suggerire all'imperatore l'assegnazione di titoli nobiliari, di
riscuotere i tributi.

Nel territorio dei Potentati esistono delle citt\`a stato, chiamate
Signorie, governate dai Signori e dalle loro famiglie e gestite con
leggi da loro emanate (e che sono state accettate dai Reggenti). Qui i
Signori amministrano la giustizia, organizzano l'esercito e gestiscono
la raccolta delle tasse di cui trattengono il 5\%. Le Signorie
comprendono solitamente tutte le terre circostanti. Spesso si
verificano conflitti fra le Signorie che possono portare
all'eventuale annessione del territorio della Signoria sconfitta a
quello della Signoria vincitrice.

La Corte Errante completa il giro di tutte le Signorie dell'Impero in
circa 10 anni. Quando passa in una signoria, l'imperatore giudica
l'operato della famiglia del Signore che le ha governate. Se
l'amministrazione non \`e ritenuta idonea la Signoria viene attribuita
ad un Signore di nuova nomina.  L'imperatore pu\`o anche decidere di
riattribuire le Signorie conquistate al precedente Signore. Alcuni
Signori svolgono il ruolo di ambasciatori presso le altre nazioni.

\Esercito L'esercito Imperiale, costituito dagli Eserciti di ogni
Signoria, \`e comandato dall'imperatore tramite i Reggenti e i
Signori. 

I giovani Signori, educati fino ai 12 anni in seno alle famiglie, dopo
tale et\`a, che ne segna anche il passaggio alla maturit\`a, si recano
in uno dei templi dove vivranno allenandosi nelle arti belliche tutti
assieme, in modo da creare l'affiatamento necessario nei periodi di
guerra. Tale addestramento perdura almeno fino ai 19 anni, ma il
giovane Signore pu\`o decidere di restare al tempio. I templi pi\`u
conosciuti sono quelli di Otomo e Tanaka. 

Ogni Signoria ha un suo Esercito, comandato dal Signore, i cui soldati
vengo scelti tramite pubblici concorsi visto che la carriera militare
\`e molto ambita dalle popolazioni meno abbienti.

Se si verificano guerre fra Signorie, peraltro incoraggiate
dall'Imperatore per tenere allenati i Militari, la Signoria
vincitrice, oltre ad annettere la citt\`a avversaria, ha diritto a
veder risarciti da parte del Signore sconfitto i danni di guerra. La
famiglia viene assoggettata al controllo della Signoria vincente e non
le resta che attendere un'eventuale riattribuzione imperiale. La
tecnica di combattimento pi\`u usata \`e la guerriglia.

Le armi utilizzate comunemente dall'esercito sono i katana, le
alabarde e gli archi lunghi. 

Esistono inoltre 2 Accademie di Arti Marziali. Una (il tempio di
Otomo), presso il Monte Tama, l'altra (l'accademia di Sakamoto) nella
citt\`a di Sokata. A causa della rivalit\`a fra le due Accademie,
coloro che hanno frequentato una di esse non sono ammessi agli
insegnamenti dell'altra.

\Economia L'imperatore, i Reggenti e i Signori vivono delle imposte
sul reddito di ogni individuo (la pressione fiscale si aggira intorno
al 35\%). Essi, tuttavia, possono essere proprietari di attivit\`a
commerciali. Le coltivazioni principali sono costituite da cereali,
alberi da frutto e alberi da legna. La pesca e l'allevamento di ovini,
bovini, suini ed equini soddisfa il fabbisogno interno.

Dalle miniere si estraggono pietre preziose e marmo.

L'artigianato \`e tra i migliori soprattutto per ci\`o che concerne
pietre e metalli preziosi.

\Religione La religione \`e di tipo animista, ogni cosa ha il suo spirito. Le divinit\`a
maggiori sono il Cielo e la Terra, sotto la cui protezione vive l'Imperatore
(Figlio del Cielo).

I riti hanno luogo attorno a vasti cerchi di pietre che rappresentano
l'universo, al centro di esso vi \`e un tamburo (il cuore pulsante
della Terra), il cui suono richiama i fedeli alla funzione. Durante il
rito si elevano preghiere agli spiriti della Citt\`a, della Famiglia e
dell'individuo.  Nella tradizione, questi ultimi proteggono la persona
cui sono legati ed il loro emblema e colore sono portati in guerra dal
loro protetto.

A tutti gli spiriti sono offerti cibo e oggetti di valore che vengono
bruciati, sepolti ed in parte fatti propri dagli officianti. 

Gli officianti sono chiamati Monaci e nella scala sociale sono pari ai
Signori. Alcuni di essi sono addirittura figli di Signori che hanno
rinunciato ai diritti e ai doveri che li legavano alla famiglia per
abbracciare la ricerca spirituale.

Gli Orientali sono molto legati alla ricerca interiore, praticano la
meditazione e si cimentano nelle discussioni filosofiche.

Nel Pantheon delle divinit\`a, grande importanza hanno anche i Demoni
(Oni), che possono portare fortuna o sciagura.

\Magia La magia viene insegnata in accademie affiliate al Conclave, che
richiedono rette elevate per cui solo i ceti pi\`u abbienti possono
permettersi la loro frequentazione.  Tuttavia i meno abbienti possono
fare ricorso all'Anticonclave, che qui \`e abbastanza influente.

Le scuole di magia insegnate nelle accademie sono Demonologia,
Esistenza ed Illusionismo.

\Moda L'abbigliamento pi\`u comune \`e costituito, per le donne, da un
panno avvolto su tutto il corpo, mentre per gli uomini da una semplice
tunica. Nei periodi pi\`u caldi, gli uomini indossano una sorta di
gonna. I Signori utilizzano una gonna-pantalone molto larga, i
Reggenti e l'Imperatore indossano kimono di tessuti preziosi. Nei
periodi pi\`u freddi tutti utilizzano mantelli di pelli.

L'imperatore utilizza, come simbolo del suo potere, un mantello di
piume d'uccello, i Signori utilizzano le piume di uccello come diadema
sull'elmo.  Nei centri urbani l'armatura pu\`o essere indossata solo
dal Signore e dai suoi militari.

I capelli sono portati solitamente lunghi e rasati sulle tempie.  I
signori usano spesso il taglio a criniera di cavallo. Il divertimento
pi\`u comune \`e rappresentato dalle competizioni marziali e dai
giochi di intelligenza come Scacchi e Dama.

I cibi sono elaborati e solitamente basati sui vegetali.

Particolare \`e il modo di salutare che \`e tipico per ogni
titolo. L'imperatore non saluta mai, i Reggenti alzando una mano, I
Signori avvolgendo il pugno destro con la mano sinistra. Il popolo
saluta con un inchino. Davanti all'imperatore tutti sono tenuti a
prostrarsi. 

La lingua ufficiale \`e l'Orientale.

\Fisico Il popolo orientale \`e composto da due genti con diversi
caratteri somatici. I primi sono caratterizzati da una bassa statura
(circa 170 cm), occhi scuri a mandorla, carnagione giallastra e
capelli lisci e scuri. La seconda, di statura pi\`u elevata (185 cm di
media), ha carnagione scura, capelli scuri, occhi scuri e tratti
somatici levanti. La gran parte della popolazione \`e costituita da
meticci.

\minmaxumani

\subsubsection{Classi sociali e abilit\`a di Classe sociale}

\classesoc{01-05}{Schiavi}{0.1}

Non vantano alcuna libert\`a, e devono lavorare per tutta la vita alle
dipendenze del padrone, i loro figli nasceranno schiavi. I padroni
pi\`u benevoli possono concedere agli schiavi piccole somme mensili
con cui gli schiavi potranno, un giorno, affrancarsi.

\begin{abilist}
\item Cantare
\item Conoscere
\item Piante 
\item Lingua Madre
\item Lingua Levante
\item Resistenza alla fatica
\item Resistenza al freddo
\item Resistenza alle malattie
\item Sollevare pesi
\end{abilist}

\classesoc{06-15}{Servi}{0.2} 

Vivono nella casa del proprio signore, ricevendo vitto e alloggio e un
compenso mensile variabile da 1 a 3 Corone.

\begin{abilist}
\item Calcolare
\item Cantare o Danzare o Suonare o Cucinare o Abilit\`a professionale
  collegata alla mansione
\item Cavalcare quadrupedi
\item Galateo
\item Lingua Madre
\item Lingua Levante
\item Muoversi in silenzio
\end{abilist}

\classesoc{16-30}{Pescatori}{0.3}

Vivono della pesca e dello sfruttamento del mare e dei suoi proventi.
Producono cibo per la propria famiglia e possono guadagnare al netto
delle spese di sostentamento da 1 a 20 scudi al mese.

\begin{abilist}
\item Carpenteria
\item Conoscere animali
\item Controllo respirazione
\item Fare nodi
\item Lingua levante
\item Lingua madre
\item Navigare
\item Nuotare
\item Orientamento
\end{abilist}

\classesoc{31-40}{Contadini}{0.3} 

Coltivano campi ed allevano bestiame, producendo cibo e vestiti
per la propria famiglia. Possono guadagnare al netto delle spese di
sostentamento da 1 a 10 scudi al mese.

\begin{abilist}
\item Ammaestrare 
\item Carpenteria
\item Cavalcare quadrupedi
\item Conoscere animali
\item Conoscere piante
\item Lingua levante
\item Lingua madre
\item Resistenza alla fatica
\end{abilist}

\classesoc{41-60}{Professionisti}{0.5} 

Praticano i pi\`u svariati mestieri e vivono di essi. Possono essere:
artigiani, giuristi, cartografi, erboristi, carpentieri, medici,
macellai, guerrieri ecc. Il loro introito mensile pu\`o variare da 5 a
30 corone al mese.

\begin{abilist}
\item Abilit\`a professionale \#1 
\item Abilit\`a professionale \#2
\item Abilit\`a professionale \#3
\item Lingua levante
\item Lingua madre 
\item Altre abilit\`a a scelta del Master
\end{abilist}

\classesoc{61-75}{Benestanti}{0.5}

Praticano molteplici mestieri. Possono essere proprietari di grandi
esercizi commerciali, come alberghi, gioiellerie, fattorie, e via
dicendo, che forniscono gran parte del loro reddito, variabile tra 15 e
200 corone al mese.

\begin{abilist}
\item ATL o ADT
\item Calcolare
\item Cavalcare bipedi o quadrupedi
\item Danzare
\item Galateo
\item Lingua levante
\item Lingua madre
\item Lingua straniera
\item Mitologia
\item Religione
\item Storia Generale
\item Trattare
\item Valutare
\end{abilist}

\classesoc{76-88}{Borghesi ricchi}{0.6}

Sono proprietari di molteplici attivit\`a in cui investono il proprio
capitale. Sono perlopi\`u banchieri, armatori, grossi allevatori,
possessori di miniere, mercanti di schiavi, ecc.  Ricavano al mese da
50 a 1500 corone.

\begin{abilist}
\item ATL o ADT
\item Calcolare
\item Cavalcare bipedi o quadrupedi
\item Danzare Galateo
\item Lingua levante
\item Lingua madre
\item Lingua straniera
\item Mitologia
\item Religione
\item Storia Generale
\item Trattare
\item Valutare
\end{abilist}




\classesoc{89-94}{Signori}{0.7} 

Sono a capo di una signoria e oltre alle entrate dovute al gettito
fiscale possono incamerare da 200 a 2000 corone al mese come rendita
delle loro propriet\`a.

\begin{abilist}
\item ATL
\item ADT 
\item AIA 
\item CAC 
\item Astronomia 
\item Calcolare 
\item Cavalcare quadrupedi 
\item Galateo
\item Letteratura
\item Lingua levante
\item Lingua madre
\item Medicina 
\item Musica
\item Nuotare
\item Religione
\item Storia Generale
\end{abilist}

\classesoc{95-99}{Reggenti}{0.7}

Seguono la corte imperiale e dipendono da essa; i loro introiti
derivano dalla riscossione delle imposte del Potentato. Possono
incamerare mensilmente fino a 300 piastre.

\begin{abilist}
\item ATL
\item ADT
\item Calcolare
\item Cavalcare quadrupedi
\item Danzare
\item Galateo
\item Lingua levante
\item Lingua madre
\item Lingua straniera
\item Mitologia
\item Oratoria
\item Religione
\item Storia Generale
\item Valutare
\end{abilist}


\classesoc{100}{Imperatore}{0.8}

Vive degli introiti derivanti dalle tasse pagate da tutti i suoi
sudditi. Pu\`o avere introiti mensili fino a 2000 piastre

\begin{abilist}
\item ATL
\item Calcolare
\item Cavalcare quadrupedi
\item Danzare
\item Galateo
\item Letteratura
\item Lingua levante
\item Lingua madre
\item Lingua straniera
\item Mitologia
\item Musica
\item Oratoria
\item Religione
\item Storia Generale
\end{abilist}

\filler{libro1.eps}

\subsection{I Nomadi}

\pinupbig{romdellesabbiebig.eps}{Rom delle sabbie}{tb}

I nomadi sono popolazioni non stanziali, appartenenti ad ogni razza ed
etnia.  Essi hanno abbandonato da generazioni la loro societ\`a
originaria per abbracciare la vita da girovaghi. \`E possibile che
ogni gruppo di nomadi sia formato da genti di razze ed etnie diverse.
Solitamente a legarli alle loro origini \`e il solo aspetto religioso.

Per evitare di essere sopraffatti da altri popoli, sono soliti riunire
i clan in comunit\`a pi\`u grandi.  I nomadi sono diffusi in tutte le
terre conosciute, ma le comunit\`a pi\`u importanti sono:

\begin{description}
\item{\textbf{Rom delle sabbie},} che si trovano nel territorio desertico di
  Kush e si spostano di oasi in oasi per nutrire buoi e vacche. La
  loro cavalcatura tipica \`e il dromedario.  La loro arma \`e la
  scimitarra (spada) e si vestono con abiti chiari leggeri e ampi. Si
  spostano con carri dotati di slitte, che hanno una struttura a
  cupola realizzata con pelli conciate.
  
\item{\textbf{Abitanti delle pianure},} dislocati nelle Grandi Pianure
  dell'isola di Bahuney, vivono di caccia e raccolta di frutti e
  radici. Sono fieri della loro libert\`a, sono abili cavalieri e
  guerrieri; usano Archi Lunghi, Alabarde, Lance, Spade e Asce.
  Si spostano con cavalli ed i loro accampamenti sono costituiti da
  tende leggere e facilmente trasportabili. Vestono abiti di pelli.

\end{description}

Un nomade che appartenga ad una trib\`u che ha uno sciamano pu\`o
avere nozioni magiche limitate ad una sola scuola di magia con un VAL
di max. 5. Il Conclave tende a reprimere queste usanze, perci\`o
l'insegnamento e l'uso della magia sono segreti.

\Fisico I nomadi possono appartenere a qualsiasi razza, frequentemente
sono meticci. In maggioranza si tratta di umani, anche se \`e facile
trovare tra loro gnomi (in particolare di origine Kunetha), orchi ed
elfi di origine Umbra.

\subsubsection{Abilit\`a di razza}

\begin{abilist}
\item Resistenza al caldo 7
\end{abilist}

\subsubsection{Classi sociali e abilit\`a di Classe sociale}

\classesoc{1-40}{Pastori}{0.3}

Portano al pascolo il gregge, che d\`a latte, carne e pelli a tutta la
trib\`u.  Il ricavato appartiene alle casse comuni.

\begin{abilist}
\item Ammaestrare 
\item Cavalcare bipedi 
\item Cavalcare quadrupedi 
\item Conoscere Animali 
\item Conoscere Piante 
\item Empatia 
\item Fare nodi 
\item Lingua levante 
\item Lingua madre 
\item Orientamento 
\item Trattare 
\end{abilist}

\classesoc{41-75}{Guerrieri/Cacciatori}{0.4}

Difendono la trib\`u e procurano carni e pelli per i suoi bisogni.

\begin{abilist}
\item Allarme ADT o ADL ATC o ABC 
\item Cavalcare quadrupedi 
\item Conoscere animali 
\item Conoscere piante 
\item Fare nodi 
\item Fiutare 
\item Imitare 
\item Individuare tracce 
\item Lingua levante 
\item Lingua madre 
\item Muoversi in silenzio 
\item Trappole 
\item Udire 
\end{abilist}

\classesoc{76-95}{Artigiani}{0.4} 

Svolgono lavori utili a tutta la trib\`u e vendono i loro prodotti.
Possono essere: falegnami, pellai, fabbri o stagnini, cartografi,
gioiellieri, intagliatori (osso, legno) o addirittura ladri.

\begin{abilist}
\item Abilit\`a professionale \#1 
\item Abilit\`a professionale \#2 
\item Abilit\`a professionale \#3 
\item Lingua levante 
\item Lingua madre 
\item Altre abilit\`a a scelta del Master 
\end{abilist}

\classesoc{96-98}{Sciamani}{0.7} 

Sono custodi della memoria e dei riti religiosi dei clan. Vivono di
ci\`o che viene loro offerto dal popolo, della raccolta di erbe e dei
loro sottoprodotti.

\begin{abilist}
\item Alchimia 
\item Astronomia 
\item Astrologia 
\item Cavalcare bipedi 
\item Conoscere piante 
\item Lingua Madre 
\item Lingua Levante 
\item Medicina 
\item Mitologia 
\item Religione 
\end{abilist}

\classesoc{99-100}{Capitrib\`u}{0.6} 

Provvedono alle esigenze del clan mantenendovi l'ordine e amministrandovi la
giustizia. Gestiscono la cassa comune, che pu\`o variare da 1 a 5 corone
al mese per membro della trib\`u.

\begin{abilist}
\item AIA 
\item ADT o ADL 
\item Calcolare 
\item Cavalcare quadrupedi 
\item Giurisprudenza 
\item Lingua levante 
\item Lingua madre 
\item Lingua straniera \#1 
\item Lingua straniera \#2 
\item Mitologia 
\item Oratoria 
\item Raggirare 
\item Recitare 
\item Religione
\end{abilist}



\subsection{La Confederazione del Sud}

\Stendardo{stendardo_chire.eps}

La Confederazione del Sud \`e abitata dai Chire, umani di pelle scura.

\Geografia La Confederazione del Sud si trova nella parte Sud Ovest
dell'Arcipelago a circa 750 km dalle coste di Bahuney. \`E formata da
9 chiatte (Khares, singolare Khare) rettangolari, ancorate al fondo
del mare (in corrispondenza di una montagna sottomarina, la cui cima
si trova 650 metri al di sotto della superficie), ognuna delle quali
occupa approssimativamente 100 km quadrati. Otto di loro sono
disposte in cerchio a formare una rosa dei venti, mentre la nona si
trova al centro di esso. Non esistono ponti che collegano fra loro le
Khares, tutti i trasferimenti vengono fatti mediante imbarcazioni.
Ogni Khare perci\`o \`e dotata di imponenti porti da dove navi
enormi (Bahn-yere) prendono il largo, o dove esse attraccano.

Il clima \`e di tipo tropicale con precipitazioni molto scarse ed
estati torride. Non esiste vegetazione spontanea in quanto le Khares
sono artificiali; gli abitanti delle Khares riescono comunque a
coltivare all'interno di serre specializzate tutto ci\`o che serve
alla loro sopravvivenza.

In casi di emergenza le Khares possono essere disancorate e possono
muoversi come enormi imbarcazioni.

\Politica La Confederazione del Sud \`e abitata dai Chire, ed \`e una
repubblica parlamentare. Le autorit\`a pi\`u importanti di ogni Khare
sono i Governatori. Ogni Governatore viene eletto in ogni Khare fra
tutti i Possessori (proprietari di Bahn-yere) e viene votato tramite
suffragio universale da tutti gli abitanti con diritto di cittadinanza
che abbiano superato il venticinquesimo anno di et\`a.

Ogni governatore resta in carica 15 anni e pu\`o essere rieletto per
una sola volta.  Il governatore ha diritto a scegliere i suoi ministri
per il Governo di Khare (Organo esecutivo), ma questi devono essere un
numero massimo di 9 e devono essere cittadini della Khare. Essi
restano in carica per 6 anni.

L'organo legislativo \`e rappresentato dal Consiglio di Khare formato
da 35 membri eletti fra tutti coloro che abbiano avuto nell'ultimo
anno un reddito non inferiore alle 50 piastre. Esso resta in carica
per 6 anni, controlla l'operato del governo e promulga le leggi. Ogni
Consiglio di Khare vota il suo Presidente a maggioranza relativa.

Il potere giudiziario spetta al Consiglio di Giustizia, che resta in
carica per 7 anni ed \`e formato da 9 membri eletti fra gli iscritti
all'albo dei giudici praticanti la professione da almeno 15 anni.
Il pi\`u anziano dei membri \`e il Presidente del consiglio di
Giustizia.

Esiste poi un organismo superiore che viene convocato per le decisioni
riguardanti la confederazione nella sua integrit\`a, per questioni di
interesse pubblico, per dichiarazioni di guerra, ecc. Esso \`e il
Consiglio Generale ed \`e formato da tutti i Governatori di ogni
Khare, e dai Presidenti del Consiglio di Khare e del Consiglio di
Giustizia. Il Governatore pi\`u anziano \`e detto Sowhmel-hoos e
svolge le fun\-zioni di pre\-si\-den\-te.

Ogni Khare ha un suo Codice di Leggi, ma esiste un Codice Confederale
che regola i reati pi\`u gravi. In tempo di guerra in tutta la
Confederazione ha validit\`a solo il Codice di Guerra. Ogni Khare
assume il nome che ne contraddistingue la sua posizione nella rosa dei
venti: Khare del Nord, Khare del Sud, dell'est, dell'Ovest, del
Sud-Est e cos\`{\i}  via. La Khare Centrale \`e la capitale ed \`e la
sede del Consiglio Generale.

\pinup{chire.eps}{Vedetta Chire}

\Esercito L'esercito della Confederazione \`e formato dai
Professionisti e dalla Leva.  La leva \`e obbligatoria per tutti i
cittadini maschi che abbiano compiuto il 18 anno di et\`a e vengano
ritenuti idonei. Ad essi viene insegnato l'utilizzo di Spada bastarda,
Sciabola e Arco lungo.

La leva dura 18 mesi. Tutti coloro che alla fine di tale periodo
vogliano diventare Professionisti devono sostenere una prova e giurare
fedelt\`a alla Confederazione e alla Khare di appartenenza.
L'esercito viene utilizzato prettamente per mantenere l'ordine
interno.

Ogni Khare ha il suo esercito guidato da un Comandante Generale,
nominato dal Ministro della Difesa, che resta in carica per 4 anni.
Esistono corpi speciali ben addestrati che vengono mandati sulle Bahn-yere
per difenderle da eventuali attacchi esterni e che vengono pagati
direttamente dai Possessori.

Il settore pi\`u importante dell'esercito \`e ovviamente la Flotta che
riveste funzioni di controllo all'interno della Confederazione e ne
difende i confini. Essa appartiene alla Confederazione nella sua
totalit\`a ed \`e capeggiata dall'ammiraglio Generale nominato dal
Consiglio Generale.

In caso di dichiarazione di guerra l'esercito viene unificato e viene
nominato dal Consiglio Generale un Comandante di Guerra che ne regola
gli interventi.  Egli rimane in carica fino alla fine della Guerra, o
fino alla sua morte. Il Consiglio Generale si riserva il diritto di
destituirlo in casi particolari.  Al momento dell'entrata in guerra
scatta in tutto il territorio della Confederazione il Codice di
Guerra.

\Economia L'economia \`e praticamente controllata
dai Possessori.  Essi sono i proprietari delle Bahn-yere (le uniche
navi capaci di navigare al di fuori del mare Interno), e commerciano
con tutta la parte Sud-occidentale dell'Arcipelago. I loro commerci si
spingono ad Est fino a Loydi-Genya e alle coste meridionali delle
Terre Orientali, mentre ad occidente raggiungono le coste di
Dyshu-Thandy. Solitamente le imbarcazioni Chire non navigano nel Mare
Interno per evitare le navi dei Levanti e degli Elfi Luxi che in quel
tratto di mare hanno il monopolio sul commercio e la pesca, ma si
spostano lungo le coste che si affacciano sul mare esterno, senza
allontanarsi troppo da esse.  La ragione della loro ricchezza sta
nella notevole differenza di prezzo d'acquisto delle merci e il prezzo
di vendita, giustificato anche dalla variet\`a e dei loro prodotti e
dalla loro affidabilit\`a.

I Possessori comprano merci su richiesta degli acquirenti direttamente
nel luogo di produzione e ne controllano la qualit\`a. Riescono a
procurare le pi\`u stravaganti invenzioni degli gnomi ai facoltosi
Borghesi di Dyshu-Thandy e le affilatissime lame degli Orientali ai
nobili di Bahuney. Le Bahn-yere oltre ad avere delle stive
particolarmente capienti sono in grado di trasportare un grande numero
di persone, e nel corso degli anni sono diventate grandi al punto di
diventare piccoli villaggi mobili. Tutta la famiglia del Possessore
vive sulla nave.

La propriet\`a della Bahn-yere \`e ereditaria e viene trasferita al
figlio maschio pi\`u grande in caso di morte del Possessore o qualora
egli venisse eletto Governatore. I figli cadetti possono decidere di
restare sulla nave oppure possono trasferirsi in una Khare. A titolo
di imposta i Possessori devono pagare alla Confederazione una
quantit\`a di denaro pari al 20\% del loro guadagno annuo.

Importante per l'economia della Confederazione \`e ovviamente la pesca
d'alto mare che viene praticata non lontano dai porti delle Khares. Il
settore terziario \`e abbastanza sviluppato: sono presenti le scuole
(che sono obbligatorie fino ai 10 anni di et\`a) e i servizi postali.
L'agricoltura si svolge all'interno di serre specializzate e consente
la coltivazione di ortaggi e alberi da frutto.

Molti di questi prodotti vengono esportati nell'Arcipelago come
prodotti tropicali.  Importante \`e anche l'allevamento del baco da
seta, la produzione di fibre tessili vegetali, e la coltivazione del
tabacco.

Famosa \`e la produzione dei Frutti Gialli, che i Chire coltivano e
portano sempre con loro in viaggio perch\'e nutrendosi di essi si
previene lo scorbuto.  

Ogni abitante della Confederazione che percepisce reddito \`e tenuto a
pagare un'imposta proporzionale alla sua capacit\`a contributiva.

\Religione Non esiste una Religione ufficiale, quella pi\`u praticata
\`e di tipo monoteistico. I Chire credono in un Dio creatore di tutte le cose,
che decide sulla vita e sulle vicende di ogni persona.

Non esistono Sacerdoti cosicch\'e le funzioni religiose vengono
presiedute da coloro che hanno indirizzato i loro studi verso la
Religione.  I Chire credono nell'avvento di un Messia che assolve
tutti i peccati e che purifica le anime. Per loro la salvezza di
ognuno sta nelle sue opere ed il Dio decide sul futuro degli uomini in
relazione a ci\`o che hanno compiuto nel corso della loro vita.

Sono tollerate tutte le altre religioni. \`E abbastanza consistente la
comunit\`a dei Rileani e si possono trovare soprattutto sulle Khares
dell'est, del Nord-Est e del Sud-Est dei templi del Mildril edificati
da sacerdoti o sacerdotesse Luxi.

\Magia Nella Confederazione del Sud sono ubicate tre grandi Accademie
facenti parte del Conclave. Esse si trovano nella Khare centrale,
nella Khare del Nord ed in quella del Sud. L'accademia della Khare
centrale \`e retta da un Arcimago.

L'utilizzo della magia \`e perfettamente legale ed il lavoro dei maghi
\`e ben pagato; i maghi possono essere mandati ad esercitare le loro
capacit\`a sulle Flotte militari, sulle Bahn-yere e come rinforzo per
mantenere l'ordine interno sulle Khares. \`E tuttavia vietato
l'insegnamento della magia al di fuori delle Accademie.

Le scuole insegnate sono Esistenza, Stregoneria ed Elementalismo. La
retta da pagare per poter frequentare un'accademia \`e di 10 piastre
annuali, cosicch\'e soltanto i pi\`u ricchi (di norma i figli dei
Possessori o dei Governatori) possono accedervi. Le Accademie
concedono comunque delle borse di studio a tutti coloro che dimostrino
grandi capacit\`a e che riescano a superare con il massimo dei voti il
test d'ingresso.

Il fatto che le chiatte galleggino sulla superficie dell'acqua
impedisce ai maghi di acquisire l'energia magica dei Venti della
Terra. I Chire risolvono questo problema in diversi modi: estraendo
l'energia magica da alcune specie di animali marini (stelle marine
della specie Kalareyna, tartarughe giganti Tsyu-betch) oppure
immagazzinando energia all'interno di oggetti magici durante i viaggi;
questo permette di riconoscere i maghi Chire, visto che portano detti
oggetti alla cintura.

Un terzo mezzo \`e costituito dai ``pozzi'', lunghissimi e sottili
pali metallici che vengono conficcati sul fondo del mare e arrivano
fino alle Khares in superficie, dai quali si pu\`o attingere energia a
pagamento. Esiste un pozzo per ogni Khare e viene gestito da
funzionari della Confederazione, la quale fa suoi tutti gli introiti.

\Moda I Chire vestono abiti molto sgargianti e molto leggeri. Sono
molto socievoli e molto ospitali con gli stranieri e questi vengono
sempre trattati nel migliore dei modi. 

In linea di massima ad i lavoratori sono concessi pochi svaghi a causa
del lungo periodo lavorativo. Ai dipendenti della Confederazione come
insegnanti, medici e militari vengono concesse delle ferie
facoltative, che non vengono remunerate.  Il periodo massimo di ferie
\`e di 25 giorni distribuite per tutto il corso dell'anno.

La scuola \`e obbligatoria per tutti i bambini fra i 6 ed i 10 anni ed
\`e completamente gratuita. Esistono poi delle scuole propedeutiche a
pagamento che introducono alla professione che viene svolta nelle
scuole Professionali (frequentate dai ragazzi fra i 14 ed i 18 anni).
Le universit\`a sono a disposizione soltanto di pochi in quanto
abbastanza costose; la durata media dei corsi \`e di 3 anni e nelle
facolt\`a si pu\`o scegliere fra due indirizzi: Religione e Lingue.

In linea di massima all'interno delle confederazione il tenore di vita \`e
abbastanza elevato cosicch\'e quasi tutti vivono abbastanza agiatamente, e pressoch\'e
nessuno ha problemi di sussistenza. 

Lo sport nazionale \`e la Regata a Vela e consiste in gare con
velocissime imbarcazioni. \`E stato istituito un Campionato Confederale
e vi sono anche dei navigatori professionisti, anche se molti di
coloro che partecipano alle gare sono militari della flotta
confederale. Particolarmente diffuse sono le gare di nuoto.

Passeggiando per le strade delle Khare \`e molto frequente vedere
musicisti, cantanti o ballerini che rallegrano l'ambiente. Molti di
loro vengono assoldati da compagnie di Canto Popolare e suonano nelle
taverne o fanno tourn\'ee per tutto l'Arcipelago, rendendo proverbiali
le capacit\`a musicali dei Chire.

La lingua ufficiale \`e il Chire. 

\Fisico I Chire sono umani di pelle nera abbastanza scura. La loro
altezza media \`e di 1 metro e 80 e peso medio 75 chili, hanno capelli
neri e crespi e occhi neri ed il loro fisico \`e particolarmente
atletico ed asciutto.

\minmaxumani

\subsubsection{Abilit\`a di Razza}

\begin{abilist}
\item Nuotare 10 
\item Resistenza al caldo 8 
\end{abilist}

\subsubsection{Classi sociali e abilit\`a di Classe sociale}

\classesoc{01-80}{Professionisti}{0.5}

Praticano i pi\`u svariati mestieri e vivono di essi. Possono essere:
artigiani, commercianti, cartografi, giuristi, medici, erboristi,
carpentieri, pescatori, macellai ecc. Il loro introito mensile pu\`o
variare da 5 a 100 corone al mese.

\begin{abilist}
\item ATL o ADT 
\item Abilit\`a professionale \#1 
\item Abilit\`a professionale \#2
\item Abilit\`a professionale \#3 
\item Lingua levante
\item Lingua madre 
\item Altre abilit\`a a scelta del Master 
\end{abilist}

\classesoc{81-95}{Possessori}{0.6}

Sono i proprietari di navi Bahn-yere e vivono di commercio. Il
loro introito mensile pu\`o variare da 10 a 60 piastre.

\begin{abilist}
\item ATL o ADT
\item Calcolare 
\item Danzare
\item Galateo
\item Lingua levante
\item Lingua madre
\item Lingua straniera 
\item Mitologia
\item Navigare
\item Religione
\item Storia 
\item Trattare
\item Valutare
\end{abilist}

\classesoc{96-00}{Governatori}{0.7}

Sono coloro che sono stati eletti governatori. Guadagnano da 20 a 70
piastre al mese.

\begin{abilist}
\item ATL o ADT
\item Calcolare
\item Cavalcare quadrupedi
\item Danzare
\item Galateo
\item Giurisprudenza
\item Lingua levante
\item Lingua madre
\item Lingua straniera
\item Mitologia
\item Navigare
\item Oratoria
\item Religione
\item Storia
\item Valutare
\end{abilist}


\section{Gli Elfi} 

\subsection{Lo Stato Elfico Di Luxi}

\Stendardo{stendardo_luxi.eps}

Lo Stato Elfico di Luxi \`e abitato dagli elfi Luxi.

\Geografia Lo Stato Elfico di Luxi occupa interamente l'isola di Bahuney, situata
a Sud-Ovest dell'Arcipelago. Le coste settentrionali sono abbastanza
basse e non eccessivamente frastagliate, il che favorisce
l'insediamento e la costruzione di porti.

Il territorio di Bahuney \`e prevalentemente pianeggiante e coperto da
foreste di latifoglie, che diventano in certi tratti impercorribili per
via della densit\`a della vegetazione, e macchia.

 A Sud, la met\`a occidentale della costa che si affaccia
al Mare Esterno \`e alta con falesie a picco per via della presenza dei Monti
Simeinel, catena lunga circa 600 km, la cui cima pi\`u alta \`e il Monte
Shauti (2026 metri).

Altre piccole catene e massicci montuosi punteggiano il territorio, ma
nessuna di esse supera i 2000 metri. La costa Nord-Orientale \`e
ricca di bellissime spiagge, molte delle quali inaccessibili da terra,
frequentate dai ricchi orientali e dagli stranieri facoltosi durante
l'estate (i Luxi odiano abbronzarsi!).

I fiumi sono larghi, lenti e spesso navigabili. Il pi\`u
importante, il Maisef, attraversa Bahuney da Sud a Nord, nascendo dai
Simeinel e sfociando nel Golfo di Levante, nel Mare Interno, per un
totale di 1500 km.

La regione sud-orientale dell'Isola, inospitale e scarsamente
produttiva, ha una densit\`a di popolazione bassissima ed \`e abitata
da clan di nomadi. I Luxi sfruttano queste caratteristiche utilizzando
il territorio come sede di numerose prigioni e di alcuni campi di
lavori forzati per detenuti.

\Politica Lo Stato Elfico di Luxi \`e una monarchia costituzionale
ereditaria.  Il Re \`e una persona di grande prestigio, \`e comandante
in capo dell'esercito, detiene il potere esecutivo e svolge compiti
istituzionali.

Gli organi legislativi dello Stato sono il Consiglio dei Nobili e il
Consiglio degli Anziani. Il Consiglio dei Nobili \`e formato da 64
membri, ciascuno dei quali \`e Signore di un Territorio (sul quale
per\`o non esercita alcun potere), ed \`e presieduto dal Re. Il
mandato al Consiglio dei Nobili \`e ereditario e a vita e
strettamente legato al titolo nobiliare. Il consiglio degli Anziani
\`e composto da un massimo di 49 Senatori, eletti ogni 10 anni tra
tutti gli Elfi che hanno compiuto 100 anni e che hanno un avuto
reddito minimo di 60000 Corone nei 10 anni del mandato precedente.

Possono votare tutti coloro che nell'anno in corso hanno compiuto 25
anni e fanno parte di una famiglia che ha un reddito minimo di 2000
Corone al mese. I piccoli possidenti, gli artigiani e i ``poveri'' non
hanno diritto al voto, ma possono accedervi nel momento in cui,
migliorando le loro condizioni, arrivano al reddito minimo.

Il Consiglio degli Anziani detiene il potere legislativo, mentre il
consiglio dei Nobili e il Re hanno potere di Veto sulle leggi votate
dal Consiglio degli Anziani se queste sono state approvate da una
maggioranza inferiore ai due terzi dei votanti. I membri del consiglio
dei Nobili possono proporre le leggi ed esprimere pareri sulle
proposte del Consiglio degli Anziani, ma non possono votare per la
loro approvazione.

Il Re pu\`o nominare 5 ministri di sua fiducia: economia,
esteri, difesa, interni, giustizia, scegliendo tre membri del Consiglio degli
Anziani e due del Consiglio dei Nobili. I Ministri devono ricevere dopo la nomina
il Benestare del Consiglio degli Anziani a maggioranza relativa. 

Il Tribunale Maggiore \`e l'organo che presiede a tutte le
attivit\`a giudiziarie.  \`E composto da 25 Giudici eletti ogni 3
anni dai Giudici stessi tra tutti i Giudici del Regno che hanno
compiuto 100 anni. 

I Giudici sono le autorit\`a che hanno il compito
di interpretare le leggi, giudicare i crimini ed infliggere le
condanne. Qualunque cittadino Luxi di almeno 25 anni pu\`o diventare
Giudice seguendo un corso gratuito di 3 anni, al quale si accede
mediante un esame di ammissione. Durante il corso l'aspirante Giudice
deve esibire una condotta di vita integerrima e al termine deve
superare una difficilissima prova. L'aspirante giudice viene espulso
dal corso se lo si ritiene ``indegno''. Dopo 25 anni di attivit\`a
un Giudice pu\`o diventare istruttore ai corsi per Giudice.  Il
Tribunale Maggiore nomina i nuovi Giudici e toglie il mandato ai
Giudici che vengono giudicati ``indegni'', inoltre pu\`o giudicare e
condannare un membro del Consiglio dei Nobili, del Consiglio degli
Anziani, o addirittura il Re. I Consigli in seduta congiunta possono
espellere dal Tribunale Maggiore un Giudice che si \`e ritenuto
indegno del suo mandato. Il Ministro della Giustizia ha diritto di
parola ma non di voto a tutte le riunioni del Tribunale Maggiore. 

Dal punto di vista amministrativo, Luxi \`e diviso in 64 Territori,
ciascuno dei quali \`e retto da un Sindaco con poteri esecutivi
eletto ogni 11 anni e dal Signore del Territorio che ha soltanto
compiti di rappresentanza, oltre a far parte del Consiglio dei Nobili.

La Capitale dello Stato di Luxi, Pert, \`e la sede di tutti i
Consigli, nonch\'e del Palazzo del Re.

Le donne possono entrare a far parte di qualunque organo istituzionale
escluso il Consiglio dei Nobili e hanno gli stessi diritti degli
uomini di pari censo ed et\`a.

\Esercito L'esercito Luxi non \`e
particolarmente potente n\`e numeroso. Il suo compito principale \`e
quello di mantenere l'ordine pubblico collaborando con i Giudici. Le
citt\`a Luxi pi\`u importanti sono per\`o fortificate da spesse
mura.

I Luxi non sono amanti delle armi e dei combattimenti, preferendo in
caso di guerra affidarsi ai loro alleati. Praticano per\`o
volentieri (\`e quasi uno sport nazionale) la scherma, con la
Sciabola Elfica, il Fioretto, e la Spada.

 Capita spesso di trovare tra
le fila del piccolo esercito cavalieri di notevole bravura ed
eleganza, tanto che esiste una leggenda che narra che l'equitazione
\`e insegnata ai Luxi direttamente dal Dio dei Venti.  Esiste
inoltre un corpo speciale di spie e guastatori di cui si sa soltanto
che viene addestrato in maniera esemplare dai migliori maestri
Orientali e in cui \`e pressoch\'e impossibile entrare (si viene
scelti in seguito a segnalazione).  Di notevole importanza \`e la
Marina, che svolge in tempo di pace funzioni di polizia e controllo su
tutta la parte est del Mare Interno.

\Economia Le coste basse al settentrione e la presenza di fiumi in pianura
giustificano l'importanza dei settori economici del commercio navale,
della pesca, dell'agricoltura estensiva in prossimit\`a dei corsi
d'acqua. 

La flotta mercantile \`e importantissima sia per il commercio
interno che per quello estero. Nel Golfo di Levante si pescano pesci
azzurri e pesce pregiato, destinato prevalentemente al mercato
interno. L'artigianato copre interamente il fabbisogno dei trasporti
via acqua per quanto riguarda la produzione di imbarcazioni.  Gli
stessi Chire ammirano l'operato dei carpentieri Luxi, le cui barche e
navi sono famose per eleganza ed affidabilit\`a.

Fiorenti l'allevamento del baco da seta e degli ovini da lana, che
forniscono materia prima per la produzione di vestiti e tessuti
pregiati, che assieme alle barche sono la principale voce delle
esportazioni.

Gli allevatori Luxi
sono famosi per i loro cavalli di razza purissima, che spesso vengono
comprati a caro prezzo dai nobili, dai borghesi pi\`u ricchi e dai
Re, anche stranieri.

Tra gli altri prodotti: ortaggi, cereali (soprattutto grano), frutta,
vini pregiati, birra, liquori (anche esportati). 

I Luxi importano metallo grezzo e armi da
Terranova e dalle Terre Orientali, materiale magico da Dyshu Thandy,
macchinari da Loydi-Genya, spezie e profumi da tutto l'Arcipelago. Al
Re \`e destinata una parte del gettito fiscale, mentre i nobili
vivono delle rendite terriere.

\Magia La principale Accademia di Magia si trova a Silays ed \`e retta da
un Arcimago. Le scuole insegnate nel territorio sono Illusionismo, della
quale i Luxi sono maestri, ed Esistenza.

Poich\'e la retta di frequenza \`e molto elevata, alle Accademie
possono accedere soltanto i nobili e i borghesi ricchi (min. 0.6 di
fattore istruzione) oppure i meno abbienti che vengono presentati e
``sponsorizzati'' da un Nobile o da un Borghese, che paga loro i
corsi.  Per costoro \`e motivo di grande orgoglio e prestigio avere
un ``pupillo'' che si distingue nelle arti magiche. Per accedere ai
corsi \`e comunque necessario superare un esame di ammissione.

\Religione I Luxi sono generalmente indifferenti alle questioni religiose, tanto
che non esiste una vera e propria religione di stato. Molti di loro
professano comunque un culto che identifica in un Dio detto Mildril
l'essenza delle cose e l'origine di tutto. Le loro cerimonie sono di
fatto sedute di meditazione collettiva, presiedute da un Sacerdote,
durante le quali cercano di mettere in contatto l'essenza del loro
corpo con l'essenza dello spirito, quindi con il Mildril stesso,
cercando l'armonia Suprema e la Bellezza Cosmica. Non sono rare le
occasioni in cui piccole comunit\`a organizzano cerimonie di tipo
tantrico.

I Luxi tollerano tutte le altre religioni, tanto che \`e facile
trovare a fianco ad un Tempio di Mildril templi della Terra costruiti
dagli elfi Umbra oppure chiese Rileane. Le regole di tipo religioso
prevedono soltanto il rispetto per il prossimo e consigliano l'amore
libero, che aiuta ``a trovare l'armonia delle Cose e a Canalizzare lo
Spirito verso il Mildril'' come citato nel principale libro ``sacro''
chiamato ``Il Libro delle Essenze''.

La religione condanna
comunque i rapporti fra consanguinei; inoltre molti Elfi Luxi che sono
entrati in contatto con la religione Rileana non sono molto d'accordo
col principio dell'amore libero, specie quando hanno mogli giovani,
carine ed inquiete. 

\pinup{elfaluxi2.eps}{Elfa Luxi}

\Moda Tutti i Luxi tendono, per una questione di cultura, se ne hanno la
possibilit\`a, a curare molto il loro aspetto fisico,
l'abbigliamento, la pettinatura, il trucco, le movenze, il portamento.
Sono caratteristiche le loro parrucche bianche legate sulla nuca da un
fiocco nero o azzurro, il fondotinta bianco col quale i Luxi di
carnagione meno chiara si truccano il viso e la tintura rossa con la
quale mettono in risalto il disegno delle labbra. I ricchi (non
c'\`e praticamente distinzione tra nobili e borghesi) si vestono con
abbondanza di pizzi e merletti, con ampi colletti e maniche.

Gli elfi Luxi godono della fama, assolutamente giustificata, di
``belli'' ed ``esteti''. Coloro che ne hanno la possibilit\`a
coltivano infatti le abitudini del libertinaggio e della ``ricerca del
bello'' in senso assoluto, complice la religione. I ricchi organizzano
spesso feste con balli a cui partecipano personalit\`a di rilievo.
Non \`e raro che tali feste degenerino in qualcosa che agli occhi di
Levanti e Rileani pu\`o essere considerato scandaloso, ma che per gli
elfi Luxi \`e assolutamente normale. I Sacerdoti Rileani che vivono a
Bahuney tentano costantemente di opporsi a tali costumi, senza
riuscirvi, tanto sono radicate queste abitudini nella cultura e nella
mentalit\`a Luxi.

L'arte in tutti i suoi aspetti \`e tenuta in grande considerazione,
anche se per questioni economiche sono pochi coloro che possono
goderne appieno. I meno ricchi si accontentano di organizzare feste di
paese o nelle locande, al ritmo di danze popolari, durante le quali i
giovani e le giovani non ancora sposati formano coppie spesso
destinate a durare una sola notte.

Visto che
\`e semplice per un giovane o una giovane Luxi fare ``esperienza'',
i Luxi sono famosi per essere bravi amatori. L'esperienza in questo
settore \`e tenuta in grande considerazione per la scelta dello
sposo o della sposa. Una volta sposati, i Luxi diventano monogami, ma
le ``scappatelle'' sono quasi sempre tollerate se giustificate (es.
malattia di un partner, viaggi, lontananza). 

Gli sport pi\`u
praticati dai ricchi sono la scherma e l'equitazione; i meno
fortunati organizzano tornei di combattimento col bastone, gare di
corsa, gare a cavallo senza sella.

Importante la presenza a Bahuney di una minoranza di Elfi Umbra che
spesso lavorano come servi o si mischiano alla popolazione meno
abbiente, tanto che sono disprezzati dai ricchi, per i quali ``Umbra''
\`e sinonimo di ``Servo'' o ``Pezzente''.

La Lingua ufficiale \`e l'Elfico, nella sua caratteristica variante
dal suono dolce e flautato.


\Fisico I Luxi di sesso maschile sono alti in media 1.80, magri e
asciutti.  Sono di aspetto delicato e hanno lineamenti molto dolci.
Generalmente hanno un portamento elegante e felino, tanto che agli
occhi degli altri abitanti dell'Arcipelago risultano ``effeminati''.

Le donne Luxi hanno una figura molto sottile e sinuosa, un modo di
muoversi molto sensuale ed emanano un particolare sex-appeal. Sia
gli uomini che le donne hanno le caratteristiche orecchie a punta, la
carnagione chiarissima, quasi albina, gli occhi di un azzurro
brillante molto chiaro e i capelli scuri e lisci.

Come per tutti gli elfi, il loro invecchiamento si svolge in due fasi:
durante la prima, che va dalla nascita a circa 25-26 anni di et\`a,
crescono e si sviluppano come gli esseri umani; dopodich\'e, nella
seconda fase, il loro invecchiamento rallenta tanto che 5 dei loro
anni corrispondono circa ad 1 anno di et\`a degli umani, consentendo
loro di raggiungere facilmente un'et\`a di 200-250 anni. 

La loro longevit\`a \`e compensata da una scarsa fertilit\`a: una
famiglia elfica arriva a procreare al massimo 2-3 figli nel corso
della vita.

Gli elfi sono in grado di vedere al buio purch\'e questo non sia
totale (realizzando un normale TOSS).

\minmaxelfi

\subsubsection{Abilit\`a di razza}

\begin{abilist}
\item Resistenza al caldo 4
\item Visione Notturna
\item Sedurre 3
\end{abilist}

\subsubsection{Classi sociali e abilit\`a di Classe
  Sociale}

\classesoc{01-09}{Contadini}{0.3}

Coltivano campi ed allevano bestiame.  Producono cibo e vestiti per la
propria famiglia e possono guadagnare al netto delle spese di
sostentamento da 1 a 20 scudi al mese.

\begin{abilist}
\item Ammaestrare
\item Carpenteria
\item Cavalcare bipedi o quadrupedi
\item Conoscere animali
\item Conoscere piante
\item Lingua levante
\item Lingua madre
\item Resistenza alla fatica
\end{abilist}

\classesoc{10-19}{Pescatori}{0.3}

Vivono della pesca e dello sfruttamento del mare e dei suoi proventi.
Producono cibo per la propria famiglia e possono guadagnare al netto
delle spese di sostentamento da 1 a 20 scudi al mese.

\begin{abilist}
\item Carpenteria
\item Conoscere animali
\item Controllo respirazione
\item Fare nodi
\item Lingua levante
\item Lingua madre
\item Navigare
\item Nuotare
\item Orientamento
\end{abilist}

\classesoc{20-29}{Cacciatori}{0.3}

Vivono della caccia dei suoi proventi. Producono cibo e vestiti per la
propria famiglia e possono guadagnare da 2 a 8 corone al mese.

\begin{abilist}
\item Allarme
\item ADT o ADL
\item ATC o ABC
\item Conoscere animali
\item Conoscere piante
\item Fare nodi
\item Fiutare
\item Imitare
\item Individuare tracce
\item Lingua levante
\item Lingua madre
\item Muoversi in silenzio
\item Trappole
\item Udire
\end{abilist}

\classesoc{30-59}{Professionisti}{0.5}

Praticano i pi\`u svariati mestieri e vivono di essi. Possono
essere: artigiani, giuristi, cartografi, erboristi, carpentieri,
medici, macellai ecc. Il loro introito mensile pu\`o variare da 5
a 40 corone al mese.

\begin{abilist}
\item Abilit\`a professionale \#1
\item Abilit\`a professionale \#2
\item Abilit\`a professionale \#3
\item Lingua levante
\item Lingua madre
\item Altre abilit\`a a scelta del Master
\end{abilist}

\classesoc{60-89}{Benestanti}{0.6}

Praticano molteplici mestieri. Hanno grandi esercizi commerciali, come
alberghi, gioiellerie, medi allevamenti e via dicendo, che gli
forniscono gran parte del reddito, variabile tra 15 e 300 corone al
mese.

\begin{abilist}
\item ATL o ADT
\item Calcolare
\item Cavalcare bipedi o quadrupedi
\item Danzare
\item Galateo
\item Lingua levante
\item Lingua madre
\item Lingua straniera
\item Mitologia
\item Religione 
\item Storia
\item Generale
\item Trattare
\item Valutare
\end{abilist}

\classesoc{90-95}{Borghesi ricchi}{0.7}

Sono proprietari di molteplici attivit\`a, in cui investono il proprio
capitale. Sono perlopi\`u banchieri, armatori, grossi allevatori,
possessori di miniere, mercanti di schiavi, ecc. Ricavano al mese da
100 a 3000 corone.

\begin{abilist}
\item ATL
\item ADT
\item Calcolare
\item Cavalcare bipedi o quadrupedi
\item Danzare
\item Galateo
\item Lingua levante
\item Lingua madre
\item Lingua straniera
\item Mitologia
\item Religione
\item Storia Generale
\item Trattare
\item Valutare
\end{abilist}

\classesoc{96-99}{Nobili}{0.8}

Hanno estesi possedimenti terrieri ed attivit\`a collegate ad essi.
Possono ricavare, a seconda del rango, da 500 a 3000 corone al mese.

\begin{abilist}
\item ATL
\item ADT
\item Calcolare
\item Cavalcare bipedi o quadrupedi
\item Danzare
\item Galateo
\item Lingua levante
\item Lingua madre
\item Lingua straniera
\item Mitologia
\item Oratoria
\item Religione
\item Storia Generale
\end{abilist}

\classesoc{100}{Re}{0.8}

Una parte dell'erario \`e destinata a foraggiare la famiglia reale.
Il re (e la sua famiglia) hanno un introito mensile fisso di 3000
corone al mese.

\begin{abilist}
\item ATL
\item ADT
\item Calcolare
\item Cavalcare bipedi o quadrupedi
\item Danzare
\item Galateo
\item Lingua levante
\item Lingua madre
\item Lingua straniera
\item Mitologia
\item Oratoria
\item Religione
\item Storia Generale
\end{abilist}


\subsection{Umbrosa}

\Stendardo{stendardo_umbra.eps}

Gli elfi Umbra vivono nel territorio di Umbrosa

\Geografia Il territorio di Umbrosa copre la quasi totalit\`a dell'isola di Malagana. A
Nord confina con le terre brulle (Tehramala) dove sono dislocate le trib\`u
orchesche.

Le coste si presentano alte e frastagliate in tutta la zona Nord-Ovest
e Nord-Est mentre degradano pi\`u dolcemente verso Sud (da Capo
Ponente ad Est e Punta Mow-Lentys a Ovest sino alla Baia delle Grandi
Corna). Il territorio si innalza progressivamente da Sud a Nord ma non
supera mai la quota massima di 837 metri, raggiunti da Punta Alta.

Il clima \`e caratterizzato da inverni rigidi e nevosi ed estati
tiepide, ma a Sud una corrente calda (Kahl-Henty) concede temperature
pi\`u miti.

Tutto il territorio \`e coperto dalla grande foresta di Segram
(alberi dal tronco largo coi rami che si estendono orizzontalmente)
che formano una sorta di pianura di fronde, in alcune zone
cos\`{\i} fitta da poterci camminare.  Queste pianure si sviluppano ad
altezze differenti, a circa 10 metri l'una dall'altra, fino a
raggiungere punte di 50-60 metri.

Le citt\`a vengono erette sulle pianure dei Segram e spesso distinte
in citt\`a alte e citt\`a basse a seconda del livello su cui sono
dislocate. 

La capitale Jar-Gelar sorge nell'entroterra sulla riva del
fiume Maj-Jiahn. Le citt\`a di Sey-yan, Mut, e Felyan sono le citt\`a
di frontiera che dividono Umbrosa dalla terra degli Orchi.  Altre
citt\`a molto importanti sono Kella, Elra, Fey, Partica e Syray.

Circa al centro della Foresta cresce il Grande albero della Vita
(Sahmatha) il cui tronco ha 50 metri di diametro alla base ed \`e alto
pi\`u di 120 metri.

I corsi d'acqua che sfociano su ogni costa sono di piccola portata. Il
Maj-Jiahn, il grande fiume che nasce nei monti Thyan (nella terra
degli orchi) e che sfocia al centro della Baia delle Grandi Corna \`e
il fiume pi\`u importante ed \`e quasi interamente navigabile da
Jar-Gelar fino alla foce. Tutti i fiumi sono generalmente calmi ma
durate la stagione delle piogge e col disgelo divengono tumultuosi ed
a causa delle piene del Maj-Jiahn si verificano delle inondazioni che
fertilizzano il terreno.

\Politica Umbrosa \`e una sorta di repubblica democratica,
amministrativamente divisa in 6 zone, i cui organi principali sono il
Consiglio degli Anziani e l'assemblea del Popolo.  Il Consiglio \`e
composto di 40 membri e il pi\`u anziano di essi \`e il presidente.
Viene eletto ogni 15 anni fra coloro che hanno pi\`u di 100 anni di
et\`a.

Il Consiglio avoca a s\`e il pieno potere giudiziario quale sommo
tribunale. Il Consiglio nomina un giudice primario per zona con il
compito di amministrare la giustizia.  

In realt\`a i Giudici sono anche i ``creatori'' della legge poich\'e
esistono pochissime norme scritte per lo pi\`u religiose. Essi quindi
giudicano in base all'equit\`a ed ai precedenti giudiziari. 
Le udienze avvengono di solito a porte chiuse e si \`e ammessi
mediante iscrizione ad un'apposita lista. Le decisioni dei giudici
sono appellabili presso il Consiglio degli Anziani che per ogni
diatriba nomina quattro giudici fra i suoi membri.

L'assemblea \`e composta da 60 membri, 10 per ogni zona, ed \`e eletta
ogni dieci anni fra i maggiori di 45 anni. L'assemblea detiene il ruolo
diplomatico ed il potere Amministrativo Centrale, che esercita demandando 
l'amministrazione di ogni zona ad un commissario.
All'interno di ogni zona i 10 membri dell'assemblea sono eletti a
suffragio universale. L'assemblea nomina anche i generali posti a capo
dell'esercito e, in caso di guerra, i suoi componenti sono ridotti da
60 a 6 (uno per zona).

\Esercito L'esercito Umbra \`e una sorta di corpo d'\'elite poich\'e
preposto alla sola difesa dell'isola e quindi addestrato a combattere
nella foresta.  Esso \`e famoso per la bravura dei suoi arcieri.

\`E diviso in 6 bracci, uno per ogni zona, ognuno dei quali \`e
comandato da un generale eletto dai 10 membri zonali dell'assemblea.
Ogni braccio \`e diviso in 3 mani (fanti, arcieri e cavalieri) ognuna
con un proprio comandante.

Direttamente dall'assemblea dipendono la flotta, composta da due
onde, ognuna comandata da un ammiraglio e il corpo dei Cavalieri del
Cielo che usano i Reptis come cavalcatura e sono solitamente armati
con gli archi elfici.

Chiunque pu\`o entrare a far parte dell'esercito se \`e abbastanza
in gamba. L'unico corpo elitario \`e quello dei Cavalieri del Cielo
poich\'e \`e necessario un grande affiatamento tra cavalcatura e
cavaliere per cui i Reptis vanno addestrati da quando sono cuccioli.
Entrare in armi \`e motivo di orgoglio per il prescelto; a nessuno
\`e precluso il raggiungimento di un alto grado.

Ogni generale, eletto a vita, provvede ad istituire scuole d'armi in
ogni zona al fine di addestrare i giovani all'uso di: Spada, Spada
Bastarda, Alabarda, Arco Elfico e alle tecniche di ``combattimento
nascosto''.

In tempo di pace l'esercito svolge mansioni di ordine pubblico sotto
il comando del primo giudice. Il compito di difendere le citt\`a dagli
attacchi dal suolo \`e affidata al corpo dei Ranger, che
tramandano il mestiere di padre in figlio. Questa carica e quella di sciamano
sono le sole cariche ereditarie ammesse, per via della lunga preparazione
necessaria a svolgere il compito e per il fatto che solo pochi
accetterebbero di vivere al suolo.

Gli arcieri sono spesso assoldati a caro prezzo, come mercenari, dagli
eserciti stranieri.

Sotto la citt\`a di Kella, a livello del Suolo, si trova l'accademia
di Arti Marziali pi\`u importante dell'isola.

\Economia Gli elfi Umbra sono un popolo quasi totalmente autarchico,
le uniche esportazioni riguardano una sorta di lattice bianco e
gommoso (Sute) secreto dai Segram ed utilizzato per impermeabilizzare
le imbarcazioni; i funghi la cui raccolta \`e incentivata dalla loro
proliferazione che se non fosse tenuta sotto controllo impoverirebbe
troppo la foresta; e i gioielli prodotti a partire dall'oro presente
nei fiumi, di cui gli Umbra sono degli abili lavoratori.

L'economia si basa sugli scambi interni e l'agricoltura \`e
incentrata sulla semplice raccolta. 

Ogni 6 mesi
tutti gli Umbra versano alla Gabella Zonale una tassa corrispondente a
circa il 30\% degli introiti semestrali. Un terzo della stessa viene
inviato all'assemblea per l'amministrazione centrale e per gli
eserciti non dipendenti dalle Zone. I restanti due terzi vengono
amministrati dal delegato di Zona, che li ripartisce, a seconda delle
esigenze, fra amministrazione zonale ed esercito zonale.

Un'altra tassa, nella misura del 5\% semestrale, viene versata
direttamente al consiglio degli anziani per l'amministrazione della
giustizia.

Gli unici cittadini a non pagare le imposte sono i soldati, gli
sciamani, i Ranger ed i primi giudici, ai quali viene versata una
somma mensile gi\`a decurtata delle imposte. 

Gli esportatori sono inoltre tenuti al pagamento di una tassa di
carico a seconda del peso e del valore della merce trasportata.

Per quanto non esista una sostanziale sperequazione fra le varie
classi sociali \`e da notare che le famiglie dei raccoglitori di Sute
e dei gioiellieri sono le pi\`u ricche ed importanti e riescono ad
influire non poco sulla politica del regno.

Senza influsso politico ma importanti per l'economia sono i
fabbricanti di archi Elfici rinomati in tutte le terre conosciute.

L'unica importazione degna di nota \`e quella del metallo per armi ed
armature dalla terra degli orchi poich\'e gli Umbra non possiedono
giacimenti minerari particolarmente ricchi.

\Religione Gli Umbra riconoscono come divinit\`a la Grande Madre
Terra, creatrice di tutta la vita e dispensatrice di morte. Questa
fede li rende molto attaccati e rispettosi della natura e di tutte le
creature viventi (ad eccezione dei Krenay poich\'e si dice che siano
sfuggiti ``all'abbraccio della Dea'').

Le altre religioni sono tollerate. In alcune citt\`a costiere
occidentali esse sono professate ma risentono notevolmente dell'influsso
dell'antica religione divenendo cos\`{\i} delle specie di ibridi. Molte
regole religiose sono riprese dal Consiglio e trasformate in leggi.

\pinup{madre.eps}{Icona umbra}

Il Capo della Religione \`e il Signore dell'albero della Vita
(Shay-Han) che trasmette la sua carica ad un prescelto, prima di
morire, con un rito sacro che riunisce tutti gli Umbra ai piedi
dell'albero della vita.

Altra occasione di riunione ai piedi dell'albero \`e la festa di
primavera che si tiene il primo giorno di Kalhentee. Generalmente le
festivit\`a sono legate al mutare delle stagioni e le cerimonie
vengono celebrate dallo sciamano di ogni citt\`a.

Quando gli Umbra uccidono un animale per cibarsene \`e usanza recitare
una preghiera.

Quasi tutti gli Umbra si recano almeno una volta nella vita presso la
Roccia (nell'isola di Norda) dove si dice sia stato sepolto Shalaran,
portavoce degli Elfi quando ancora Luxi e Umbra erano un solo popolo
e principale fautore della pace con gli Umani.  

La carica di Sciamano, il Sacerdote del culto della Grande Madre, \`e
ereditaria.

\Magia Ad Umbrosa esistono 3 scuole di magia: la Stregoneria,
l'Elementalismo e l'Esistenza. Esse sono insegnate da 3 accademie
aderenti al Conclave ma non direttamente dipendenti da esso.

Le 3 accademie sono situate nelle citt\`a di Mut e Kella e nella
capitale Jar-Gelar. L'accademia di Mut \`e retta da un Arcimago.
Ammessi all'accademia sono solitamente i benestanti (Fattore
Istruzione 0.5) ma ogni anno avviene una gara fra tutti gli aspiranti
che non possono permettersi di pagare le rette ed ogni accademia ne
sceglie 3.

La magia viene insegnata anche dalle famiglie, previa comunicazione
alle autorit\`a. L'insegnamento in famiglia non permette comunque di
raggiungere un alto grado di abilit\`a magica. 

La scuola di Stregoneria \`e considerata la pi\`u prestigiosa. La
scuola della Necromanzia, invece, \`e bandita dal Regno.  Chi venisse
scoperto nel suo utilizzo sarebbe sottoposto a giudizio del consiglio;
la massima pena che pu\`o essere inflitta \`e l'esilio.

\Moda Gli Umbra sono un popolo molto sobrio e molto fiero della sua
sobriet\`a.  Le occasioni mondane sono ben poche e non capita spesso
che si organizzino feste, eccezion fatta per i matrimoni.

Le classi sociali sono di fatto chiuse e le attivit\`a sono tipiche
per ogni classe. Sono pochi coloro che possono concedersi grandi
svaghi e questi sono immancabilmente i ricchi. Essi di solito
organizzano la caccia notturna al Krenay, per cui \`e necessario
essere guidati da un Ranger. Le femmine ritengono, forse a ragione,
che sia un passatempo stupido, data la pericolosit\`a di queste
lucertole umanoidi. Non \`e infatti infrequente che in queste battute
di caccia, che si svolgono al suolo, qualche ricco elfo ci lasci la
pelle.

Il ceto medio preferisce trascorrere le nottate in osteria a bere e ad
ascoltare musici girovaghi. Fra i pi\`u giovani \`e frequente un gioco
chiamato ``Mamakowa'' dove due o pi\`u persone hanno il compito di
trovare gli altri partecipanti, nascosti, prima che questi raggiungano
la ``casa'' dei cercatori.

Le donne, nel loro tempo libero, si cimentano in un gioco il cui scopo
\`e cercare di far cadere dei sassi colorati (nero, rosso e verde)
all'interno di caselle, ciascuna delle quali avente un determinato
punteggio.

Gli Umbra sono un popolo sessualmente libero, anche se strettamente
monogamo. Un Umbra, nel corso della sua pur lunga vita, genera al
massimo 2-3 figli, data la scarsa fertilit\`a delle coppie.

Gli abiti, scuri, sono ricavati da pelli, lana o fibre vegetali
grezze. \`E frequente l'uso del cappuccio, soprattutto in inverno.

Sia le donne che gli uomini portano pantaloni; se la temperatura lo
permette le donne indossano corte gonne e gli uomini perizomi che
arrivano fino a met\`a della coscia.

\pinup{haryell.eps}{Haryell, Elfa Umbra}

\Fisico Gli Umbra hanno statura medio-alta (altezza media 1.75 m), e
non sono eccessivamente robusti. Il loro aspetto \`e gradevole, hanno
solitamente capelli scuri e occhi chiari, la carnagione naturalmente
chiara (ma \`e diffusa l'usanza di abbronzarsi).

Sono in grado di vedere al buio purch\'e questo non sia totale (Con un
TOSS).

Il loro ciclo vitale si svolge in due fasi: durante la prima, che va
dalla nascita a circa 25-26 anni di et\`a, crescono e si sviluppano
come gli esseri umani; dopodich\'e, nella seconda fase, il loro
invecchiamento rallenta tanto che 5 dei loro anni corrispondono circa
ad 1 anno di et\`a degli umani. Gli elfi facilmente arrivano ad
un'et\`a di 200-250 anni.

Come i Luxi, la loro prolificit\`a \`e molto bassa; essi arrivano ad
avere massimo due o tre figli per coppia.

La lingua ufficiale \`e l'Elfico.

\minmaxelfi

\subsubsection{Abilit\`a di razza}
\begin{abilist}
\item Resistenza al freddo 6
\item Scalare 3
\item Visione Notturna
\end{abilist}

\subsubsection{Classi sociali e abilit\`a di Classe sociale}

\classesoc{01-04}{Contadini}{0.3}

Sfruttano le risorse naturali dei boschi. Producono cibo e vestiti per la
propria famiglia e possono guadagnare al netto delle spese di
sostentamento da 1 a 20 scudi al mese.

\begin{abilist}
\item Carpenteria
\item Cavalcare bipedi
\item Conoscere animali
\item Conoscere piante
\item Lingua Levante
\item Lingua Madre
\item Resistenza alla fatica
\end{abilist}

\classesoc{05-19}{Pescatori}{0.3}

Vivono della pesca e dello sfruttamento del mare e dei suoi proventi.
Producono cibo per la propria famiglia e possono guadagnare al netto
delle spese di sostentamento da 1 a 20 scudi al mese.

\begin{abilist}
\item Carpenteria
\item Conoscere animali
\item Controllo respirazione
\item Fare nodi 
\item Lingua levante
\item Lingua madre
\item Navigare
\item Nuotare
\item Orientamento
\end{abilist}

\classesoc{20-35}{Cacciatori}{0.3}

Vivono della caccia e dei suoi proventi. Producono cibo e vestiti per
la propria famiglia e possono guadagnare da 2 a 8 corone al mese.

\begin{abilist}
\item Allarme
\item ADT o ADL
\item ATC
\item Conoscere animali
\item Conoscere piante
\item Fare nodi
\item Fiutare
\item Imitare
\item Individuare tracce
\item Lingua levante
\item Lingua madre
\item Muoversi in silenzio
\item Trappole
\item Udire
\end{abilist}

\classesoc{36-39}{Sciamani}{0.7}

Sono custodi della memoria e dei riti religiosi dei clan. Vivono di
ci\`o che viene loro offerto dal popolo e della raccolta di erbe e
dei loro prodotti. Il loro introito pu\`o variare da 1 a 50 scudi al
mese.

\begin{abilist}
\item Alchimia
\item Astronomia
\item Astrologia
\item Cavalcare bipedi
\item Conoscere piante
\item Lingua Madre
\item Lingua Levante
\item Medicina
\item Mitologia
\item Religione
\end{abilist}

\classesoc{40-59}{Ranger}{0.5}

Vivono di caccia e raccolta e della paga data loro dal Commissario di
Zona. I loro introiti variano da 2 a 20 corone al mese.

\begin{abilist}
\item Allarme
\item Alchimia
\item ATC
\item ADT
\item Conoscere animali
\item Conoscere piante
\item Fare nodi
\item Fiutare
\item Imitare
\item Individuare tracce
\item Individuare trappole
\item Lingua levante
\item Lingua madre
\item Muoversi in silenzio
\item Orientamento
\item Trappole 
\item Udire
\end{abilist}

\classesoc{60-85}{Professionisti}{0.5}

Praticano i pi\`u svariati mestieri e vivono di essi. Possono
essere: artigiani, giuristi, cartografi, erboristi, carpentieri,
medici, macellai ecc. Il loro introito mensile pu\`o variare da 5
a 30 corone.

\begin{abilist}
\item Abilit\`a professionale \#1
\item Abilit\`a professionale \#2
\item Abilit\`a professionale \#3
\item Lingua levante
\item Lingua madre
\item Altre abilit\`a a scelta del Master
\end{abilist}

\classesoc{86-95}{Benestanti}{0.6}

Praticano molteplici mestieri. Hanno grandi esercizi
commerciali, come alberghi, gioiellerie, medi allevamenti e via
dicendo, che gli forniscono gran parte del reddito, variabile tra 15 e
150 corone al mese.

\begin{abilist}
\item ATL o ADT
\item Calcolare
\item Cavalcare bipedi
\item Danzare
\item Lingua levante
\item Lingua madre
\item Lingua straniera
\item Religione
\item Storia
\item Generale
\item Trattare
\item Valutare
\end{abilist}

\classesoc{96-100}{Borghesi ricchi}{0.7}

Sono proprietari di molteplici attivit\`a, in cui investono il proprio
capitale. Sono perlopi\`u banchieri, armatori, grossi allevatori,
possessori di miniere, mercanti di schiavi, ecc.  Ricavano al mese da
100 a 2000 corone.

\begin{abilist}
\item ATL
\item ADT
\item Calcolare
\item Cavalcare bipedi
\item Danzare
\item Lingua levante
\item Lingua madre
\item Lingua straniera
\item Mitologia
\item Religione
\item Storia Generale
\item Trattare  
\item Valutare
\end{abilist}

\section{Gli Gnomi}

\subsection{Loydi-Genya}

\Stendardo{stendardo_gnomi.eps}

Gli gnomi sono divisi in due etnie, Beedha e
Kunetha, che condividono il territorio di Loydi-Genya.

\Geografia Loydi-Genya \`e abitata da due etnie di gnomi: i \textbf{Beedha} e i
\textbf{Kunetha}.  Essa \`e un'isola molto particolare, la sua morfologia \`e
stata modificata dai suoi operosi abitanti moltissimo tempo fa. Ora
\`e costituita dall'enorme altopiano di Tahula che la occupa quasi
interamente. Solo le coste non hanno subito modifiche ad eccezione
della zona Sud-Est dove sorge la principale citt\`a Loy (abitata quasi
totalmente dai Kunetha). Essa \`e l'unica citt\`a sull'isola
principale nonch\'e unico porto, e unica citt\`a con un nome normale.
Le altre citt\`a F12, B32, K47, e la capitale A1, sorgono ognuna su
un'isola dell'Arcipelago a Sud-Est.

Le isole, fra loro e con la terraferma, sono unite da
mastodontici ponti fatti in un materiale sconosciuto, e la cui
costruzione ogni gnomo attribuisce alla propria famiglia. I ponti sono
alti circa 100 metri s.l.m., larghi 40 metri e ricoperti da una fitta
vegetazione che ne riduce notevolmente la carreggiata (circa 5 m).

Al centro dell'altopiano vi \`e un vastissimo lago perfettamente e stranamente
circolare da cui si dipartono nove fiumi Vyl, Ta, Nuhra, Laoh, Ossy,
Teesy, Vehra, Mowra, e Laska. Questi precipitano dall'altopiano da
un'altezza di 140 metri e formano dei laghi cos\`{\i}  piccoli da far
supporre che l'acqua scenda direttamente nel sottosuolo. Sia
l'altopiano che le coste sono coperti da una fitta foresta pluviale.

\Politica Loydi-Genya \`e una sorta di repubblica democratica
fortemente burocratizzata.  Per quanto riguarda i Beedha, gli
organi principali sono le due Camere: la camera del commercio con 100
membri eletti a suffragio universale fra le 1000 famiglie di
commercianti; e la camera degli inventori composta da 10 membri
volontari. Nel caso non vi siano volontari o ve ne siano troppi
(ipotesi pressoch\'e impossibile, dato che a nessun inventore
interessa la politica) la camera del Commercio provvede a integrare o
a tagliare l'altra camera.

Le decisioni sono prese dalla Camera del Commercio e devono avere il
benestare dell'altra camera (che di solito approva qualunque mozione o
progetto senza neanche leggerlo).  La Camera del Commercio affida poi
l'attuazione del programma a Commissioni; solitamente in numero di 10:
Edilizia, Trasporti, Idrica, Esteri, Istruzione, Relazioni Interne,
Mineraria, Tessile, Caccia e Pesca, Brevetti. Ogni Commissione \`e a
sua volta divisa in Sottocommissioni, che saranno a loro volta divise
in Sotto-sottocommissioni, e via dicendo. Da menzionare sono: la
commissione Brevetti che deve decidere a chi appartiene la paternit\`a
di un'invenzione contesa fra pi\`u gnomi (caso abbastanza
frequente) e la Commissione Relazioni interne che gestisce i rapporti
con i Kunetha a cui propone i programmi che li riguardano.

Il consiglio dei Kunetha, composto dai Capoclan dei 20 clan, deve
esprimere il suo parere, che non \`e comunque vincolante.  Inoltre il
consiglio elegge tre candidati fra i suoi membri, tra i quali la
camera del Commercio elegger\`a il Supremo Comandante Militare che
riunir\`a a s\`e tutti i poteri in caso di guerra.

I Kunetha sono divisi in 20 Clan che abitano la citt\`a di Loy, sono
fortemente uniti e le lotte fra clan sono rare. Sono dotati di una
lungimiranza sconosciuta ai Beedha ma sono considerati una sorta di
reietti.

\Esercito L'esercito \`e il pi\`u malandato di tutte le terre
conosciute, ma \`e tenacissimo.  Non vi \`e un vero addestramento
militare, ma tutti sanno usare le armi da taglio corte: daga, pugnale,
stiletto.  

Ci\`o che rende ``temibile'' l'esercito gnomico \`e la presenza di
stranissimi macchinari bellici (carri armati, carri anfibi, navi
corazzate a vapore, lanciafrecce a ripetizione ed armi da fuoco) che
rendono imprevisto, anche per gli stessi gnomi, l'esito della
battaglia.

Il Comando \`e affidato al supremo comandante ai suoi due generali che
scelgono i capi delle sette divisioni di: fanteria, cavalleria,
artiglieria, corazzate anfibie, corazzate terrestri, corazzate marine
e guastatori. Dato lo scarso utilizzo per pochi conflitti che hanno
visto partecipi gli gnomi, tali macchine presentano spesso delle
disfunzioni pi\`u o meno catastrofiche.

\Economia I Beedha commercianti producono e scambiano pressoch\'e
qualunque tipo di merce, dai generi alimentari ai gioielli, utilizzando
le macchine degli inventori per passare dalla materia prima al
prodotto finito.  Allevano principalmente bovini, ovini e fortral da
cui ricavano pelle, latte e carne. I Kunetha allevano suini.  \`E
fiorente anche la gestione di bordelli.

Le esportazioni riguardano tutte le merci ma in maggior misura l'erba
Cucha (una droga allucinogena molto gradita ai Chire), tabacco, the e
alcune invenzioni (ad esempio il ``Su e Gi\`u'' della citt\`a di
Rylex). Le importazioni sono limitate ai metalli.

Le invenzioni degli gnomi, sorprendentemente, funzionano, anche se a
volte sono molto grandi e il loro utilizzo molto complicato. Quelle
pi\`u complesse rimangono invece al livello di progetto perch\'e
vengono considerate irrealizzabili anche per la mancanza dei
materiali dotati delle caratteristiche meccaniche richieste dal progetto.

Alcune invenzioni particolarmente complicate possono essere messe in
funzione solo dagli gnomi, nonostante neanche loro sappiano come.

\pinup{gnomo_beedha.eps}{Gnomo Beedha al lavoro}

\Religione Gli gnomi sono atei, tutt'al pi\`u agnostici; uno gnomo
che professi una fede religiosa viene solitamente deriso ed
emarginato.

I morti vengono
cremati in grossi inceneritori di rifiuti presenti in ogni citt\`a,
chiamati Skalhatoyow.
 
\Magia La magia viene insegnata in accademie, aderenti al Conclave,
aperte gratuitamente a tutti i ceti sociali e a tutte le et\`a (non
\`e raro trovare qualche gnomo anzianotto che frequenta le accademie).

Gli insegnanti vengono scelti dal Conclave, solitamente fra maghi
gnomi dato che le altre razze trovano difficolt\`a ad abituarsi al
modello scolastico seguito dagli gnomi.  Le accademie impartiscono gli
insegnamenti di Demonologia, Illusionismo, Esistenza.

\Moda Gli gnomi vestono abiti pratici leggeri dai colori sgargianti. I
Kunetha sono soliti portare collane e orecchini vistosi. Le citt\`a
sono costruite su pi\`u livelli con passeggiate coperte e ponti di
corda ad unire livelli e settori. Il ``Su e gi\`u'' \`e il mezzo pi\`u
usato per cambiare livello.

Le strade sono illuminate da strani globi che emanano una luce color
ambra, e che non si possono n\'e spegnere n\'e rompere. Questi
globi si accendono automaticamente all'imbrunire; nessuno gnomo sa
come funzionino ma nessuno vuole ammetterlo.

I giovani gnomi vengono istruiti dalla propria famiglia o dal clan
secondo le proprie inclinazioni. 

Gli gnomi (e le gnome) apprezzano moltissimo la bellezza ``esotica''
delle razze umana ed elfica.

I passatempi preferiti sono: il gioco d'azzardo, la frequentazione dei
bordelli sia per gnomi che per gnome, la morra (gioco molto sentito,
che causa delle incomprensioni che sovente sfociano in accoltellamenti
non punibili come crimini) e lo Skowdee-Askhu-Ryow detto
anche ``buio matto'' di cui esiste addirittura un campionato. 

Questo gioco si svolge all'interno di un campo circolare di 10 metri
di raggio coperto da una cupola per far si che vi sia un buio totale.
All'interno del campo da gioco 50 gnomi devono cercarsi, catturarsi e
tramite prese dolorose costringere l'avversario alla resa. Colui che
si arrende perde un punto in classifica, chi costringe alla resa ne
guadagna uno. Il gioco prosegue finch\'e un solo gnomo non rimane in
piedi. Il campo da gioco ha una peculiarit\`a, ogni 10 eliminazioni si
restringe di 2 metri. La lingua ufficiale \`e lo Gnomico.

\Fisico Gli gnomi sono alti massimo 130 cm e pesano circa 30-40 chili.
Le femmine sono leggermente pi\`u basse e snelle. Hanno carnagione
olivastra e occhi e capelli scuri.

\subsection{Gnomi Beedha}

\minmaxgnomi

\subsubsection{Abilit\`a di razza}

\begin{abilist}
\item Resistenza al caldo 7
\item Blaterare 10
\end{abilist}

\subsubsection{Classi sociali e abilit\`a di Classe sociale}

\classesoc{1-69}{Inventori}{0.7}

Gli inventori producono, oltre a quantit\`a industriali di oggetti
pi\`u o meno inutili, anche parte dei prodotti necessari alla loro
sussistenza e parte delle materie prime. Hanno un introito mensile
variabile tra 5 e 500 corone mensili.

\begin{abilist}
\item ATC 
\item Abilit\`a professionale \#1
\item Abilit\`a professionale \#2
\item Architettura
\item Balistica
\item Calcolare
\item Carpenteria
\item Geologia
\item Ingegneria
\item Lingua levante
\item Lingua madre
\item Storia Generale
\end{abilist}

\classesoc{70-100}{Commercianti}{0.7}

Si preoccupano del commercio interno ed estero scambiando ogni genere
di prodotto, comprese le improbabili macchine prodotte dagli
inventori. l loro introito mensile pu\`o variare da 5 corone a 4
piastre al mese.

\begin{abilist}
\item ATC
\item Calcolare
\item Cartografia
\item Cavalcare bipedi
\item Geografia
\item Individuare menzogne
\item Lingua levante
\item Lingua madre
\item Lingua straniera \#1
\item Lingua straniera \#2
\item Raggirare
\item Trattare
\item Valutare (a scelta)
\end{abilist}

\subsection{Gnomi Kunetha}

\minmaxgnomi

\subsubsection{Abilit\`a di Razza} 
\begin{abilist}
\item Udire 5
\item  Blaterare 10
\end{abilist}

\subsubsection{Classi sociali e abilit\`a di Classe sociale}

\classesoc{1-24}{Mendicanti}{0.2}

Si procurano da vivere chiedendo l'elemosina e ricavano
proporzionalmente alla bont\`a della gente.

\begin{abilist}
\item ATC
\item Astrologia o Lettura della mano
\item Cantare o Suonare o Danzare
\item Geografia
\item Lingua levante
\item Lingua madre
\item Raggirare
\item Resistenza al freddo
\item Resistenza alle malattie
\item Travestirsi
\end{abilist}

\classesoc{25-39}{Ladri}{0.4}

Si procurano da vivere rubando e ricavano proporzionalmente
alla loro abilit\`a e alla ricchezza delle loro vittime.

\begin{abilist}
\item ADL
\item ADT o ATC
\item Acrobazia
\item Borseggiare
\item Cavalcare bipedi
\item Geografia 
\item Individuare passaggi segreti
\item Individuare trappole
\item Lingua levante
\item Lingua madre
\item Muoversi al buio
\item Muoversi in silenzio
\item Nascondersi
\item Scassinare
\item Scalare
\item Travestirsi
\end{abilist}

\classesoc{40-59}{Saltimbanchi}{0.4}

Girovagano nelle terre conosciute organizzando spettacoli di giochi, balli e
musica e riescono a guadagnare dalle 2 alle 20 corone al mese per famiglia.

\begin{abilist}
\item ATC
\item Acrobazia
\item Ammaestrare
\item Barare o Contorsionismo o Prestidigitazione o Giocoleria o
  Ventriloquio
\item Cadere
\item Cavalcare quadrupedi
\item Geografia
\item Lingua levante
\item Lingua madre
\item Saltare
\item Scalare
\end{abilist}

\classesoc{60-92}{Artigiani}{0.5}

Svolgono lavori di utilit\`a e vendono i loro prodotti. Possono
essere: falegnami, pellai, fabbri o stagnini, cartografi, gioiellieri,
intagliatori (osso, legno). Possono ricavare da 2 a 20 corone al mese.

\begin{abilist}
\item ATC
\item Abilit\`a professionale \#1
\item Abilit\`a professionale \#2
\item Abilit\`a professionale \#3
\item Lingua levante
\item Lingua madre
\item Altre abilit\`a a scelta del Master
\end{abilist}

\classesoc{93-96}{Saggi}{0.7}

Sono custodi della tradizione del loro popolo. Vivono di ci\`o che
viene loro offerto dal popolo e della raccolta di erbe e dei loro
prodotti. Possono ricavare da 1 a 50 scudi al mese.

\begin{abilist}
\item ATC
\item Alchimia
\item Astrologia o
\item Lettura della mano
\item Conoscere piante
\item Lingua Madre
\item Lingua Levante
\item Medicina
\item Mitologia
\item Occultismo
\item Raggirare
\item Religione
\end{abilist}

\classesoc{97-100}{Capiclan}{0.6}

Sono preposti al mantenimento dell'ordine e all'amministrazione della giustizia.
Provvedono alle esigenze della propria famiglia ricevendo uno stipendio mensile di
circa 10 corone oltre al necessario per vivere.

\begin{abilist}
\item ATC
\item Calcolare
\item Cavalcare quadrupedi
\item Giurisprudenza
\item Lingua levante
\item Lingua madre
\item Lingua straniera \#1
\item Lingua straniera \#2
\item Mitologia
\item Oratoria
\item Raggirare
\item Recitare
\item Religione
\item Travestirsi
\end{abilist}


\section{I Nani}

\subsection{Gennoria}

\Stendardo{stendardo_nani.eps}

I Nani popolano il Regno di Gennoria.

\Geografia Il regno di Gennoria si sviluppa nella parte montuosa dell'isola di
Kush, occupando tutto il massiccio montuoso del Genna-ruxi, la cui
cima pi\`u alta \`e punta Alabarda, alta 7020 metri. Altra catena
montuosa importante, se non altro sul piano culturale, \`e il
Monteferro.

Numerosi sono i fiumi, nessuno dei quali \`e navigabile per via
delle forti pendenze e dell'irregolarit\`a del corso. La capitale
del regno \`e Gennor, situata alla profondit\`a di 100 metri nel
sottosuolo. Da essa si dipanano i tunnel principali che conducono alle
altre citt\`a sotterranee, fra cui: Horun, Solorg, Bithi. L'unica
citt\`a dotata di un porto, e situata sulla superficie, \`e Neoh
Nely. La vegetazione sulle montagne \`e scarsa e di tipo arbustivo;
sulle coste, invece, predominano le pinete.

\Politica Tra i nani, solo coloro che hanno partecipato ad una battaglia
acquistano i diritti politici e vengono chiamati firmatari. Tra
costoro, chi riesce a raccogliere almeno 1000 firme di suoi pari
diviene Console a vita, guida in guerra i suoi firmatari, li difende
in giudizio di fronte al re e amministra personalmente la giustizia
nei confronti dei non firmatari che a lui si rivolgono. 

Il Re ha la funzione di amministrare la giustizia sui firmatari e di
regolare e organizzare le attivit\`a dei Consoli. La carica di Re
\`e ereditaria, ma nel caso egli si riveli un inetto, viene
spodestato e sostituito con uno eletto dal Concilio dei Consoli.
Finora il numero maggiore di generazioni di una stessa famiglia che
hanno detenuto la carica di re \`e di 3, e 3 sono gli anni che
solitamente si aspettano prima di spodestarne uno. Il sistema
giuridico \`e basato sulla legge del taglione: occhio per occhio,
dente per dente.  Non esiste alcun controllo economico.

\Esercito Ad ogni maggiorenne \`e consentito partecipare alle battaglie per
acquisire i diritti politici. A tal fine, in periodo di pace,
i giovani nani prestano servizio come mercenari presso altri popoli.

Il controllo delle operazioni militari \`e realizzato dal re che, di
volta in volta, decide quali consoli, con relativi firmatari,
mobilitare.  

L'insieme di un console e dei suoi firmatari costituisce
una Schiera Consolare.  

I fronti di battaglia preferiti sono le
gallerie e le montagne, ove \`e temuta la loro cavalleria di mufloni
giganti. 

Le armi tradizionalmente pi\`u usate sono: l'Ascia da
Battaglia, il Martello da Guerra, e l'Alabarda. I Nani sono in guerra
con i Reuben per il controllo di alcune gallerie e odiano
profondamente i Rom del deserto che cercano costantemente di depredare
il loro cimitero.

Nella citt\`a di Gowr-nhen si trova la principale
Accademia di Arti Marziali.

\Economia La ricerca del successo personale si sposa perfettamente con il tipo
di economia sviluppata dai nani. Non esiste alcun controllo da parte
delle autorit\`a tranne una tassazione del 3\% sul reddito, di cui il 2\% per 
i consoli e l'1\% per il re. Ogni nano organizza il proprio lavoro al
fine di trarne il massimo guadagno. Per tale ragione, spesso, i nani
si uniscono in societ\`a.

Le importazioni e le esportazioni sono frutto delle contrattazioni tra
singoli nani o societ\`a e mercanti, siano essi stranieri o nani.

Estraggono dal sottosuolo oro, argento, piccole quantit\`a di
ferro, insufficienti per l'esportazione, e il rarissimo Myrtium, un
metallo raro, leggerissimo, resistente e molto, molto prezioso.
Inoltre producono i migliori marmi e graniti delle isole, che
normalmente vengono esportati da squadre di costruttori al fine
di curarne la posa.

Coltivano i vegetali necessari alla produzione
del Phyluf-Herhu (il liquore tipico) e della birra, la pi\`u
alcolica dell'Arcipelago.

Allevano e ammaestrano I M\`uvara, i mufloni giganti, come
cavalcature e al fine di ricavarne carne e latte (che fanno cagliare
producendo lo Jodhu).  Le corna vengono lavorate per ricavarne il
manico dei pugnali. Importano cereali, ortaggi, legno, metalli vili,
pesce e tessuti.

\Religione La religione professata dai nani \`e basata sul culto dei
morti. La sua particolarit\`a \`e che non esistono sacerdoti, ma ogni
famiglia venera privatamente i propri antenati.  

Ogni nano cerca con
le sue azioni di essere degno dei propri antenati.  Se la sfortuna
vuole che tra essi si annoveri un individuo di cui essere poco fieri,
il nano far\`a di tutto per celarlo agli altri e per cancellare tale
macchia dalla storia della famiglia. 

La vita del nano \`e permeata dal pensiero della morte, infatti inizia
prestissimo a preparare la propria tomba, costituita da una stanza
scavata nella catena montuosa del Monteferro, adibita a cimitero, e a
conservare le erbe, gli unguenti e le garze per la propria
mummificazione. Per un nano \`e fondamentale avere dei discendenti che
curino la sua tomba.

Le cerimonie pi\`u importanti sono: il passaggio alla maturit\`a (a 16
anni) consistente nel dare il primo colpo di scalpello al sito ove
verr\`a edificata la tomba, e il funerale, a cui solitamente partecipa
anche il Console.

\Magia Il regno di Gennoria \`e affiliato al Conclave. Le accademie
richiedono una retta piuttosto alta per cui soltanto i pi\`u abbienti
possono permettersi di studiare la magia, suscitando l'invidia degli
altri, poich\'e \`e da loro vista come un'affascinante aiuto sia
economico che militare.  Questo ha portato allo svilupparsi
dell'Anticonclave, e di conseguenza ad una massiccia repressione delle
attivit\`a magiche irregolari.

Le scuole insegnate sono: Elementalismo e Necromanzia, quest'ultima in
particolare viene utilizzata per chiedere consigli ai propri antenati.

\Moda Maschi e femmine hanno presso i nani pari diritti e doveri.
Bench\'e sia importantissima la storia della famiglia, ogni individuo
viene considerato solo in funzione del successo che riesce a
conseguire; pertanto la vita di un nano consiste in un gareggiare
continuamente con gli altri nani al fine di dimostrarsi il migliore.
Un'offesa personale \`e solitamente lavata col sangue dell'offensore e
questo porta a faide interminabili.  

L'abbigliamento tipico consiste
in: stivali neri, brache di tela pesante bianchi, corpetto di panno
nero, camicia di lino bianca e sul capo l'immancabile berretto. Sul
lavoro si utilizzano grembiuli di cuoio.

  I divertimenti pi\`u comuni
sono i giochi di carte e i dadi. Lo stress causato dalla ricerca del
successo \`e spesso affogato nel forte liquore nanico, il Phyluf-Herhu. 

La lingua ufficiale \`e il Nanico.

\pinup{nano.eps}{Un tipico abitante di Gennoria}

\Fisico Alti fino a 140 cm e pesanti
fino a 80 kg, hanno una struttura ossea e muscolare imponente. Bench\'e
la carnagione sia olivastra e gli occhi nerissimi, i capelli sono
solitamente grigi e mossi. Il volto di un nano maschio, anche giovane,
\`e sempre carico di rughe e la peluria sul suo corpo \`e
abbondante.  I Nani possono vedere nel buio pi\`u totale percependo
le fonti di calore (con un TOSS). 


\minmaxnani

\subsubsection{Abilit\`a di razza}
\begin{abilist}
\item Resistenza al caldo 10
\item Visione termica
\end{abilist}

\subsubsection{Classi sociali e abilit\`a di Classe sociale}

\classesoc{01-24}{Minatori}{0.3}

 Vivono dei proventi ottenuti dall'estrazione dei metalli
e dalla loro lavorazione, oppure vengono pagati dai proprietari delle
miniere. I loro introiti variano tra 5 e 500 scudi al mese.

\begin{abilist}
\item ABL
\item CAC
\item AO
\item Carpenteria
\item Geologia
\item Fare nodi
\item Ingegneria
\item Lavorare Pietre
\item Lingua Madre
\item Lingua Levante
\item Muoversi al buio
\item Orientamento
\item Resistenza alla Fatica
\item Sollevare Pesi
\item Valutare Pietre
\end{abilist}

\classesoc{25-44}{Costruttori}{0.3}

Si occupano di edilizia e carpenteria, edificano a pagamento case,
palazzi, strade, gallerie ecc. Guadagnano da 10 e 1000 scudi al
mese.

\begin{abilist}
\item ABL o ATL
\item CAC
\item AO
\item Architettura
\item Carpenteria
\item Geologia
\item Fare nodi
\item Ingegneria
\item Lingua Madre
\item Lingua Levante
\item Resistenza alla Fatica
\item Scalare
\item Sollevare Pesi
\end{abilist}

\classesoc{45-64}{Contadini}{0.3}

Coltivano campi ed allevano bestiame. Producono cibo e vestiti per la
propria famiglia e possono guadagnare al netto delle spese di
sostentamento da 1 a 10 scudi al mese.

\begin{abilist}
\item AIA o ABL o ATL
\item CAC
\item AO
\item Ammaestrare
\item Carpenteria
\item Conoscere animali
\item Conoscere piante
\item Lingua levante
\item Lingua madre
\item Resistenza alla fatica
\end{abilist}

\classesoc{65-89}{Professionisti}{0.5}

Praticano i pi\`u svariati mestieri e vivono di essi. Possono essere: artigiani,
militari, cartografi, giuristi, medici, erboristi, carpentieri, macellai ecc.
Il loro introito mensile pu\`o variare da 5 a 100 corone al mese.

\begin{abilist}
\item CAC
\item AO
\item Abilit\`a professionale \#1
\item Abilit\`a professionale \#2
\item Abilit\`a professionale \#3
\item Lingua Levante
\item Lingua Madre
\item Altre abilit\`a a scelta del Master
\end{abilist}

\classesoc{90-97}{Consoli}{0.7}

I consoli ricevono un appannaggio fisso dai loro firmatari: da 100 a
500 corone al mese.

\begin{abilist}
\item ABL
\item ATL
\item AIA
\item CAC
\item AO
\item Giurisprudenza
\item Lingua Levante
\item Lingua Madre
\item Oratoria
\item Tattica Militare
\item Trattare
\item Altre abilit\`a a scelta del Master
\end{abilist}

\classesoc{98-100}{Re}{0.7}

Il re guadagna dal fisco 5000 o pi\`u corone al mese.

\begin{abilist}
\item ABL
\item ATL
\item AIA
\item CAC
\item AO
\item Giurisprudenza
\item Lingua Levante
\item Lingua Madre
\item Oratoria
\item Tattica Militare
\item Trattare
\item Altre abilit\`a a scelta del Master
\end{abilist}


\section{Gli Orchi}

\subsection{Terhamala}

\Stendardo{stendardo_orchi.eps}

Gli orchi, divisi nelle due etnie Pran-Horu e Haroka, abitano il
territorio di Terhamala.

\Geografia Terhamala occupa la parte settentrionale di Malagana, dal confine con
Umbrosa a Sud fino a Capo Ondoso a nord. Il territorio \`e
prevalentemente aspro e montuoso anche se a Sud, in prossimit\`a del
Massiccio di Mugrun, si estende un ampio pianoro. Le montagne del
Massiccio raggiungono altitudini ragguardevoli, con ben tre cime che
superano i 4000 metri e una che supera i 5000 (il monte Khan Hooga,
5109 metri).  Gran parte del Massiccio \`e perennemente coperta
dalle nevi. Le foreste di conifere sono poco estese.

I corsi d'acqua sono brevi, veloci e non navigabili e le coste ove
sfociano sono alte e frastagliate. Sul Pianoro \`e situato il lago
di Moh Lehnt Arjus, la cui superficie gela durante gli inverni pi\`u
rigidi. Il lago \`e alimentato dal fiume Maj-Jiahn. L'unica
citt\`a degna di nota, poich\'e l'unica ad essere costruita in
muratura, \`e Bahrrak.

\Politica Terhamala \`e occupata interamente dalle nazioni degli Orchi (Haroka
e Pran-Horu), il cui ordine politico \`e fortemente frammentato.

Ogni trib\`u ha nel capotrib\`u, che pu\`o anche essere femmina,
il supremo capo militare, giudiziario e legislativo. Esso ha il potere
di vita e di morte su tutti i componenti della trib\`u. 

Le varie
trib\`u sono spesso in guerra fra loro e la pace costa sempre molte
vite. Questa \`e la ragione principale per cui gli orchi sono stati
facilmente sconfitti e costretti a spostarsi progressivamente a nord
dall'esercito degli Umbra; di quella che una volta era la loro isola
agli orchi \`e rimasto solo quanto loro concesso. 

Ogni anno le trib\`u nominano un guerriero (Kah Dohts), uno per
ciascuna di esse.  I prescelti dovranno sfidarsi in una competizione
che provi la loro abilit\`a. La competizione ha di solito carattere
militare. Il vincitore consegner\`a al proprio capotrib\`u il
titolo di Capoguerra, attribuito a colui che ha il compito di guidare l'esercito in
guerra contro nemici esterni. 

Il capotrib\`u trasmette
solitamente la sua carica al proprio primogenito, ma chiunque pu\`o
sfidarlo a duello e prenderne il posto se riesce a sconfiggerlo.

Accade molto spesso che a tenere realmente le redini del governo di
ciascuna trib\`u sia in realt\`a il Sacerdote, che funge da
consigliere e da interprete dei segni delle divinit\`a. La nazione
dei Pran-Horu occupa il Pianoro, mentre gli Haroka vivono sul
Massiccio di Mugrun.

\Esercito Ogni trib\`u ha il suo esercito, formato da maschi e femmine abili
al combattimento. Solo dopo la dura sconfitta subita nella guerra
contro gli Umbra gli orchi hanno pensato a creare un esercito
unitario. In caso di conflitto il Capoguerra ha li ruolo di supremo
comandante di questo esercito.  L'esercito orchesco \`e penalizzato
dall'esiguo numero di componenti, ma in quanto a pericolosit\`a sul
campo di battaglia \`e temuto da tutti i regni. Gli orchi si
addestrano fin da piccoli all'uso delle armi (Ascione, Ascia Bipenne,
Spadone a 2 mani, Morning Star, Martello da Guerra e Alabarda); sono
inoltre abili nel combattimento.

L'esercito unitario \`e suddiviso in: fanteria, fanteria pesante e
cavalleria (questo corpo \`e quasi interamente composto da orchi
Pran-Horu). Ogni capotrib\`u comanda il suo contingente.
l'addestramento al combattimento riveste una grande importanza per
tutta la popolazione.

\Economia L'economia degli orchi \`e indirizzata verso l'autarchia. Per quello
che riguarda gli Haroka le attivit\`a principali sono l'estrazione e
la lavorazione dei metalli per la fabbricazione di armi ed armature
sobrie, solide e di ottima fattura. l'attivit\`a principale degli
orchi Pran-Horu \`e l'allevamento, praticato comunque anche dagli
Haroka.  Entrambe le etnie si dedicano comunque anche ad attivit\`a
quali la tessitura, l'agricoltura e il baratto. Nessuno \`e tenuto
al pagamento delle tasse e ognuno \`e arbitro della propria
condizione economica. 

Gli orchi esportano in tutto l'Arcipelago le loro armature, le loro
armi e i loro robusti cavalli da combattimento.  Il loro ponte per il
commercio \`e costituito dalle imbarcazioni degli elfi Luxi e dei
Chire. Riveste una certa importanza anche la caccia alle foche. Le
importazioni riguardano soltanto il Sute (sostanza
gommosa ricavata dai Segram, usata per impermeabilizzare le barche) e
degli Archi Elfici.

\Religione Ogni trib\`u venera il fondatore della trib\`u come il proprio
dio. I Sacerdoti (o le Sacerdotesse) sono i depositari di riti e
segreti del culto e i principali interpreti della volont\`a del dio.
Le festivit\`a variano da trib\`u a trib\`u e spesso anche di
anno in anno, a seconda della mutevole volont\`a divina. Le
cerimonie prevedono il sacrificio di animali e talvolta anche di
membri della trib\`u o di prigionieri di guerra. I Sacerdoti
tramandano la carica al proprio primogenito, sia esso maschio o
femmina.

\Magia Non esistono accademie di magia a Terhamala.
Tutta la conoscenza magica viene tramandata fra i membri delle
famiglie (max 3 punti) e in misura maggiore dagli sciamani (max 8
punti). Periodicamente, data l'assenza di accademie affiliate al
Conclave, vengono effettuati dei controlli segreti da ispettori del
Conclave per accertare il rispetto delle Leggi. Le uniche scuole
insegnate sono Necromanzia e Demonologia. 

\Moda Gli orchi vestono rozzi indumenti di pelle o di lana coperti da pezzi
di armatura. Sono poco socievoli e ``poco monogami''.  Solitamente un
maschio possiede un harem di due o tre femmine, ma pu\`o accadere,
anche se raramente, che una femmina particolarmente forte possieda un
harem di maschi. 

 La tessitura \`e un
ruolo prettamente femminile, mentre la concia e la cucitura delle
pelli spettano ai maschi.

I loro giochi preferiti sono la S'thruh-mppa, che \`e uno stile di lotta
molto violento, e il gioco d'azzardo.

\pinup{orco_small.eps}{Kah Dohts}

Alla sera gli orchi si concedono spesso delle solenni sbronze a base
di Mutta, un liquore ricavato dalla fermentazione delle bacche
dell'omonima pianta, tipica del posto.  Sotto gli effetti dell'alcool
accadono degli scontri fra i maschi o le femmine dominanti per
l'accaparramento dell'harem altrui.  Il perdente di solito muore.

\`E caratteristica la loro pettinatura: si radono la sommit\`a del
cranio e raccolgono in una treccia dietro la nuca i capelli che
portano lunghi. La lingua ufficiale \`e l'Orchesco.

\Fisico Gli orchi sono molto pi\`u alti e robusti degli esseri umani e degli
elfi: hanno un'altezza media di 2 metri e 10, ma non \`e difficile
incontrare orchi che superano i 2 metri e $\frac{1}{2}$.

La loro massa muscolare \`e generalmente in linea con l'altezza,
questo conferisce loro notevole forza e agilit\`a.  Il loro aspetto
fisico \`e in compenso abbastanza sgradevole: sono quasi
completamente coperti di peli scuri, hanno la pelle giallastra e la
fronte bassa, una dentatura molto robusta e canini grandi, aguzzi e sporgenti
come vere e proprie zanne.

La  mandibola, gli zigomi e le arcate sopracciliari sono
pronunciatissime.

Gli orchi parlano numerosi dialetti simili tra loro, 
riconducibili ai due rami Pran-Horu e Haroka della lingua Orchesca.

\subsection{Orchi Pran-Horu}

\minmaxorchi

\subsubsection{Abilit\`a di razza}

\begin{abilist}
\item Resistenza al freddo 10
\end{abilist}


\subsubsection{Classi sociali e abilit\`a di Classe sociale}

\classesoc{1-9}{Contadini}{0.3}

Coltivano campi ed allevano bestiame. Producono cibo e vestiti per la
propria famiglia e possono guadagnare al netto delle spese di
sostentamento da 1 a 10 scudi al mese.

\begin{abilist}
\item Carpenteria
\item Cavalcare
\item Conoscere Animali
\item Conoscere Piante
\item Lingua Levante
\item Lingua Madre
\item Resistenza alla fatica
\end{abilist}

\classesoc{10-19}{Pescatori}{0.3}

Vivono della pesca e dello sfruttamento del mare e dei suoi proventi.
Producono cibo per la propria famiglia e possono guadagnare al netto
delle spese di sostentamento da 1 a 10 scudi al mese.

\begin{abilist}
\item Carpenteria
\item Conoscere Animali
\item Controllo
\item Respirazione
\item Fare nodi
\item Lingua Levante
\item Lingua Madre
\item Navigare
\item Nuotare
\item Orientamento
\end{abilist}

\classesoc{20-34}{Cacciatori}{0.3}

Vivono della caccia e dei suoi proventi. Producono cibo e vestiti per
la propria famiglia e possono guadagnare da 2 a 7 corone al mese.

\begin{abilist}
\item Allarme
\item ADT o ADL
\item ATC o ABC
\item Conoscere animali
\item Conoscere piante
\item Fare nodi
\item Fiutare
\item Imitare
\item Individuare tracce
\item Lingua Levante
\item Lingua Madre
\item Muoversi in Silenzio
\item Trappole Udire
\end{abilist}

\classesoc{35-54}{Guerrieri}{0.4}

Ricevono dal clan un contributo mensile variabile tra 2 e 4 corone
oltre al vitto. Hanno poi diritto a prendere parte alla spartizione
del bottino di guerra.

\begin{abilist}
\item ADT
\item AIA
\item ATC o ABC
\item ABL
\item CAC
\item Maestria in un'arma
\item Allarme
\item Cavalcare Quadrupedi
\item Lingua Levante
\item Lingua Madre
\item Prontezza
\item Tattica militare
\end{abilist}

\classesoc{55-69}{Professionisti}{0.5}

Praticano i pi\`u svariati mestieri e vivono di essi. Possono essere:
artigiani, giuristi, cartografi, erboristi, carpentieri, medici,
macellai ecc. Il loro introito mensile pu\`o variare da 5 a 10 corone
al mese.

\begin{abilist}
\item Abilit\`a professionale \#1
\item Abilit\`a professionale \#2
\item Abilit\`a professionale \#3
\item Lingua Levante
\item Lingua Madre
\item Altre abilit\`a a scelta del Master
\end{abilist}

\classesoc{60-89}{Benestanti}{0.5}

Praticano molteplici mestieri. Hanno grandi esercizi
commerciali o possiedono grosse botteghe artigiane che forniscono loro
gran parte del reddito, variabile tra 10 e 100 corone al mese.

\begin{abilist}
\item ATL o ADT
\item Calcolare
\item Cavalcare quadrupedi
\item Danzare
\item Lingua levante
\item Lingua madre
\item Lingua straniera
\item Religione
\item Storia Generale
\item Trattare Valutare
\end{abilist}

\classesoc{90-96}{Borghesi ricchi}{0.6}

Sono proprietari di molteplici attivit\`a ed interessi, in cui
investono il proprio capitale. Sono perlopi\`u grossi allevatori,
possessori di miniere, mercanti di schiavi, ecc. Ricavano al mese da
50 a 500 corone.

\begin{abilist}
\item ATL o ADT
\item Calcolare
\item Cavalcare quadrupedi
\item Danzare
\item Lingua levante
\item Lingua madre
\item Lingua straniera
\item Religione
\item Storia
\item Generale
\item Trattare 
\item Valutare
\end{abilist}

\classesoc{93-96}{Sciamani}{0.7}

Sono custodi della memoria e dei riti religiosi dei clan. Vivono di
ci\`o che viene loro offerto dal popolo e della raccolta di erbe e
dei loro prodotti.

\begin{abilist}
\item Alchimia
\item Astronomia
\item Astrologia
\item Conoscere piante
\item Lingua Madre
\item Lingua Levante
\item Medicina
\item Mitologia
\item Religione
\end{abilist}

\classesoc{97-100}{Capitrib\`u}{0.6}

Provvedono alle esigenze del clan mantenendovi l'ordine e
amministrandovi la giustizia. Guidano la trib\`u in guerra.
Guadagnano da 5 a 15 corone di tributi al mese

\begin{abilist}
\item ABL
\item ADT
\item AIA
\item Calcolare
\item Cavalcare quadrupedi
\item Giurisprudenza
\item Lingua levante
\item Lingua madre
\item Lingua straniera
\item Mitologia
\item Oratoria
\item Religione
\end{abilist}

\subsection{Orchi Haroka}
\minmaxorchi

\subsubsection{Classi sociali e abilit\`a di Classe sociale}

\classesoc{1-30}{Pastori}{0.3}

Portano al pascolo il gregge, che d\`a latte, carne e pelli a tutta la
trib\`u. Possono guadagnare da 5 a 50 scudi al mese.

\begin{abilist}
\item ATC
\item Ammaestrare
\item Cavalcare quadrupedi
\item Conoscere animali
\item Conoscere piante
\item Individuare tracce
\item Lingua madre
\item Lingua levante
\item Orientamento 
\item Resistenza alla fatica
\end{abilist}

\classesoc{31-65}{Cacciatori}{0.3} 

Vivono della caccia e dei suoi
proventi. Producono cibo e vestiti per la propria famiglia e possono
guadagnare da 5 a 60 scudi al mese.

\begin{abilist}
\item Allarme
\item ADT o ADL
\item ATC o ABC
\item Conoscere animali
\item Conoscere piante
\item Fare nodi
\item Fiutare
\item Imitare
\item Individuare tracce
\item Lingua levante
\item Lingua madre
\item Muoversi in silenzio
\item Trappole
\item Udire
\end{abilist}

\classesoc{66-90}{Guerrieri}{0.4}

Ricevono dal clan un contributo mensile variabile tra 1 e 4 corone
oltre al vitto. Hanno poi diritto a prendere parte alla spartizione
del bottino di guerra.

\begin{abilist}
\item ADT
\item AIA
\item ATC o ABC
\item ABL
\item CAC
\item Maestria in un'arma
\item Allarme
\item Cavalcare quadrupedi
\item Lingua levante
\item Lingua madre
\item Prontezza
\item Tattica militare
\end{abilist}

\classesoc{91-95}{Sciamani}{0.6}

Sono custodi della memoria e dei riti religiosi dei clan. Vivono di
ci\`o che viene loro offerto dal popolo, della raccolta di erbe e
della vendita loro derivati.

\begin{abilist}
\item Alchimia
\item Astronomia
\item Astrologia
\item Conoscere piante
\item Lingua
\item Madre
\item Lingua Levante
\item Medicina
\item Mitologia
\item Religione
\end{abilist}

\classesoc{96-100}{Capitrib\`u}{0.5}

Provvedono alle esigenze
del clan mantenendovi l'ordine e amministrandovi la giustizia. Guidano la trib\`u
in guerra. Guadagnano da 5 a 10 corone di tributi al mese.

\begin{abilist}
\item ABL
\item ADT
\item AIA
\item Calcolare
\item Cavalcare quadrupedi
\item Giurisprudenza
\item Lingua levante
\item Lingua madre
\item Lingua straniera
\item Mitologia
\item Oratoria
\item Religione
\end{abilist}
\vskip 1cm
\section{I Reuben}

\subsection{Alanm-El-Shaddai}
\vskip 1cm
\Stendardo{stendardo_reuben.eps}
\vskip 1cm
I Reuben sono una razza umanoide abitante il Territorio di Alanm-El-Shaddai.

\Geografia Alanm-El-Shaddai si estende nella parte nord-occidentale
dell'isola di Kush, fra il massiccio di Genna-ruxi e il mare
interno e confina a Sud e a est con Gennoria.

Il clima \`e torrido, il territorio \`e totalmente occupato da un
deserto sabbioso punteggiato da oasi, frequentemente abitate dalle
trib\`u nomadi.

La citt\`a pi\`u importante sulla superficie \`e Amath, un porto
franco, protettorato di Bahuney, abitato in prevalenza da mercanti.
Questo porto rappresenta l'unico contatto dei Reuben con le altre
isole.

Il territorio abitato dai Reuben \`e costituito da un complesso di
enormi gallerie sotterranee che collegano grotte naturali dove sorgono
le cinque citt\`a: Eilekan, Elezer, Ennora, Rephidim, Sin e la
capitale Shaddai. Le gallerie sono illuminate da particolari muschi
fluorescenti.

L'uscita delle gallerie coincide spesso con le oasi, contese ai
nomadi, e una di esse sbocca nei pressi della citt\`a di Amath.
L'estensione delle gallerie sotterranee \`e causa di tensioni e
conflitti territoriali con i Nani.

\Politica La popolazione \`e suddivisa nelle cinque trib\`u Eilekan, Elezer,
Ennora, Rephidim, Sin, ciascuna avente come capitale l'omonima
citt\`a. Ogni trib\`u \`e specializzata in un ristretto insieme di
attivit\`a.

Il Re, eletto a suffragio universale, regna fino alla morte o fino
alle sue dimissioni. L'elezione prevede un periodo di campagna
elettorale autofinanziata da parte dei candidati dei partiti
conservatore e liberale e dei candidati indipendenti. 

Il compito del Re \`e coordinare le attivit\`a economiche, compito nel
cui esercizio \`e coadiuvato dal consiglio delle Arti e dei Mestieri,
i cui rappresentanti sono eletti dai membri delle Arti e dei Mestieri
con le stesse modalit\`a del Re.  Il consiglio si riunisce una volta
al mese per discutere i prezzi e la quantit\`a delle mercanzie da
esportare e importare. Alla fine di ogni seduta, che pu\`o durare
pi\`u giorni, un resoconto \`e fornito al Re che decide se avvalorare
ci\`o che \`e stato deciso in assemblea
  
Il Re ricopre la carica di somma autorit\`a religiosa, che esercita
basandosi sulle interpretazioni del libro sacro (Ras-El-Shaddai, la
testa dello splendente) date dal sommo sacerdote della trib\`u Sin.

Il Re \`e anche il supremo comandante dell'Esercito, per la direzione
del quale si avvale delle consulenza del Generale Nero (Per le trib\`u
Sin, Rephidim, Elezer, Ennora) e del Generale Bianco (Eilekan),
entrambi eletti a suffragio universale.

Il potere giudiziario \`e amministrato dai sacerdoti della trib\`u
Sin, facendo riferimento al Ras-El-Shaddai.  Gli Eilekan amministrano
l'ordine pubblico applicando le sentenze.

Il partito Conservatore persegue la vittoria finale contro i nomadi
del deserto e contro i nani. \`E sostenitore di un'economia di
sussistenza.

Il partito Liberale considera le altre razze degli sfortunati che non
conoscono la vera via segnata da El-Shaddai. Sostengono l'utilit\`a
anche bellica di un'economia aperta.

\pinup{reuben_guerriero.eps}{Guerriero Reuben}

\Esercito Ogni individuo \`e parte dell'esercito e viene addestrato
scrupolosamente al combattimento sin dalla pi\`u tenera et\`a,
entrando a far parte dell'esercito all'et\`a di 10 anni. 

Per un giorno
a settimana ogni Reuben partecipa alle esercitazioni di gruppo; ogni
trib\`u combatte unita secondo la propria specializzazione comandata
dal pi\`u anziano, che ha il compito di definire la gerarchia
militare. 

Le armi dei Reuben sono caratterizzate dalla presenza di elementi
meccanici che ne modificano la struttura con semplici movimenti.  Tra
queste, lo Shi-Rik (lama di scorpione) e lo Shi-Tai (lama da coda) sono
le pi\`u diffuse.

La tipica armatura Reuben \`e il Du-Ar (Armatura di Ragno), costruita
con esoscheletri di un grosso ragno sotterraneo chiamato
Marhya-Far-Hank e bende ricavate dalla sua tela. Caratteristico \`e
l'elmo per via delle sue grandi dimensioni.

Branca particolare dell'esercito sono gli Eilekan, costituita da tutti
i nati albini (circa 1 decimo della popolazione) che addestrati
rigorosamente sin dalla nascita, senza poter conoscere i loro genitori,
ricoprono il ruolo di killer o guastatori creando il caos all'interno
degli accampamenti o delle citt\`a che si intende conquistare.

\Economia I Reuben sono i principali produttori di droghe di tutto il
Continente, la pi\`u preziosa delle quali \`e la Khan che per loro
ha anche un significato religioso. 

La coltivazione delle erbe da cui si ricavano le droghe
avviene all'interno di serre ricavate da caverne poste su pi\`u
livelli, ognuno con condizioni termiche diverse in maniera da simulare
i loro habitat sulla superficie. 

Importante \`e anche la coltura di tuberi, muschi e funghi (da cui
si ricavano balsami e oli), e foraggio (quest'ultimo in appositi campi
predisposti all'ingresso dei condotti sotterranei). Allevano in
prevalenza cammelli, cavalli, pecore per il trasporto, il pellame e la
carne.

 Le Marhya-Far-Hank vengono allevate
per la produzione di filo da tessitura, ma in cattivit\`a sono poco
longeve e di minori dimensioni. I Reuben esportano droghe (il cui uso
\`e diffusissimo presso tutti i popoli eccetto che tra i
Rileani), tessuti lavorati e balsami e oli profumati. Importano
principalmente metalli e legname.

\Religione La religione \`e la base della vita di un Reuben che nasce, vive e
muore in funzione di essa. Il Ras-El-Shaddai infatti \`e testo sacro
giuridico e organizzativo di tutta la societ\`a e viene tatuato sul
corpo di ogni individuo al compimento del decimo anno. I Reuben sono
monoteisti e adorano El-Shaddai (lo Splendente), creatore di tutta la
terra, che li salv\`o dal tradimento delle altre razze che
privarono i Reuben del dominio sulla superficie.

Secondo il libro sacro chiunque si converta alla loro fede diviene
parte integrante del popolo e in funzione della sua preparazione viene
adottato da una delle trib\`u.

A capo della religione si trova il Sommo Sacerdote della trib\`u Sin
che organizza il culto su tutto il territorio e personalmente esegue
le cerimonie presso l'altare principale di Shaddai.

I Giorni Sacri sono il primo giorno dell'anno, il primo giorno del
mese e il primo di ogni settimana.

Secondo la religione di El-Shaddai l'anima di un Reuben rispettoso
della propria fede si reincarna per continuare a servire El-Shaddai e
per tale ragione i piccoli sono tenuti in grande considerazione. 

La sorte per un Reuben infedele \`e la dispersione dell'anima. Non
esiste alcuna forma di matrimonio e i figli sono di colui che li
partorisce.

\Magia Non esistono accademie di magia ad Alanm-El-Shaddai, e perci\`o
sono diffusi gli insegnamenti da parte dell'Anticonclave all'insaputa
dei membri ufficiali del Conclave. Per quanto il Conclave sospetti
l'esistenza di tale influenza non \`e mai riuscito, nonostante l'invio
di ispettori, ad avere prove a riguardo.

I depositari della conoscenza magica sono gli Ennora che insegnano le
basi della magia a coloro che ritengono pi\`u meritevoli. Solamente
gli Ennora pi\`u anziani condividono fra loro le alte vette delle
conoscenze magiche, al fine di garantire la peculiarit\`a della loro
trib\`u, pertanto le uniche scuole insegnate ai giovani usufruitori
sono Elementalismo e Demonologia.

\Moda I Reuben
indossano ampie tuniche scure stette in vita da una cintura
caratterizzata dalla presenza di numerose tasche. Sotto di essa
coprono tutto il corpo con garze di tela di Marhya-Far-Hank che oltre
ad avere la funzione di unire i pezzi dell'armatura hanno la
propriet\`a di isolare il corpo dall'ambiente esterno. 

Tale aspetto
\`e particolarmente importante nelle terre dove sono presenti
frequenti precipitazioni al fine di evitare l'eccessivo contatto con
l'acqua, che pu\`o provocare il Male Freddo. Le scarpe, molto
flessibili grazie alla suola fatta di tela di Marhya-Far-Hank, permettono
ai piedi di espletare la funzione prensile.  

I Reuben portano sempre con s\`e
garze di tela di ragno e i balsami per la pulizia personale. Per
rimuovere il balsamo e le impurit\`a si servono di particolari
insetti pulitori chiamati Ploes (che in antica lingua Reuben significa
centopiedi spiritosi e divertenti che svolgono la funzione di
pulitori). 

I Reuben vivono in case scavate sulle pareti e sul soffitto
delle caverne con gli ingressi collegati da strutture in legno. I
divertimenti pi\`u comuni sono i combattimenti simulati con armi
innocue, la caccia alle Marhya-Far-Hank selvatiche e la caccia ai
membri delle trib\`u nomadi e ai nani, che vengono rapiti
appositamente. 

La loro dieta \`e caratterizzata prevalentemente da
arrosti e non disdegnano di nutrirsi della carne dei nemici uccisi in
battaglia. 

La lingua ufficiale \`e il Reuben.

\Fisico I Reuben hanno un'altezza media di 240 cm con la coda lunga 70/80
cm, il loro peso medio \`e di 130 kg; hanno spalle e torace ampi fino ad
un metro e mezzo e braccia muscolosissime che giungono oltre le
ginocchia; la carnagione i capelli e gli occhi sono sempre scuri
tranne per gli Eilekan che sono albini. I capelli, lunghi, sono
raccolti in treccine; lo stesso vale per la barba portata da tutti
tranne che dagli Eilekan, che rasano tutto il corpo. 

I Reuben sono ermafroditi, presentano cio\`e entrambi i sessi anche
se l'aspetto esteriore \`e quello di un maschio. I piccoli nascono
da uova. 

Piedi e coda sono prensili e le loro membra sono facilmente
estensibili e contorcibili. 

Le loro uniche fonti d'acqua sono l'umidit\`a atmosferica, che
assorbono attraverso la pelle, e il cibo che ingeriscono.
L'ingestione diretta dell'acqua o anche un semplice eccesso di
umidit\`a provocano il Male Freddo.

I Reuben sono in grado di vedere al buio pi\`u totale rilevando le
fonti di calore (con un TOSS). I Reuben non ibridano con nessun'altra
razza.

\minmaxreuben

\subsubsection{Abilit\`a di razza}

\begin{abilist}
\item Contorsionismo 10
\item Orientamento 5
\item Religione 5
\item Resistenza al caldo 10
\item Scalare 5
\item Visione Notturna
\end{abilist}

\subsubsection{Handicap}

\begin{abilist}
\item Idrofobia
\item Male Freddo (a contatto con l'acqua)
\end{abilist}

\subsubsection{Classi sociali e abilit\`a di Classe sociale}

\classesoc{01-20}{Sin}{0.5}

\textbf{Sacerdoti}. Svolgono il compito di diffondere la parola di
El-Shaddai, amministrare il culto e la giustizia. Essi svolgono
inoltre una loro propria attivit\`a da cui traggono il loro
sostentamento. Guadagnano da 10 a 60 corone al mese.

\begin{abilist}
\item ABL
\item AIA
\item CAC
\item Cantare
\item Conoscere Piante
\item Danzare
\item Giurisprudenza
\item Mitologia
\item Oratoria
\item Religione
\item Resistenza al Caldo
\item Storia Generale
\item Abilit\`a professionale
\end{abilist}

\classesoc{21-40}{Eilekan}{0.5}

\textbf{Difensori della Fede}. Hanno il compito di infliggere le pene
comminate dai Sin. Essi svolgono inoltre una loro propria attivit\`a
da cui traggono il loro sostentamento. Guadagnano da 10 a 60 corone al
mese.

\begin{abilist}
\item ATC
\item ATL
\item CAC
\item Acrobazia
\item Cavalcare Quadrupedi
\item Giurisprudenza
\item Maestria in 1 Arma
\item Muoversi al Buio
\item Muoversi in Silenzio
\item Prontezza
\item Religione
\item Tattica militare
\item Torture
\item Abilit\`a professionale
\end{abilist}

\classesoc{41-60}{Rephidim}{0.5}

\textbf{Cavalieri}. Allevano e addestrano il bestiame. Nell'esercito hanno ruolo di
Cavalleria. Essi svolgono inoltre una loro propria attivit\`a da cui
traggono il loro sostentamento. Guadagnano da 20 a 50 corone al mese.

\begin{abilist}
\item ADT
\item AIA
\item CAC
\item Ammaestrare
\item Cavalcare Alati
\item Cavalcare Bipedi
\item Cavalcare Quadrupedi
\item Imitare
\item Individuare Tracce
\item Individuare Trappole
\item Resistenza al Caldo
\item Trappole
\item Abilit\`a professionale
\end{abilist}

\classesoc{61-80}{Elezer}{0.5}

\textbf{Ladri}. Il loro compito principale \`e quello di realizzare
furti e razzie a danno di Nomadi del deserto e dei Nani. Essi svolgono
inoltre una loro propria attivit\`a da cui traggono il loro
sostentamento. Guadagnano da 10 a 60 corone al mese.

\begin{abilist}
\item ADL
\item ATC
\item CAC
\item Acrobazia
\item Borseggiare
\item Cadere
\item Giocoleria
\item Individuare Passaggi Segreti
\item Individuare Trappole
\item Muoversi in Silenzio
\item Prestidigitazione
\item Resistenza al Caldo
\item Scassinare
\item Abilit\`a professionale
\end{abilist}

\classesoc{81-100}{Ennora}{0.5}

\textbf{Occultisti}. Hanno il compito di detenere il sapere, in
particolare quello magico. Essi svolgono inoltre una loro propria
attivit\`a da cui traggono il loro sostentamento. Guadagnano da 20 a
50 corone al mese.

\begin{abilist}
\item ABL
\item ATC
\item CAC
\item Alchimia
\item Conoscere Demoni
\item Conoscere Animali Magici
\item Conoscere Esseri Onirici
\item Conoscere Esseri del Limbo
\item Crittografia
\item Medicina
\item Occultismo
\item Resistenza al Caldo
\item Valutare Oggetti Magici
\item Abilit\`a professionale
\end{abilist}


\section{I Giganti}

\subsection{Mahn-Hesa}

\Stendardo{stendardo_giganti.eps}

I Giganti vivono nel territorio di Mahn-Hesa 

\Geografia Mahn-Hesa occupa la parte inferiore dell'isola di Norda.
Confina ad Ovest con Terranova, a Nord e a est con le Terre del Nord
(attraverso l'Eisenval) ed a Sud con il Mare Interno. Il suo sbocco
sul mare si limita a comprendere la Baia Dei Giganti fra Capo
Bahs-Shyu e Punta Deys-Mathas.

Il clima nell'entroterra \`e di tipo continentale freddo con inverni
rigidi e nevosi ed estati tiepide e secche, mentre a Sud la
temperatura media assume valori leggermente pi\`u alti grazie al mare
che mitiga notevolmente il clima. Qui le precipitazioni sono frequenti
anche in estate.

La superficie del territorio \`e prevalentemente collinare con lievi
pendenze che degradano dolcemente fino alla Baia dei Giganti che
presenta un litorale basso e sabbioso.

A Nord, Mahn-Hesa \`e separata da Terranova dal Massiccio del Dehrberg
dal quale ha origine il Rijsh-Tahl (grande fiume) il fiume pi\`u
importante della regione.  Nella sua parte pi\`u meridionale il
Rijsh-Tahl \`e navigabile per 20 km. In questo tratto si possono
trovare i resti della vecchia citt\`a portuale di Neuphlow,
praticamente disabitata, residuo della dominazione Levante.

A Nord la vegetazione naturale \`e costituita in gran parte da foreste
di conifere, in cui predomina la tipica sagoma imponente a forma di
cono dell'abete Novese; a Sud occupano gran parte del territorio le
brughiere e i pascoli.

\Politica Il territorio di Mahn-Hesa \`e abitato dai Giganti. Non esiste un
vero e proprio stato, i giganti sono divisi in Clan composti da una
sola famiglia allargata (comprendente solitamente il padre, la madre,
i figli naturali o acquisiti ed i nipoti). Non esiste nessun ordine
gerarchico all'interno delle famiglie; in generale vi \`e per\`o
una certa superiorit\`a del padre sia sulla madre che sui figli e
sui nipoti. \`E lui che di solito prende le decisioni e coordina la
vita del Clan. Alla morte del padre nella maggior parte dei casi il
Clan si scioglie e si vanno a formare dei Clan separati capeggiati dai
figli, altrimenti il figlio maggiore assume il titolo di Capoclan.

I Clan sono riuniti in villaggi, che sono di solito formati da un
numero di famiglie variabile fra 10 e 15. Le eventuali controversie
fra gli abitanti sono risolte dal Druido del villaggio, che \`e
considerato il capo spirituale e giudiziario; non sono rari comunque
gli scontri fisici fra contendenti che non sono riusciti a mettersi
d'accordo (anche se quasi mai questi terminano con la morte di uno di
essi).

\Esercito Non esiste un esercito unitario di Mahn-Hesa.  In caso di guerre fra
villaggi vicini, i Druidi provvedono a reclutare tutti coloro in grado
di utilizzare armi come clave o bastoni (solitamente coloro che fanno
\label{armigiganti} parte di Clan di cacciatori) e a dargli una
parvenza di unit\`a, ma questa eventualit\`a \`e piuttosto remota
visto che i giganti vivono in pace con tutti.

Per scongiurare ogni tipo di invasione da parte di Novesi o Nordici e
per evitare che questi popoli possano sottomettere gli abitanti dei
villaggi di Mahn-Hesa utilizzandoli come ``macchine da guerra'', lungo
tutti i confini del territorio con Terranova e le Terre del Nord \`e
stato eretto un muro alto una decina di metri presidiato da
guarnigioni novesi e nordiche che regolano il traffico dei commerci.
Esso \`e chiamato Eisenval ed \`e il risultato di un patto tra le
due nazioni confinanti con Mahn-Hesa. Per tutta la sua lunghezza
(circa 750 km) L'Eisenval \`e fortificato con grosse torri di
avvistamento posizionate a distanze differenti. Le Porte pi\`u
importanti sono la Porta di Norda e la Porta di Harus rispettivamente
sul lato est e sul lato Ovest del muro.

\Economia L'economia di Mahn-Hesa \`e molto arretrata e basata sul baratto. Le
attivit\`a sono prettamente a carattere familiare, quelle pi\`u
importanti sono l'allevamento e la caccia. l'allevamento, che
comprende anche la pastorizia, \`e svolto soprattutto dai clan che
abitano le regioni meridionali (Forth-lants) dove \`e pi\`u
abbondante la presenza di pascoli.

L'attivit\`a di caccia \`e invece concentrata nella parte nord del
territorio (Norhlants). Sono abbastanza frequenti, pertanto, i
baratti fra i Clan delle due regioni. Gli allevatori solitamente si
dedicano anche alla conciatura delle pelli e alla loro cucitura (che
\`e svolta in modo particolare dalle donne).

I giganti non hanno molti contatti con l'esterno se si eccettuano i
Clan di muratori che sono molto richiesti dai Novesi e dai Levanti per
la costruzione dei palazzi pi\`u imponenti. Non \`e praticamente
conosciuto l'uso della moneta ed anche il pagamento dei muratori da
parte degli stranieri solitamente avviene in natura, con prodotti
agricoli, utensili o, pi\`u generalmente, con qualcosa che possa
essere facilmente scambiato all'interno del territorio di Mahn-Hesa.

\Religione La religione \`e di tipo totemico. I giganti venerano la
natura ed associano ad ogni clan un animale che lo controlla e lo
protegge. Ogni clan venera un diverso animale che solitamente \`e
molto feroce, molto forte o molto agile. Esistono pertanto il clan del
lupo, del falco, dell'orso e via dicendo.

Sono previsti dei sacrifici in favore
degli animali venerati e di norma ad essi vengono offerte le loro
prede naturali.

 Il Druido \`e colui che regola questi sacrifici e li
concentra in particolari giorni dell'anno, come il solstizio d'inverno
o d'estate. I Druidi non appartengono a nessun clan e scelgono i loro
successori fra i bambini pi\`u meritevoli all'interno del villaggio.
Abitano in disparte e nessuno, che non sia da loro autorizzato,
pu\`o venire a conoscenza dei loro riti e dei loro segreti.

\Magia A Mahn-Hesa non ci sono accademie di magia. A tutti coloro che
i Druidi ritengano meritevoli, vengono insegnate le basi della magia
(Max 3 punti). Solo i successori dei Druidi hanno diritto ad una
conoscenza pi\`u approfondita (Max 8 punti).

Data l'assenza di accademie affiliate al Conclave, vengono effettuati
periodicamente dei controlli segreti da ispettori del Conclave per
accertare il rispetto delle Leggi. Le uniche scuole insegnate sono
Esistenza e Stregoneria.

\Moda I giganti indossano rozzi vestiti di pelliccia o di lana. Le
loro abitazioni sono costituite da mura circolari costruite a secco
con grosse pietre, e da tetti di arbusti e legna. 

I giganti sono prettamente monogami e molto fedeli. I matrimoni
vengono celebrati dal Druido ed in occasione di essi vengono
organizzate cerimonie che possono durare anche due settimane.

Durante il periodo estivo, si radunano dopo il tramonto attorno al
fuoco intonando canti propiziatori e ballando danze tribali.

Di notevole interesse per la loro particolarit\`a sono le tecniche di
caccia. Questa viene praticata a mani nude o con le clave e
solitamente le fasi principali sono tre:

\begin{enumerate}
\item Avvistare l'animale
\item Corrergli incontro
\item Ucciderlo a colpi di clava o soffocarlo con le mani
\end{enumerate}

Di solito i giganti non hanno prole numerosa. La mortalit\`a infantile
\`e abbastanza elevata, cosicch\'e famiglie con pi\`u di quattro figli
sono delle vere e proprie eccezioni. La lingua ufficiale \`e un idioma
molto primitivo detto Mahn.

\Fisico I giganti hanno un'altezza media di 3 metri/ 3 metri e mezzo e
peso medio di 200 chili. Hanno pelle bianca, carnagione chiara, occhi
piccoli e scuri e capelli neri. La loro vita media \`e di circa 120
anni.

\minmaxgiganti

\subsubsection{Abilit\`a di razza}
\begin{abilist}
\item Fiutare 7
\item Resistenza al freddo 10
\end{abilist}

\subsubsection{Classi sociali e abilit\`a di Classe sociale}

\classesoc{01-34}{Cacciatori}{0.2}

Vivono della caccia dei suoi proventi. Producono cibo e vestiti per la
propria famiglia e scambiano le pelli e i vestiti con cibo e altri
beni.

\begin{abilist}
\item Allarme
\item ADT o ADL 
\item ATC o ABC
\item Conoscere animali
\item Conoscere piante
\item Fare nodi
\item Fiutare 
\item Imitare
\item Individuare tracce
\item Lingua levante
\item Lingua madre
\item Muoversi in silenzio
\item Trappole
\item Udire
\end{abilist}

\classesoc{35-69}{Allevatori}{0.2}

Portano al pascolo il proprio gregge, che
d\`a latte, carne e pelli, che vengono utilizzate per il sostentamento e
lo scambio. Pu\`o guadagnare da 5 a 50 scudi.

\begin{abilist}
\item Ammaestrare
\item Cavalcare bipedi
\item Cavalcare quadrupedi
\item Conoscere Animali
\item Conoscere Piante
\item Empatia
\item Fare nodi
\item Lingua levante
\item Lingua madre
\item Orientamento
\item Trattare
\end{abilist}

\classesoc{70-95}{Muratori}{0.3}

 Costruiscono case e artefatti in muratura per conto di altri in
cambio di cibo, pelli e altri beni.

\begin{abilist}
\item Architettura
\item Calcolare
\item Carpenteria
\item Geologia
\item Ingegneria
\item Lavorare pietre
\item Lingua levante
\item Lingua madre
\item Resistenza alla fatica
\item Sollevare pesi
\end{abilist}

\classesoc{95-100}{Druidi}{0.5}

Sono custodi della memoria e dei riti religiosi dei clan. Vivono di ci\`o
che viene loro offerto dal popolo e della raccolta di erbe e dei loro prodotti.

\begin{abilist}
\item Alchimia
\item Astronomia
\item Astrologia
\item Conoscere piante
\item Lingua Madre
\item Lingua Levante
\item Medicina
\item Mitologia
\item Religione
\end{abilist}
\fi

\vfill
\pagebreak
\section{Cenni Storici}

Le conoscenze storiche sull'Arcipelago in epoche antecedenti la
Dominazione Levante sono incerte e nebulose. Si presume che le razze
vivessero organizzate in trib\`u alquanto primitive e con scarsi
contatti fra loro. Con il progresso delle tecniche di navigazione, i
Levanti, che allora abitavano tutta l'Isola di Dyshu-Thandy,
intrapresero dei viaggi di esplorazione delle isole vicine, che si
trasformarono poi in imprese militari volte alla conquista.

L'opera Levante si concluse con la formazione di un vastissimo impero
comprendente oltre la met\`a dell'Arcipelago, non riuscendo per\`o a
occupare  Malagana, Loydi-Genya e le Terre Orientali. Il
periodo storico durante il quale ci\`o avvenne viene chiamato
Dominazione Levante.

In quel periodo il popolo Elfico viveva unito nell'isola di Bahuney ed
ancora non esisteva la distinzione fra le etnie Luxi e Umbra.  Essi
accolsero passivamente la dominazione Levante ritirandosi
nell'entroterra.

Poich\'e il dominio Levante lasciava ampia libert\`a Religiosa e
Culturale, le Trib\`u umane dell'isola di Norda permisero una fusione
fra la loro cultura e quella Levante. Non fu cos\`{\i}  invece nell'Isola
di Kush. I Levanti, che riuscirono a conquistare la superficie
scacciando i Reuben nelle profondit\`a del sottosuolo, furono
successivamente costretti alla ritirata dalla strenua resistenza
opposta da Reuben e Nani alleati per riconquistare l'Isola.

Il progresso portato dai Levanti su Bahuney scaten\`o tensioni fra
gli elfi che volevano vivere nel lusso abbandonando le foreste, e gli
elfi conservatori. Queste tensioni sfociarono nella scissione pacifica
che si concluse con la nascita delle due etnie Luxi e Umbra e l'esodo di
questi ultimi verso l'isola di Malagana. 

I Levanti stipularono con i Giganti dei patti d'amicizia che
consentirono loro di evitare scontri contro tali possenti avversari.

Successivamente, a causa delle tensioni interne dovute al passaggio
dalla Monarchia di tipo Feudale alla Democrazia ed anche al
diffondersi del nuovo credo che si rifaceva agli insegnamenti di
Ryless, il controllo sulle colonie si indebol\`{\i}. Ci\`o port\`o
alla nascita del Regno di Terranova, composto dalle trib\`u Nordiche
che mantennero l'ordinamento feudale, del Pontificato, dove il clero
di Ryless prese il potere e della Terra di Levante, dove vinse la
democrazia.

Nel caos di tale periodo, chiamato Periodo della Nascita delle Nazioni,
le colonie di Bahuney e Norda furono completamente abbandonate. A
ricordare il fiorente passato rimangono solo i trattati commerciali e
pochi insediamenti.

Le alte sfere del Pontificato galvanizzate dalla loro fede decisero di
intraprendere azioni di conversione religiosa forzata in tutte le
terre, cercando l'appoggio del popolo dei Giganti. Tali mire furono
sventate dai Novesi e dalle trib\`u delle Terre del Nord, che
stipularono il patto del Muro edificando l'Eisenval a delimitare il
territorio di Mahn-Hesa. Le missioni Rileane non riuscirono nel loro
intento, limitandosi all'edificazione di avamposti. 

La fede Rileana, fondendosi con le religioni autoctone, port\`o al
formarsi del sincretismo religioso considerato eresia dalla Chiesa del
Pontificato. In questo periodo nacque il corpo degli Inquisitori.

Gli Umbra incontrarono a Malagana la ferrea resistenza delle autoctone
trib\`u orchesche.  Buona parte delle ricchezze provenienti da Bahuney
venne impiegata per assoldare le milizie di alcuni Signori orientali.
Grazie al loro aiuto, la guerra si concluse con il confinamento degli
Orchi nella parte settentrionale chiamata Terhamala e la nascita di
Umbrosa.  Gli Orchi si spartirono le terre montuose e pianeggianti
formando due diverse Etnie: Haroka e Pran-Horu. \`E al termine della
guerra tra Orchi e Elfi Umbra che si fa risalire la nascita del
Conclave.

Con la caduta della Dominazione Levante. i Chire, un popolo di
navigatori dalla pelle nera, proveniente dal Sud, cominciarono ad
intraprendere rapporti commerciali con le nazioni dell'Arcipelago
riuscendo a mettere in contatto gli Gnomi con gli altri popoli. Il
Conteggio Commerciale degli anni fa riferimento alla Nascita di
Ryless. 

Ad oggi ci troviamo nell'anno 1327 Dopo la Nascita di Ryless (DNR).

%%% Local Variables: 
%%% mode: latex
%%% TeX-master: "manual"
%%% End: 
